%%%%%%%%%%%%%%%%%%%%%%%%%%%%%%%%%%%%%%%%%%%%%%%%%%%%%%%%%%%%%%%%%%%%%%%%%%%%%%%%
% engenharia.tex
%
% Configurações específicas do livro "Engenharia em 500 Gramas": pacotes de
% engenharia (circuitos, gráficos de dados, diagramas de tempo), caixas
% didáticas semânticas e comandos do livro.
%
% Deve ser carregado DEPOIS de utils/preambulo.tex (que já carrega tikz,
% tcolorbox[many], listings, siunitx, xcolor[svgnames] e cleveref). Para isso,
% em livro.tex, logo após %%%%%%%%%%%%%%%%%%%%%%%%%%%%%%%%%%%%%%%%%%%%%%%%%%%%%%%%%%%%%%%%%%%%%%%%%%%%%%%%
% preambulo.tex
%
% Arquivo de configuração de packages para o LaTeX, considerando diversas
% classes que constumo utilizar.
%
% Para maiores informações, visite:
%    https://github.com/abrantesasf/latex
%    https://github.com/uvv-computacao/uvvtex2
%
% Para maiores informações, visite:
%    https://github.com/abrantesasf/latex
%
% NÃO ALTERE SE NÃO SOUBER O QUE ESTÁ FAZENDO!
%%%%%%%%%%%%%%%%%%%%%%%%%%%%%%%%%%%%%%%%%%%%%%%%%%%%%%%%%%%%%%%%%%%%%%%%%%%%%%%%


%%%%%%%%%%%%%%%%%%%%%%%%%%%%%%%%%%%%%%%%%%%%%%%%%%%%%%%%%%%%%%%%%%%%%%%%%%%%%%%%
%%% Estruturas de controle:
%%%%%%%%%%%%%%%%%%%%%%%%%%%%%%%%%%%%%%%%%%%%%%%%%%%%%%%%%%%%%%%%%%%%%%%%%%%%%%%%
% Por: Abrantes Araújo Silva Filho
%      abrantesasf@gmail.com
% URL: https://github.com/abrantesasf/latex


%%%%%%%%%%%%%%%%%%%%%%%%%%%%%%%%%%%%%%%%%%%%%%%%%%%%%%%%%%%%%%%%%%%%%%%%%%%%%%%%
%%% Carrega pacotes iniciais necessários para estrutura de controle e para a
%%% criação e o parse de novos comandos
\usepackage{ifthen}
\usepackage{xparse}
\usepackage{ifxetex}


%%%%%%%%%%%%%%%%%%%%%%%%%%%%%%%%%%%%%%%%%%%%%%%%%%%%%%%%%%%%%%%%%%%%%%%%%%%%%%%%
%%% Compilação condicional em PDF:
%%%%%%%%%%%%%%%%%%%%%%%%%%%%%%%%%%%%%%%%%%%%%%%%%%%%%%%%%%%%%%%%%%%%%%%%%%%%%%%%
% Por: Abrantes Araújo Silva Filho
%      abrantesasf@gmail.com
% URL: https://github.com/abrantesasf/latex


%%%%%%%%%%%%%%%%%%%%%%%%%%%%%%%%%%%%%%%%%%%%%%%%%%%%%%%%%%%%%%%%%%%%%%%%%%%%%%%%
%%% Pacotes para compilaçãocondicional em PDF:
\ifxetex
\else
   \usepackage{ifpdf}
\fi


%%%%%%%%%%%%%%%%%%%%%%%%%%%%%%%%%%%%%%%%%%%%%%%%%%%%%%%%%%%%%%%%%%%%%%%%%%%%%%%%
%%% Configurações de layout da página:
%%%%%%%%%%%%%%%%%%%%%%%%%%%%%%%%%%%%%%%%%%%%%%%%%%%%%%%%%%%%%%%%%%%%%%%%%%%%%%%%
% layout.tex
%
% Arquivo de configuração de packages para o LaTeX, considerando diversas
% classes que constumo utilizar.
%
% Para maiores informações, visite:
%    https://github.com/abrantesasf/latex
%    https://github.com/uvv-computacao/uvvtex2
%
% NÃO ALTERE SE NÃO SOUBER O QUE ESTÁ FAZENDO!
%%%%%%%%%%%%%%%%%%%%%%%%%%%%%%%%%%%%%%%%%%%%%%%%%%%%%%%%%%%%%%%%%%%%%%%%%%%%%%%%


%%%%%%%%%%%%%%%%%%%%%%%%%%%%%%%%%%%%%%%%%%%%%%%%%%%%%%%%%%%%%%%%%%%%%%%%%%%%%%%%
%%% Configuração do tamanho da página, margens, espaçamento entrelinhas e, se
%%% necessário, ativa a indentação dos primeiros parágrafos.

% Para article:
\makeatletter
\@ifclassloaded{article}{
  \usepackage{setspace}
  \ifxetex
    \usepackage[xetex]{geometry}
  \else
    \ifpdf
      \usepackage[pdftex]{geometry}
    \else
      \usepackage[dvips]{geometry}
    \fi
  \fi
  \usepackage{microtype}
  \usepackage{lastpage}
}{}
\makeatother

% Para book:
\makeatletter
\@ifclassloaded{book}{
  \usepackage{setspace}
  \ifxetex
    \usepackage[xetex]{geometry}
  \else
    \ifpdf
      \usepackage[pdftex]{geometry}
    \else
      \usepackage[dvips]{geometry}
    \fi
  \fi
  \usepackage{microtype}
  \usepackage{lastpage}
}{}
\makeatother

% Para exam:
\makeatletter
\@ifclassloaded{exam}{
  \usepackage{setspace}
  \ifxetex
    \usepackage[xetex]{geometry}
  \else
    \ifpdf
      \usepackage[pdftex]{geometry}
    \else
      \usepackage[dvips]{geometry}
    \fi
  \fi
  \usepackage{microtype}
  \usepackage{lastpage}
}{}
\makeatother

% Posiciona figuras e tabelas no topo da página quando adicionadas sozinhas
% em um página em branco. Ver https://github.com/abntex/abntex2/issues/170
\makeatletter
\@ifclassloaded{uvvtex2}{
  % Distância entre o topo da página e o primeiro float:
  \setlength{\@fptop}{5pt}
  % O tamanho do parágrafo é dado por:
  \setlength{\parindent}{1.3cm}
  % Controle do espaçamento entre um parágrafo e outro:
  \setlength{\parskip}{0.2cm}  % tente também \onelineskip
  % Para melhorias de justificação:
  \usepackage{microtype}
  % Para contar págians:
  \usepackage{lastpage}
}{}
\makeatother

% Posiciona figuras e tabelas no topo da página quando adicionadas sozinhas
% em um página em branco. Ver https://github.com/abntex/abntex2/issues/170
\makeatletter
\@ifclassloaded{faesatex2}{
  % Distância entre o topo da página e o primeiro float:
  \setlength{\@fptop}{5pt}
  % O tamanho do parágrafo é dado por:
  \setlength{\parindent}{1.3cm}
  % Controle do espaçamento entre um parágrafo e outro:
  \setlength{\parskip}{0.2cm}  % tente também \onelineskip
  % Para melhorias de justificação:
  \usepackage{microtype}
  % Para contar págians:
  \usepackage{lastpage}
}{}
\makeatother


%%%%%%%%%%%%%%%%%%%%%%%%%%%%%%%%%%%%%%%%%%%%%%%%%%%%%%%%%%%%%%%%%%%%%%%%%%%%%%%%
%%% Configurações para autores:
%%%%%%%%%%%%%%%%%%%%%%%%%%%%%%%%%%%%%%%%%%%%%%%%%%%%%%%%%%%%%%%%%%%%%%%%%%%%%%%%
% Por: Abrantes Araújo Silva Filho
%      abrantesasf@gmail.com
% URL: https://github.com/abrantesasf/latex


%%%%%%%%%%%%%%%%%%%%%%%%%%%%%%%%%%%%%%%%%%%%%%%%%%%%%%%%%%%%%%%%%%%%%%%%%%%%%%%%
%%% Ajustes para os autores e afiliações, em artigos
\makeatletter
\@ifclassloaded{article}{
\usepackage{authblk}
}{}
\makeatother


%%%%%%%%%%%%%%%%%%%%%%%%%%%%%%%%%%%%%%%%%%%%%%%%%%%%%%%%%%%%%%%%%%%%%%%%%%%%%%%%
%%% Cabeçalho e rodapé:
%%% Atenção! É MUITO DIFÍCIL ajustar apropriadamente os running headers
%%% e running footers da primeira vez. Você pode confiar no LaTeX e deixar que
%%% ele faça o ajuste automaticamente ou pode quebrar a cabeça e
%%% tentar melhorar algo que o LaTeX já faz muito bem, mesmo que não fique
%%% exatamente do jeito que você quer. O que você prefere? Se quiser ajustar
%%% manualmente, descomente a linha abaixo e faça as configurações no arquivo
%%% "utils/headings.tex".
%%%%%%%%%%%%%%%%%%%%%%%%%%%%%%%%%%%%%%%%%%%%%%%%%%%%%%%%%%%%%%%%%%%%%%%%%%%%%%%%%
% Por: Abrantes Araújo Silva Filho
%      abrantesasf@gmail.com
% URL: https://github.com/abrantesasf/latex


%%%%%%%%%%%%%%%%%%%%%%%%%%%%%%%%%%%%%%%%%%%%%%%%%%%%%%%%%%%%%%%%%%%%%%%%%%%%%%%%
%%% Configurações de cabeçalho e rodapé:
\usepackage{fancyhdr}
\setlength{\headheight}{1cm}
\setlength{\footskip}{1.5cm}
\renewcommand{\headrulewidth}{0.3pt}
\renewcommand{\footrulewidth}{0.0pt}
\pagestyle{fancy}
\renewcommand{\sectionmark}[1]{%
  \markboth{\uppercase{#1}}{}}
\renewcommand{\subsectionmark}[1]{%
  \markright{\uppercase{\thesubsection \hspace{0.1cm} #1}}{}}
\fancyhead{}
\fancyfoot{}
\newcommand{\diminuifonte}{%
    \fontsize{9pt}{9}\selectfont
}
\newcommand{\aumentafonte}{%
    \fontsize{12}{12}\selectfont
}
% Configura cabeçalho e rodapé para documentos TWOSIDE
\fancyhead[EL]{\textbf{\thepage}}
\fancyhead[EC]{}
\fancyhead[ER]{\diminuifonte \textbf{\leftmark}}
\fancyhead[OR]{\textbf{\thepage}}
\fancyhead[OC]{}
\fancyhead[OL]{\diminuifonte \textbf{\rightmark}}
\fancyfoot[EL,EC,ER,OR,OC,OL]{}
% Configura cabeçalho e rodapé para documentos ONESIDE
%%\lhead{ \fancyplain{}{sup esquerdo} }
%%\chead{ \fancyplain{}{sup centro} }
%%\rhead{ \fancyplain{}{\thesection} }
%%\lfoot{ \fancyplain{}{inf esquerdo} }
%%\cfoot{ \fancyplain{}{inf centro} }
%%\rfoot{ \fancyplain{}{\thepage} }


%%%%%%%%%%%%%%%%%%%%%%%%%%%%%%%%%%%%%%%%%%%%%%%%%%%%%%%%%%%%%%%%%%%%%%%%%%%%%%%%
%%% Configuração para as fontes do documento:
%%% Se você quiser usar fontes próprias ou do sistema, precisará ajustar as
%%% configurações no arquivo "utils/fontes.tex".
%%% ATENÇÃO: por padrão eu utilizo XeLaTeX com fontes proprietárias projetadas
%%% por Matthew Butterick, adquiridas em https://mbtype.com, que não podem ser
%%% distribuídas devido à licença de utilização. Se você usar XeLaTeX, você
%%% DEVERÁ OBRIGATORIAMENTE ajustar as fontes no arquivo "utils/fontes.tex".
%%%%%%%%%%%%%%%%%%%%%%%%%%%%%%%%%%%%%%%%%%%%%%%%%%%%%%%%%%%%%%%%%%%%%%%%%%%%%%%%
% fontes.tex
%
% Arquivo de configuração de packages para uso o LaTeX, conforme minhas
% preferências e modelos pessoais.
%
% Para maiores informações, visite:
%    https://github.com/abrantesasf/latex
%
% NÃO ALTERE SE NÃO SOUBER O QUE ESTÁ FAZENDO!
%%%%%%%%%%%%%%%%%%%%%%%%%%%%%%%%%%%%%%%%%%%%%%%%%%%%%%%%%%%%%%%%%%%%%%%%%%%%%%%%


%%%%%%%%%%%%%%%%%%%%%%%%%%%%%%%%%%%%%%%%%%%%%%%%%%%%%%%%%%%%%%%%%%%%%%%%%%%%%%%%
%%% Configurações de encodings e fontes:
\ifxetex
  % Se usar XeLaTeX, usa fontes específicas:
  %\usepackage[tuenc,no-math]{fontspec}
  \usepackage[tuenc]{fontspec}
  \setmainfont{equity-text-a-regular.otf}[
    Path           = /home/abrantesasf/.local/share/fonts/ ,
    BoldFont       = equity-text-a-bold.otf                ,
    ItalicFont     = equity-text-a-italic.otf              ,
    BoldItalicFont = equity-text-a-bold-italic.otf         ,
    SmallCapsFont  = equity-caps-a-regular.otf]
  \setsansfont{equity-text-b-regular.otf}[
    Path           = /home/abrantesasf/.local/share/fonts/ ,
    BoldFont       = equity-text-b-bold.otf                ,
    ItalicFont     = equity-text-b-italic.otf              ,
    BoldItalicFont = equity-text-b-bold-italic.otf]      
  \setmonofont{courier_prime.ttf}[
    Path           = /home/abrantesasf/.local/share/fonts/ ,
    BoldFont       = courier_prime_bold.ttf                ,
    ItalicFont     = courier_prime_italic.ttf              ,
    BoldItalicFont = courier_prime_bold_italic.ttf]
%  \setmonofont{triplicate-t4-regular.otf}[
%    Path           = /home/abrantesasf/.local/share/fonts/ ,
%    BoldFont       = triplicate-t4-bold.otf                ,
%    ItalicFont     = triplicate-t4-italic.otf              ,
%    BoldItalicFont = triplicate-t4-bold-italic.otf]
  \usepackage{fontawesome5}
\else
  % Se não for XeLaTeX, vai com as normais mesmo:
  \usepackage[T1]{fontenc}
  \usepackage[utf8]{inputenc}
  \usepackage{lmodern}
  % Altera a fonte padrão do documento (nem todas funcionam em modo math):
  %   phv = Helvetica
  %   ptm = Times
  %   ppl = Palatino
  %   pbk = bookman
  %   pag = AdobeAvantGarde
  %   pnc = Adobe NewCenturySchoolBook
  %\renewcommand{\familydefault}{ppl}
\fi


%%%%%%%%%%%%%%%%%%%%%%%%%%%%%%%%%%%%%%%%%%%%%%%%%%%%%%%%%%%%%%%%%%%%%%%%%%%%%%%%
%%% Linguagem e hifenização:
%%%%%%%%%%%%%%%%%%%%%%%%%%%%%%%%%%%%%%%%%%%%%%%%%%%%%%%%%%%%%%%%%%%%%%%%%%%%%%%%
% Por: Abrantes Araújo Silva Filho
%      abrantesasf@gmail.com
% URL: https://github.com/abrantesasf/latex


%%%%%%%%%%%%%%%%%%%%%%%%%%%%%%%%%%%%%%%%%%%%%%%%%%%%%%%%%%%%%%%%%%%%%%%%%%%%%%%%
%%% Configurações de linguagem e hifenização
\usepackage{polyglossia}
\setdefaultlanguage[variant=brazilian]{portuguese}
\setotherlanguage{english}

%%%%%%%%%%%%%%%%%%%%%%%%%%%%%%%%%%%%%%%%%%%%%%%%%%%%%%%%%%%%%%%%%%%%%%%%%%%%%%%%
%%% Hifenização específica de alguams palavras quando a hifenização padrão não
%%% estiver hifenizando corretamente:
\hyphenation{%
Git-Hub
LaTeX
po-ly-glos-sia
trans-gre-din-do
}


%%%%%%%%%%%%%%%%%%%%%%%%%%%%%%%%%%%%%%%%%%%%%%%%%%%%%%%%%%%%%%%%%%%%%%%%%%%%%%%%
%%% Sumário:
%%%%%%%%%%%%%%%%%%%%%%%%%%%%%%%%%%%%%%%%%%%%%%%%%%%%%%%%%%%%%%%%%%%%%%%%%%%%%%%%
% Por: Abrantes Araújo Silva Filho
%      abrantesasf@gmail.com
% URL: https://github.com/abrantesasf/latex


%%%%%%%%%%%%%%%%%%%%%%%%%%%%%%%%%%%%%%%%%%%%%%%%%%%%%%%%%%%%%%%%%%%%%%%%%%%%%%%%
%%% Ajustes de sumário:
\makeatletter
\@ifclassloaded{article}{
   \usepackage[nottoc]{tocbibind}
}{}
\@ifclassloaded{book}{
   \usepackage{tocbibind}
}{}
\makeatother


%%%%%%%%%%%%%%%%%%%%%%%%%%%%%%%%%%%%%%%%%%%%%%%%%%%%%%%%%%%%%%%%%%%%%%%%%%%%%%%%
%%% Matemática:
%%%%%%%%%%%%%%%%%%%%%%%%%%%%%%%%%%%%%%%%%%%%%%%%%%%%%%%%%%%%%%%%%%%%%%%%%%%%%%%%
% matematica.tex
%
% Arquivo de configuração de packages para uso o LaTeX, conforme minhas
% preferências e modelos pessoais.
%
% Para maiores informações, visite:
%    https://github.com/abrantesasf/latex
%
% NÃO ALTERE SE NÃO SOUBER O QUE ESTÁ FAZENDO!
%%%%%%%%%%%%%%%%%%%%%%%%%%%%%%%%%%%%%%%%%%%%%%%%%%%%%%%%%%%%%%%%%%%%%%%%%%%%%%%%


%%%%%%%%%%%%%%%%%%%%%%%%%%%%%%%%%%%%%%%%%%%%%%%%%%%%%%%%%%%%%%%%%%%%%%%%%%%%%%%%
%%% Carrega bibliotecas e símbolos matemáticos, fontes adicionais e configura
%%% algumas outras opções
\usepackage{amsmath}
\usepackage{amssymb}
\usepackage{amsthm}
\usepackage{amsfonts}
\usepackage{mathrsfs}
\usepackage{proof}
\usepackage{siunitx}
  \sisetup{group-separator = {\,}}
  \sisetup{group-digits = {integer}}
  \sisetup{output-decimal-marker = {,}}
  \sisetup{separate-uncertainty}
  \sisetup{multi-part-units = single}
  \sisetup{binary-units = true}
\usepackage{syllogism}
  \setsylpuncpa{}
  \setsylpuncpb{}
  \setsylpuncc{}
  \setsylergosign{}
\usepackage{bm}
\usepackage{cancel}
\usepackage{esvect}
\usepackage{mathtools}
\usepackage{icomma}
\usepackage{nicefrac}
%\usepackage{units}

% Altera separador decimal via comando, se necessário (prefira o siunitx):
%\mathchardef\period=\mathcode`.
%\DeclareMathSymbol{.}{\mathord}{letters}{"3B}


%%%%%%%%%%%%%%%%%%%%%%%%%%%%%%%%%%%%%%%%%%%%%%%%%%%%%%%%%%%%%%%%%%%%%%%%%%%%%%%%
%%% Definições para teoremas, etc.

% Para article:
\makeatletter
\@ifclassloaded{article}{
  \theoremstyle{definition}
  \newtheorem{definicao}{Definição}[section]
  \newtheorem{conjecture}{Conjectura}[section]
  \newtheorem{teorema}{Teorema}[section]
  \newtheorem{lemma}{Lema}[section]
  \newtheorem{corolario}{Corolário}[section]
  \theoremstyle{remark}
  \newtheorem*{nota}{Nota}
  \newtheorem*{observacao}{Observação:}
}{}

% Para book:
\@ifclassloaded{book}{
  \theoremstyle{definition}
  \newtheorem{definicao}{Definição}[chapter]
  \newtheorem{conjecture}{Conjectura}[chapter]
  \newtheorem{teorema}{Teorema}[chapter]
  \newtheorem{lemma}{Lema}[chapter]
  \newtheorem{corolario}{Corolário}[chapter]
  \theoremstyle{remark}
  \newtheorem*{nota}{Nota}
  \newtheorem*{observacao}{Observação:}
}{}
\makeatother



%%%%%%%%%%%%%%%%%%%%%%%%%%%%%%%%%%%%%%%%%%%%%%%%%%%%%%%%%%%%%%%%%%%%%%%%%%%%%%%%
%%% Cores:
%%%%%%%%%%%%%%%%%%%%%%%%%%%%%%%%%%%%%%%%%%%%%%%%%%%%%%%%%%%%%%%%%%%%%%%%%%%%%%%%
% cores.tex
%
% Arquivo de configuração de packages para uso o LaTeX, conforme minhas
% preferências e modelos pessoais.
%
% Para maiores informações, visite:
%    https://github.com/abrantesasf/latex
%
% NÃO ALTERE SE NÃO SOUBER O QUE ESTÁ FAZENDO!
%%%%%%%%%%%%%%%%%%%%%%%%%%%%%%%%%%%%%%%%%%%%%%%%%%%%%%%%%%%%%%%%%%%%%%%%%%%%%%%%


%%%%%%%%%%%%%%%%%%%%%%%%%%%%%%%%%%%%%%%%%%%%%%%%%%%%%%%%%%%%%%%%%%%%%%%%%%%%%%%%
%%% Ativa suporte extendido a cores
%%% Opções comuns:
%%%    usenames (16)
%%%    dvipsnames (64)
%%%    svgnames (150)
%%%    x11names (300)
\makeatletter
\@ifclassloaded{beamer}{
   % Beamer sem opções
  \usepackage{xcolor} 
}{
   % Resto opção svgnames
  \usepackage[svgnames, table]{xcolor}
}
\makeatother



%%%%%%%%%%%%%%%%%%%%%%%%%%%%%%%%%%%%%%%%%%%%%%%%%%%%%%%%%%%%%%%%%%%%%%%%%%%%%%%%
%%% Computação:
%%%%%%%%%%%%%%%%%%%%%%%%%%%%%%%%%%%%%%%%%%%%%%%%%%%%%%%%%%%%%%%%%%%%%%%%%%%%%%%%
% Por: Abrantes Araújo Silva Filho
%      abrantesasf@gmail.com
% URL: https://github.com/abrantesasf/latex


%%%%%%%%%%%%%%%%%%%%%%%%%%%%%%%%%%%%%%%%%%%%%%%%%%%%%%%%%%%%%%%%%%%%%%%%%%%%%%%%
%%% Carrega packages relacionados à computação
\usepackage{algorithm2e}
\usepackage{algorithmicx}
\usepackage{algpseudocode}
\usepackage{listings}
  \lstset{literate=
    {á}{{\'a}}1 {é}{{\'e}}1 {í}{{\'i}}1 {ó}{{\'o}}1 {ú}{{\'u}}1
    {Á}{{\'A}}1 {É}{{\'E}}1 {Í}{{\'I}}1 {Ó}{{\'O}}1 {Ú}{{\'U}}1
    {à}{{\`a}}1 {è}{{\`e}}1 {ì}{{\`i}}1 {ò}{{\`o}}1 {ù}{{\`u}}1
    {À}{{\`A}}1 {È}{{\'E}}1 {Ì}{{\`I}}1 {Ò}{{\`O}}1 {Ù}{{\`U}}1
    {ä}{{\"a}}1 {ë}{{\"e}}1 {ï}{{\"i}}1 {ö}{{\"o}}1 {ü}{{\"u}}1
    {Ä}{{\"A}}1 {Ë}{{\"E}}1 {Ï}{{\"I}}1 {Ö}{{\"O}}1 {Ü}{{\"U}}1
    {â}{{\^a}}1 {ê}{{\^e}}1 {î}{{\^i}}1 {ô}{{\^o}}1 {û}{{\^u}}1
    {Â}{{\^A}}1 {Ê}{{\^E}}1 {Î}{{\^I}}1 {Ô}{{\^O}}1 {Û}{{\^U}}1
    {œ}{{\oe}}1 {Œ}{{\OE}}1 {æ}{{\ae}}1 {Æ}{{\AE}}1 {ß}{{\ss}}1
    {ű}{{\H{u}}}1 {Ű}{{\H{U}}}1 {ő}{{\H{o}}}1 {Ő}{{\H{O}}}1
    {ç}{{\c c}}1 {Ç}{{\c C}}1 {ø}{{\o}}1 {å}{{\r a}}1 {Å}{{\r A}}1
    {€}{{\euro}}1 {£}{{\pounds}}1 {«}{{\guillemotleft}}1
    {»}{{\guillemotright}}1 {ñ}{{\~n}}1 {Ñ}{{\~N}}1 {¿}{{?`}}1
  }
\definecolor{mGreen}{rgb}{0,0.6,0}
\definecolor{mGray}{rgb}{0.5,0.5,0.5}
\definecolor{mPurple}{rgb}{0.58,0,0.82}
\definecolor{backgroundColour}{rgb}{0.95,0.95,0.92}
\lstdefinestyle{CStyle}{
    backgroundcolor=\color{backgroundColour},   
    commentstyle=\color{green}\ttfamily,
    keywordstyle=\color{blue}\ttfamily,
    numberstyle=\tiny\color{mGray},
    stringstyle=\color{red}\ttfamily,
    basicstyle=\ttfamily,
    morecomment=[l][\color{magenta}]{\#},
    breakatwhitespace=false,         
    breaklines=true,                 
    captionpos=b,                    
    keepspaces=true,                 
    numbers=left,                    
    numbersep=5pt,                  
    showspaces=false,                
    showstringspaces=false,
    showtabs=false,                  
    tabsize=4,
    language=C
}


%%%%%%%%%%%%%%%%%%%%%%%%%%%%%%%%%%%%%%%%%%%%%%%%%%%%%%%%%%%%%%%%%%%%%%%%%%%%%%%%
%%% Gráficos:
%%%%%%%%%%%%%%%%%%%%%%%%%%%%%%%%%%%%%%%%%%%%%%%%%%%%%%%%%%%%%%%%%%%%%%%%%%%%%%%%
% graficos.tex
%
% Arquivo de configuração de packages para uso o LaTeX, conforme minhas
% preferências e modelos pessoais.
%
% Para maiores informações, visite:
%    https://github.com/abrantesasf/latex
%
% NÃO ALTERE SE NÃO SOUBER O QUE ESTÁ FAZENDO!
%%%%%%%%%%%%%%%%%%%%%%%%%%%%%%%%%%%%%%%%%%%%%%%%%%%%%%%%%%%%%%%%%%%%%%%%%%%%%%%%


%%%%%%%%%%%%%%%%%%%%%%%%%%%%%%%%%%%%%%%%%%%%%%%%%%%%%%%%%%%%%%%%%%%%%%%%%%%%%%%%
%%% Suporte à importação de gráficos externos
\makeatletter
\@ifclassloaded{beamer}{
  \usepackage{graphicx}
}{
  \ifxetex
    \usepackage{graphicx}
  \else
    \ifpdf
      \usepackage[pdftex]{graphicx}
    \else
      \usepackage[dvips]{graphicx}
    \fi
  \fi
}
\makeatother


%%%%%%%%%%%%%%%%%%%%%%%%%%%%%%%%%%%%%%%%%%%%%%%%%%%%%%%%%%%%%%%%%%%%%%%%%%%%%%%%
%%% Suporte à criação de gráficos proceduralmente no LaTeX:
\usepackage{tikz}
\usetikzlibrary{arrows, automata, backgrounds, matrix, patterns, positioning,
                shapes, shadows, calc, circuits.logic.US}



%%%%%%%%%%%%%%%%%%%%%%%%%%%%%%%%%%%%%%%%%%%%%%%%%%%%%%%%%%%%%%%%%%%%%%%%%%%%%%%%
%%% Ícones:
%%%%%%%%%%%%%%%%%%%%%%%%%%%%%%%%%%%%%%%%%%%%%%%%%%%%%%%%%%%%%%%%%%%%%%%%%%%%%%%%
% Por: Abrantes Araújo Silva Filho
%      abrantesasf@gmail.com
% URL: https://github.com/abrantesasf/latex


%%%%%%%%%%%%%%%%%%%%%%%%%%%%%%%%%%%%%%%%%%%%%%%%%%%%%%%%%%%%%%%%%%%%%%%%%%%%%%%%
%%% Pacotes com ícones, figuras e gráficos extras

% Ícones da Creative Commons:
\usepackage{ccicons}


%%%%%%%%%%%%%%%%%%%%%%%%%%%%%%%%%%%%%%%%%%%%%%%%%%%%%%%%%%%%%%%%%%%%%%%%%%%%%%%%
%%% Blocos:
%%%%%%%%%%%%%%%%%%%%%%%%%%%%%%%%%%%%%%%%%%%%%%%%%%%%%%%%%%%%%%%%%%%%%%%%%%%%%%%%
% Por: Abrantes Araújo Silva Filho
%      abrantesasf@gmail.com
% URL: https://github.com/abrantesasf/latex


%%%%%%%%%%%%%%%%%%%%%%%%%%%%%%%%%%%%%%%%%%%%%%%%%%%%%%%%%%%%%%%%%%%%%%%%%%%%%%%%
%%% Ambiente para boxes em posições arbitrárias no Beamer
\makeatletter
\@ifclassloaded{beamer}{
   \usepackage{textpos}
}{}
\makeatother


%%%%%%%%%%%%%%%%%%%%%%%%%%%%%%%%%%%%%%%%%%%%%%%%%%%%%%%%%%%%%%%%%%%%%%%%%%%%%%%%
%%% Boxes coloridos:
%%%%%%%%%%%%%%%%%%%%%%%%%%%%%%%%%%%%%%%%%%%%%%%%%%%%%%%%%%%%%%%%%%%%%%%%%%%%%%%%
% colorbox.tex
%
% Arquivo de configuração de packages para uso o LaTeX, conforme minhas
% preferências e modelos pessoais.
%
% Para maiores informações, visite:
%    https://github.com/abrantesasf/latex
%
% NÃO ALTERE SE NÃO SOUBER O QUE ESTÁ FAZENDO!
%%%%%%%%%%%%%%%%%%%%%%%%%%%%%%%%%%%%%%%%%%%%%%%%%%%%%%%%%%%%%%%%%%%%%%%%%%%%%%%%


%%%%%%%%%%%%%%%%%%%%%%%%%%%%%%%%%%%%%%%%%%%%%%%%%%%%%%%%%%%%%%%%%%%%%%%%%%%%%%%%
%%% Ativa suporte para boxes coloridos
\usepackage[many]{tcolorbox}


%%%%%%%%%%%%%%%%%%%%%%%%%%%%%%%%%%%%%%%%%%%%%%%%%%%%%%%%%%%%%%%%%%%%%%%%%%%%%%%%
%%% Tabelas:
%%%%%%%%%%%%%%%%%%%%%%%%%%%%%%%%%%%%%%%%%%%%%%%%%%%%%%%%%%%%%%%%%%%%%%%%%%%%%%%%
% tabelas.tex
%
% Arquivo de configuração de packages para o LaTeX, considerando diversas
% classes que constumo utilizar.
%
% Para maiores informações, visite:
%    https://github.com/abrantesasf/latex
%    https://github.com/uvv-computacao/uvvtex2
%
% NÃO ALTERE SE NÃO SOUBER O QUE ESTÁ FAZENDO!
%%%%%%%%%%%%%%%%%%%%%%%%%%%%%%%%%%%%%%%%%%%%%%%%%%%%%%%%%%%%%%%%%%%%%%%%%%%%%%%%


%%%%%%%%%%%%%%%%%%%%%%%%%%%%%%%%%%%%%%%%%%%%%%%%%%%%%%%%%%%%%%%%%%%%%%%%%%%%%%%%
%%% Packages para tabelas
\usepackage{array}
\usepackage{tabularray}
\DefTblrTemplate{contfoot-text}{default}{\footnotesize (continua na próxima página)}
\DefTblrTemplate{conthead-text}{default}{(continuação)}
\ExplSyntaxOn
\RenewDocumentCommand \TblrNote { m }
  {
    \cs_if_exist:NT \hypersetup
      { \ExpTblrTemplate { note-border }{ default } }
    \TblrOverlap
      {
        \__tblr_hyper_link:nn {#1}
          { \kern0.03em \textsuperscript { \rmfamily \UseTblrFont { note-tag } #1 } }
      }
  }
\ExplSyntaxOff
\DeclareTblrTemplate{note-tag}{normal}{
  \textsuperscript{\rmfamily \UseTblrFont{note-tag}\InsertTblrNoteTag}
}
\SetTblrTemplate{note-tag}{normal}
\usepackage{longtable}
\usepackage{tabularx}
%\usepackage{tabu}
\usepackage{lscape}
\usepackage{colortbl}  
\usepackage{booktabs}
\usepackage{multirow}
\usepackage{makecell}



%%%%%%%%%%%%%%%%%%%%%%%%%%%%%%%%%%%%%%%%%%%%%%%%%%%%%%%%%%%%%%%%%%%%%%%%%%%%%%%%
%%% Ambientes de listas:
%%%%%%%%%%%%%%%%%%%%%%%%%%%%%%%%%%%%%%%%%%%%%%%%%%%%%%%%%%%%%%%%%%%%%%%%%%%%%%%%
% Por: Abrantes Araújo Silva Filho
%      abrantesasf@gmail.com
% URL: https://github.com/abrantesasf/latex


%%%%%%%%%%%%%%%%%%%%%%%%%%%%%%%%%%%%%%%%%%%%%%%%%%%%%%%%%%%%%%%%%%%%%%%%%%%%%%%%
%%% Packages para ambientes de listas
\makeatletter
\@ifclassloaded{beamer}{
  \usepackage{listings}
}{
  \usepackage{enumitem}
  \usepackage[ampersand]{easylist}
  \usepackage{fancyvrb, bera}
}
\makeatother



%%%%%%%%%%%%%%%%%%%%%%%%%%%%%%%%%%%%%%%%%%%%%%%%%%%%%%%%%%%%%%%%%%%%%%%%%%%%%%%%
%%% Ambientes para quadros:
%%%%%%%%%%%%%%%%%%%%%%%%%%%%%%%%%%%%%%%%%%%%%%%%%%%%%%%%%%%%%%%%%%%%%%%%%%%%%%%%%
% quadros.tex
%
% Arquivo de configuração de packages para uso com a classe uvvTeX2, para a
% formatação de trabalhos acadêmicos na Universidade Vila Velha (UVV)
% (https://www.uvv.br).
%
% Para maiores informações, visite:
%    https://github.com/uvv-computacao/uvvtex2
%
% NÃO ALTERE NADA AQUI!
%%%%%%%%%%%%%%%%%%%%%%%%%%%%%%%%%%%%%%%%%%%%%%%%%%%%%%%%%%%%%%%%%%%%%%%%%%%%%%%%

% Possibilita criação de quadros
% Ver https://github.com/abntex/abntex2/issues/176

\makeatletter
\@ifclassloaded{uvvtex2}{
  \newcommand{\quadroname}{Quadro}
  \newcommand{\listofquadrosname}{Lista de quadros}

  \newfloat[chapter]{quadro}{loq}{\quadroname}
  \newlistof{listofquadros}{loq}{\listofquadrosname}
  \newlistentry{quadro}{loq}{0}

  \setfloatadjustment{quadro}{\centering}
  \counterwithout{quadro}{chapter}
  \renewcommand{\cftquadroname}{\quadroname\space} 
  \renewcommand*{\cftquadroaftersnum}{\hfill--\hfill}

  \setfloatlocations{quadro}{hbtp}
}{}
\makeatother


%%%%%%%%%%%%%%%%%%%%%%%%%%%%%%%%%%%%%%%%%%%%%%%%%%%%%%%%%%%%%%%%%%%%%%%%%%%%%%%%
%%% Ambientes floats e similares:
%%%%%%%%%%%%%%%%%%%%%%%%%%%%%%%%%%%%%%%%%%%%%%%%%%%%%%%%%%%%%%%%%%%%%%%%%%%%%%%%
% Por: Abrantes Araújo Silva Filho
%      abrantesasf@gmail.com
% URL: https://github.com/abrantesasf/latex


%%%%%%%%%%%%%%%%%%%%%%%%%%%%%%%%%%%%%%%%%%%%%%%%%%%%%%%%%%%%%%%%%%%%%%%%%%%%%%%%
%%% Packages para suporte a ambientes floats, captions, etc.:
\usepackage{float}
\usepackage{wrapfig}
\usepackage{placeins}
\usepackage[justification=centering]{caption}
\usepackage{sidecap}
\usepackage{subcaption}


%%%%%%%%%%%%%%%%%%%%%%%%%%%%%%%%%%%%%%%%%%%%%%%%%%%%%%%%%%%%%%%%%%%%%%%%%%%%%%%%
%%% Notas de rodapé:
%%%%%%%%%%%%%%%%%%%%%%%%%%%%%%%%%%%%%%%%%%%%%%%%%%%%%%%%%%%%%%%%%%%%%%%%%%%%%%%%
% Por: Abrantes Araújo Silva Filho
%      abrantesasf@gmail.com
% URL: https://github.com/abrantesasf/latex


%%%%%%%%%%%%%%%%%%%%%%%%%%%%%%%%%%%%%%%%%%%%%%%%%%%%%%%%%%%%%%%%%%%%%%%%%%%%%%%%%
%%% Carrega pacotes para maior controle de notas de rodapé, inclusive dando
%%% um label para uma nota e fazendo referência cruzada.
%%% NÃO ALTERE a ordem de carregamento dos packages!
\usepackage{footmisc}


%%%%%%%%%%%%%%%%%%%%%%%%%%%%%%%%%%%%%%%%%%%%%%%%%%%%%%%%%%%%%%%%%%%%%%%%%%%%%%%%
%%% Referências cruzadas, links, citações:
%%%%%%%%%%%%%%%%%%%%%%%%%%%%%%%%%%%%%%%%%%%%%%%%%%%%%%%%%%%%%%%%%%%%%%%%%%%%%%%%
% crossrefs.tex
%
% Arquivo de configuração de packages para o LaTeX, considerando diversas
% classes que constumo utilizar.
%
% Para maiores informações, visite:
%    https://github.com/abrantesasf/latex
%    https://github.com/uvv-computacao/uvvtex2
%
% NÃO ALTERE SE NÃO SOUBER O QUE ESTÁ FAZENDO!
%%%%%%%%%%%%%%%%%%%%%%%%%%%%%%%%%%%%%%%%%%%%%%%%%%%%%%%%%%%%%%%%%%%%%%%%%%%%%%%%


%%%%%%%%%%%%%%%%%%%%%%%%%%%%%%%%%%%%%%%%%%%%%%%%%%%%%%%%%%%%%%%%%%%%%%%%%%%%%%%%
%%% Carrega pacotes para referências cruzadas, citações dentro do documento,
%%% links para internet e outros.Configura algumas opções.
%%% NÃO ALTERE a ordem de carregamento dos packages!

% Para article:
\makeatletter
\@ifclassloaded{article}{
  \usepackage{varioref}
  \ifxetex
    \usepackage{hyperref}
    \usepackage[brazilian,nameinlink,noabbrev]{cleveref}
  \else
    \ifpdf
      \usepackage[pdftex]{hyperref}
    \else
      \usepackage[dvips]{hyperref}
    \fi
      \usepackage{cleveref}
  \fi
  \usepackage{xurl}
}{}
\makeatother

% Para book:
\makeatletter
\@ifclassloaded{book}{
  \usepackage{varioref}
  \ifxetex
    \usepackage{hyperref}
    \usepackage[brazilian,nameinlink,noabbrev]{cleveref}
  \else
    \ifpdf
      \usepackage[pdftex]{hyperref}
    \else
      \usepackage[dvips]{hyperref}
    \fi
      \usepackage{cleveref}
  \fi
  \usepackage{xurl}
}{}
\makeatother

% Ajustes para citações com o cleveref:
\newcommand{\SeCref}[1]{%
  ``\hyperref[#1]{\Cref*{#1}: \nameref*{#1}}''%
}


%%%%%%%%%%%%%%%%%%%%%%%%%%%%%%%%%%%%%%%%%%%%%%%%%%%%%%%%%%%%%%%%%%%%%%%%%%%%%%%%
%%% Configurações para os hiperlinks em PDF:
%%%%%%%%%%%%%%%%%%%%%%%%%%%%%%%%%%%%%%%%%%%%%%%%%%%%%%%%%%%%%%%%%%%%%%%%%%%%%%%%
% pdfconfig.tex
%
% Arquivo de configuração de packages para uso o LaTeX, conforme minhas
% preferências e modelos pessoais.
%
% Para maiores informações, visite:
%    https://github.com/abrantesasf/latex
%
% NÃO ALTERE SE NÃO SOUBER O QUE ESTÁ FAZENDO!
%%%%%%%%%%%%%%%%%%%%%%%%%%%%%%%%%%%%%%%%%%%%%%%%%%%%%%%%%%%%%%%%%%%%%%%%%%%%%%%%

% Configurações básicas do PDF
\hypersetup{
   unicode=true,
   pdflang={pt-BR},
   bookmarksopen=false,
   bookmarksnumbered=true,
   bookmarksopenlevel=5,
   pdfdisplaydoctitle=true,
   pdfpagemode=UseOutlines,
   pdfstartview=FitH,
   pdfnewwindow=true,
   %pagebackref=true,
   pdfcreator={TeX Live 2021: XeTeX (0.999993) + LaTeX2e (2021-06-01) + eTeX (2.6) + TeX (3.141592653)},
   pdfproducer={xdvipdfmx (20210318)},
   %pdftitle={\@title}, 
   %pdfauthor={\@author},
   %pdfsubject={\imprimirpreambulo},
   %pdfkeywords={\imprimirpalavraschave, \imprimirkeywords}, 
   bookmarksdepth=5,
   citecolor=cyan,      % cor dos links para as referências bibliográficas
   linkcolor=red,       % cor dos links internos (sumário, figuras, rodapé)
   filecolor=magenta,   % dor dos links para arquivos
   urlcolor=blue,       % cor para links URL
   anchorcolor=black    % cor para o texto âncora (alvo)
}



%%%%%%%%%%%%%%%%%%%%%%%%%%%%%%%%%%%%%%%%%%%%%%%%%%%%%%%%%%%%%%%%%%%%%%%%%%%%%%%%
%%% Glossário e índice remissivo:
%%%%%%%%%%%%%%%%%%%%%%%%%%%%%%%%%%%%%%%%%%%%%%%%%%%%%%%%%%%%%%%%%%%%%%%%%%%%%%%%
% Por: Abrantes Araújo Silva Filho
%      abrantesasf@gmail.com
% URL: https://github.com/abrantesasf/latex


%%%%%%%%%%%%%%%%%%%%%%%%%%%%%%%%%%%%%%%%%%%%%%%%%%%%%%%%%%%%%%%%%%%%%%%%%%%%%%%%
%%% Pacote para índice remissivo
\makeatletter
\@ifclassloaded{beamer}{}
{
   \usepackage{makeidx}
   %\makeindex
}
\makeatother


%%%%%%%%%%%%%%%%%%%%%%%%%%%%%%%%%%%%%%%%%%%%%%%%%%%%%%%%%%%%%%%%%%%%%%%%%%%%%%%%
%%% Pacote para glossário
\makeatletter
\@ifclassloaded{beamer}{}
{
   \usepackage[toc]{glossaries}
   %\newglossary[nlg]{notation}{not}{ntn}{Notação}
   %\newglossaryentry{not:set}{
   %   type = notation,
   %   name = {$\mathbb{N}$},
   %   description = conjunto dos números naturais,
   %   sort = {N}}
   %\makeglossaries
}
\makeatother


%%%%%%%%%%%%%%%%%%%%%%%%%%%%%%%%%%%%%%%%%%%%%%%%%%%%%%%%%%%%%%%%%%%%%%%%%%%%%%%%
%%% Referências bibliográficas:
%%%%%%%%%%%%%%%%%%%%%%%%%%%%%%%%%%%%%%%%%%%%%%%%%%%%%%%%%%%%%%%%%%%%%%%%%%%%%%%%
% Por: Abrantes Araújo Silva Filho
%      abrantesasf@gmail.com
% URL: https://github.com/abrantesasf/latex


%%%%%%%%%%%%%%%%%%%%%%%%%%%%%%%%%%%%%%%%%%%%%%%%%%%%%%%%%%%%%%%%%%%%%%%%%%%%%%%%
%%% Referências bibliográficas
%\usepackage[square,numbers]{natbib}
%\bibliographystyle{plainnat}

\usepackage[alf]{abntex2cite}
%\bibliographystyle{abnt-alf}


%%%%%%%%%%%%%%%%%%%%%%%%%%%%%%%%%%%%%%%%%%%%%%%%%%%%%%%%%%%%%%%%%%%%%%%%%%%%%%%%
%%% Ajustes para a classe EXAM:
%%%%%%%%%%%%%%%%%%%%%%%%%%%%%%%%%%%%%%%%%%%%%%%%%%%%%%%%%%%%%%%%%%%%%%%%%%%%%%%%
% exam.tex
%
% Arquivo de configuração para a classe EXAM, conforme minhas
% preferências e modelos pessoais.
%
% Para maiores informações, visite:
%    https://github.com/abrantesasf/latex
%
% NÃO ALTERE SE NÃO SOUBER O QUE ESTÁ FAZENDO!
%%%%%%%%%%%%%%%%%%%%%%%%%%%%%%%%%%%%%%%%%%%%%%%%%%%%%%%%%%%%%%%%%%%%%%%%%%%%%%%%


% Utilitários para a classe Exam:
\makeatletter
\@ifclassloaded{exam}{%
   % Linhas para respostas
   \setlength\linefillthickness{0.2pt}
   \newcommand{\umalinha}{\fillwithlines{0.25in}}
   \newcommand{\duaslinhas}{\fillwithlines{0.50in}}
   \newcommand{\treslinhas}{\fillwithlines{0.75in}}
   \newcommand{\quatrolinhas}{\fillwithlines{1.00in}}
   \newcommand{\cincolinhas}{\fillwithlines{1.25in}}
   \newcommand{\seislinhas}{\fillwithlines{1.50in}}
   \newcommand{\setelinhas}{\fillwithlines{1.75in}}
   \newcommand{\oitolinhas}{\fillwithlines{2.00in}}
   \newcommand{\novelinhas}{\fillwithlines{2.25in}}
   \newcommand{\dezlinhas}{\fillwithlines{2.50in}}
   
   % Verdadeiro ou Falso
   \newcommand{\vf}[1][{}]{%
      \fillin[#1][0.25in]%
   }
   
   % Título das respostas:
   %\renewcommand{\solutiontitle}{\noindent\textbf{Resposta:}\par\noindent}
   \renewcommand{\solutiontitle}{\noindent}
   %\shadedsolutions
   %\SolutionEmphasis{\itshape}
   
   % Altera a marca de questão correta e a ênfase:
   \checkedchar{$\Rightarrow$}
   %\CorrectChoiceEmphasis{\color{red}}
}{}
\makeatother



%%%%%%%%%%%%%%%%%%%%%%%%%%%%%%%%%%%%%%%%%%%%%%%%%%%%%%%%%%%%%%%%%%%%%%%%%%%%%%%%
%%% Ajustes para textos:
%%%%%%%%%%%%%%%%%%%%%%%%%%%%%%%%%%%%%%%%%%%%%%%%%%%%%%%%%%%%%%%%%%%%%%%%%%%%%%%%
% texto.tex
%
% Arquivo de configuração de packages para o LaTeX, considerando diversas
% classes que constumo utilizar.
%
% Para maiores informações, visite:
%    https://github.com/abrantesasf/latex
%    https://github.com/uvv-computacao/uvvtex2
%
% NÃO ALTERE SE NÃO SOUBER O QUE ESTÁ FAZENDO!
%%%%%%%%%%%%%%%%%%%%%%%%%%%%%%%%%%%%%%%%%%%%%%%%%%%%%%%%%%%%%%%%%%%%%%%%%%%%%%%%


%%%%%%%%%%%%%%%%%%%%%%%%%%%%%%%%%%%%%%%%%%%%%%%%%%%%%%%%%%%%%%%%%%%%%%%%%%%%%%%%
%%% Pacote para texto dummy:
\usepackage{blindtext}


%%%%%%%%%%%%%%%%%%%%%%%%%%%%%%%%%%%%%%%%%%%%%%%%%%%%%%%%%%%%%%%%%%%%%%%%%%%%%%%%
%%% Ambientes, aliases, comandos e customizações em geral para serem
%%% utilizadas nos documentos. Defina aqui qualquer customização específica
%%% para seus documentos!
%%%%%%%%%%%%%%%%%%%%%%%%%%%%%%%%%%%%%%%%%%%%%%%%%%%%%%%%%%%%%%%%%%%%%%%%%%%%%%%%
% customizacoes.tex
%
% Arquivo de configuração de packages para uso o LaTeX, conforme minhas
% preferências e modelos pessoais.
%
% Para maiores informações, visite:
%    https://github.com/abrantesasf/latex
%
% NÃO ALTERE SE NÃO SOUBER O QUE ESTÁ FAZENDO!
%%%%%%%%%%%%%%%%%%%%%%%%%%%%%%%%%%%%%%%%%%%%%%%%%%%%%%%%%%%%%%%%%%%%%%%%%%%%%%%%


%%%%%%%%%%%%%%%%%%%%%%%%%%%%%%%%%%%%%%%%%%%%%%%%%%%%%%%%%%%%%%%%%%%%%%%%%%%%%%%%
%%% Vários comandinhos úteis e outras gambiarras

% Commando para ``italizar´´ palavras em inglês (e outras línguas!):
\newcommand{\ingles}[1]{\textit{#1}}

% Comando para escrever uma função e símbolos em fonte monoespaçada:
\newcommand{\funcao}[1]{\textbf{\texttt{#1}}}
\newcommand{\simbolo}[1]{\texttt{#1}}

% Commando para colocar o espaço correto entre um número e sua unidade:
\newcommand{\unidade}[2]{\ensuremath{#1\,\mathrm{#2}}}
\newcommand{\unidado}[2]{{#1}\,{#2}}

% Produz ordinal masculino ou feminino dependendo do segundo argumento:
\newcommand{\ordinal}[2]{%
#1%
\ifthenelse{\equal{a}{#2}}%
{\textordfeminine}%
{\textordmasculine}}

% Comando para colocar o autor da citação corretamente justificado à direita:
% Modificado de: https://latex.org/forum/viewtopic.php?p=15605&sid=216895679326b531e85caec779e1d710#p15605
%\newcommand*{\fontequote}[2]{%
%   \unskip\hspace*{1em plus 1fill}%
%   \linebreak\hspace*{\fill}\mbox{#1}\vspace{-0.15cm}
%   \linebreak\hspace*{\fill}\mbox{#2}
%}

% Comandos para os quadros padronizados dos livros de matemática:
\newcommand{\abrevideo}[1]{%
\begin{tcolorbox}[colback=yellow!5!white,colframe=yellow!75!black,
       title=\textbf{Vídeo(s) de referência: #1}]}
\newcommand{\fechavideo}{\end{tcolorbox}}

% Comandos para os ícones comuns:
\newcommand{\youtube}{\faYoutube}

% Comandos para álgebra relacional (testados em XeLaTeX):
\def\ojoin{\setbox0=\hbox{$\Join$}%
\rule[0.02ex]{.4em}{.5pt}\llap{\rule[0.95ex]{.4em}{.5pt}}}
\def\leftouterjoin{\mathbin{\ojoin\mkern-7.8mu\Join}}
\def\rightouterjoin{\mathbin{\Join\mkern-7.8mu\ojoin}}
\def\fullouterjoin{\mathbin{\ojoin\mkern-7.8mu\Join\mkern-7.8mu\ojoin}}

\usepackage{needspace}

% ============================================================
% Estrutura jurídica: artigos, cláusulas, parágrafos,
% incisos e alíneas
% ============================================================

\newcounter{artigo}
\newcounter{clausula}
\newcounter{paragrafo}
\newcounter{inciso}
\newcounter{alinea}

% ============================================================
% Medidas
% ============================================================

% Parágrafos numerados
\newlength{\recuoparagrafo}
\newlength{\larguraparagrafo}

\setlength{\recuoparagrafo}{\parindent}
\setlength{\larguraparagrafo}{2.5em}

% Parágrafo único
\newlength{\larguraparagrafounico}

\settowidth{\larguraparagrafounico}{%
  \textbf{Parágrafo único.}\enspace
}

% Incisos
\newlength{\recuoinciso}
\newlength{\largurainciso}

\setlength{\recuoinciso}{%
  \dimexpr\recuoparagrafo+\larguraparagrafo\relax
}
\setlength{\largurainciso}{3em}

% Alíneas
\newlength{\recuoalinea}
\newlength{\larguraalinea}

\setlength{\recuoalinea}{%
  \dimexpr\recuoinciso+\largurainciso\relax
}
\setlength{\larguraalinea}{2em}


% ------------------------------------------------------------
% Formatação dos números
% ------------------------------------------------------------

\renewcommand{\theartigo}{\arabic{artigo}}
\renewcommand{\theclausula}{\arabic{clausula}}
\renewcommand{\theparagrafo}{\arabic{paragrafo}}
\renewcommand{\theinciso}{\Roman{inciso}}
\renewcommand{\thealinea}{\alph{alinea}}

% ------------------------------------------------------------
% Artigo
%
% Uso:
%   \artigo
%   \artigo[Do objeto]
% ------------------------------------------------------------

\newcommand{\artigo}[1][]{%
  \par
  \Needspace{3\baselineskip}%
  \vspace{1.5em}%
  \refstepcounter{artigo}%
  \setcounter{paragrafo}{0}%
  \setcounter{inciso}{0}%
  \setcounter{alinea}{0}%
  \noindent
  \textbf{%
     Artigo~\theartigo
     \ifnum\value{artigo}<10 \textsuperscript{o}\fi%
     \if\relax\detokenize{#1}\relax
     \else\space-- #1%
     \fi
  }%
  \par\nobreak
}

% ------------------------------------------------------------
% Cláusula
%
% Uso:
%   \clausula
%   \clausula[Do objeto]
% ------------------------------------------------------------

\newcommand{\clausula}[1][]{%
  \par
  \Needspace{3\baselineskip}%
  \vspace{1.5em}%
  \refstepcounter{clausula}%
  \setcounter{paragrafo}{0}%
  \setcounter{inciso}{0}%
  \setcounter{alinea}{0}%
  \Needspace{3\baselineskip}
  \par\vspace{1.5em}
  \noindent
  \textbf{%
    Cláusula~\theclausula
    \ifnum\value{clausula}<10 \textsuperscript{a}\fi%
    \if\relax\detokenize{#1}\relax
    \else\space-- #1%
    \fi
  }%
  \par\nobreak
}


% ------------------------------------------------------------
% Parágrafo numerado
%
% Uso:
%   \paragrafo Texto...
%
% Produz:
%   § 1º Texto...
% ------------------------------------------------------------
\setlength{\recuoparagrafo}{\parindent}
\setlength{\larguraparagrafo}{2.5em}

\newcommand{\paragrafo}[1]{%
  \refstepcounter{paragrafo}%
  \setcounter{inciso}{0}%
  \setcounter{alinea}{0}%
  \par
  \begingroup
    \setlength{\parindent}{0pt}%
    \advance\leftskip by
      \dimexpr\recuoparagrafo+\larguraparagrafo\relax
    \noindent
    \llap{%
      \makebox[\larguraparagrafo][l]{%
        \textbf{%
          §~\theparagrafo
          \ifnum\value{paragrafo}<10 º\fi
        }%
      }%
    }%
    \ignorespaces#1%
    \par
  \endgroup
}

% ------------------------------------------------------------
% Parágrafo único
%
% Uso:
%   \paragrafounico Texto...
% ------------------------------------------------------------

\newcommand{\paragrafounico}[1]{%
  \setcounter{inciso}{0}%
  \setcounter{alinea}{0}%
  \par
  \begingroup
    \setlength{\parindent}{0pt}%
    \advance\leftskip by
      \dimexpr\recuoparagrafo+\larguraparagrafounico\relax
    \noindent
    \llap{%
      \makebox[\larguraparagrafounico][l]{%
        \textbf{§ único.}%
      }%
    }%
    \ignorespaces#1%
    \par
  \endgroup
}

% ------------------------------------------------------------
% Inciso
%
% Uso:
%   \inciso Texto...
%
% Produz:
%   I – Texto...
% ------------------------------------------------------------

\newcommand{\inciso}[1]{%
  \refstepcounter{inciso}%
  \setcounter{alinea}{0}%
  \par
  \begingroup
    \setlength{\parindent}{0pt}%
    \advance\leftskip by
      \dimexpr\recuoinciso+\largurainciso\relax
    \noindent
    \llap{%
      \makebox[\largurainciso][l]{%
        \textbf{\theinciso\space--}%
      }%
    }%
    \ignorespaces#1%
    \par
  \endgroup
}

% ------------------------------------------------------------
% Alínea
%
% Uso:
%   \alinea Texto...
%
% Produz:
%   a) Texto...
% ------------------------------------------------------------

\newcommand{\alinea}[1]{%
  \refstepcounter{alinea}%
  \par
  \begingroup
    \setlength{\parindent}{0pt}%
    \advance\leftskip by
      \dimexpr\recuoalinea+\larguraalinea\relax
    \noindent
    \llap{%
      \makebox[\larguraalinea][l]{%
        \textbf{\thealinea)}%
      }%
    }%
    \ignorespaces#1%
    \par
  \endgroup
}

%%%%%%%%%%%%%%%%%%%%%%%%%%%%%%%%%%%%%%%%%%%%%%%%%%%%%%%%%%%%%%%%%%%%%%%%%%%%%%%%
% engenharia.tex
%
% Configurações específicas do livro "Engenharia em 500 Gramas": pacotes de
% engenharia (circuitos, gráficos de dados, diagramas de tempo), caixas
% didáticas semânticas e comandos do livro.
%
% Deve ser carregado DEPOIS de utils/preambulo.tex (que já carrega tikz,
% tcolorbox[many], listings, siunitx, xcolor[svgnames] e cleveref). Para isso,
% em livro.tex, logo após %%%%%%%%%%%%%%%%%%%%%%%%%%%%%%%%%%%%%%%%%%%%%%%%%%%%%%%%%%%%%%%%%%%%%%%%%%%%%%%%
% preambulo.tex
%
% Arquivo de configuração de packages para o LaTeX, considerando diversas
% classes que constumo utilizar.
%
% Para maiores informações, visite:
%    https://github.com/abrantesasf/latex
%    https://github.com/uvv-computacao/uvvtex2
%
% Para maiores informações, visite:
%    https://github.com/abrantesasf/latex
%
% NÃO ALTERE SE NÃO SOUBER O QUE ESTÁ FAZENDO!
%%%%%%%%%%%%%%%%%%%%%%%%%%%%%%%%%%%%%%%%%%%%%%%%%%%%%%%%%%%%%%%%%%%%%%%%%%%%%%%%


%%%%%%%%%%%%%%%%%%%%%%%%%%%%%%%%%%%%%%%%%%%%%%%%%%%%%%%%%%%%%%%%%%%%%%%%%%%%%%%%
%%% Estruturas de controle:
\input{utils/controles}

%%%%%%%%%%%%%%%%%%%%%%%%%%%%%%%%%%%%%%%%%%%%%%%%%%%%%%%%%%%%%%%%%%%%%%%%%%%%%%%%
%%% Compilação condicional em PDF:
\input{utils/pdf}

%%%%%%%%%%%%%%%%%%%%%%%%%%%%%%%%%%%%%%%%%%%%%%%%%%%%%%%%%%%%%%%%%%%%%%%%%%%%%%%%
%%% Configurações de layout da página:
\input{utils/layout}

%%%%%%%%%%%%%%%%%%%%%%%%%%%%%%%%%%%%%%%%%%%%%%%%%%%%%%%%%%%%%%%%%%%%%%%%%%%%%%%%
%%% Configurações para autores:
\input{utils/autores}

%%%%%%%%%%%%%%%%%%%%%%%%%%%%%%%%%%%%%%%%%%%%%%%%%%%%%%%%%%%%%%%%%%%%%%%%%%%%%%%%
%%% Cabeçalho e rodapé:
%%% Atenção! É MUITO DIFÍCIL ajustar apropriadamente os running headers
%%% e running footers da primeira vez. Você pode confiar no LaTeX e deixar que
%%% ele faça o ajuste automaticamente ou pode quebrar a cabeça e
%%% tentar melhorar algo que o LaTeX já faz muito bem, mesmo que não fique
%%% exatamente do jeito que você quer. O que você prefere? Se quiser ajustar
%%% manualmente, descomente a linha abaixo e faça as configurações no arquivo
%%% "utils/headings.tex".
%\input{utils/headings}

%%%%%%%%%%%%%%%%%%%%%%%%%%%%%%%%%%%%%%%%%%%%%%%%%%%%%%%%%%%%%%%%%%%%%%%%%%%%%%%%
%%% Configuração para as fontes do documento:
%%% Se você quiser usar fontes próprias ou do sistema, precisará ajustar as
%%% configurações no arquivo "utils/fontes.tex".
%%% ATENÇÃO: por padrão eu utilizo XeLaTeX com fontes proprietárias projetadas
%%% por Matthew Butterick, adquiridas em https://mbtype.com, que não podem ser
%%% distribuídas devido à licença de utilização. Se você usar XeLaTeX, você
%%% DEVERÁ OBRIGATORIAMENTE ajustar as fontes no arquivo "utils/fontes.tex".
\input{utils/fontes}

%%%%%%%%%%%%%%%%%%%%%%%%%%%%%%%%%%%%%%%%%%%%%%%%%%%%%%%%%%%%%%%%%%%%%%%%%%%%%%%%
%%% Linguagem e hifenização:
\input{utils/linguagem}

%%%%%%%%%%%%%%%%%%%%%%%%%%%%%%%%%%%%%%%%%%%%%%%%%%%%%%%%%%%%%%%%%%%%%%%%%%%%%%%%
%%% Sumário:
\input{utils/toc}

%%%%%%%%%%%%%%%%%%%%%%%%%%%%%%%%%%%%%%%%%%%%%%%%%%%%%%%%%%%%%%%%%%%%%%%%%%%%%%%%
%%% Matemática:
\input{utils/matematica}

%%%%%%%%%%%%%%%%%%%%%%%%%%%%%%%%%%%%%%%%%%%%%%%%%%%%%%%%%%%%%%%%%%%%%%%%%%%%%%%%
%%% Cores:
\input{utils/cores}

%%%%%%%%%%%%%%%%%%%%%%%%%%%%%%%%%%%%%%%%%%%%%%%%%%%%%%%%%%%%%%%%%%%%%%%%%%%%%%%%
%%% Computação:
\input{utils/computacao}

%%%%%%%%%%%%%%%%%%%%%%%%%%%%%%%%%%%%%%%%%%%%%%%%%%%%%%%%%%%%%%%%%%%%%%%%%%%%%%%%
%%% Gráficos:
\input{utils/graficos}

%%%%%%%%%%%%%%%%%%%%%%%%%%%%%%%%%%%%%%%%%%%%%%%%%%%%%%%%%%%%%%%%%%%%%%%%%%%%%%%%
%%% Ícones:
\input{utils/icones}

%%%%%%%%%%%%%%%%%%%%%%%%%%%%%%%%%%%%%%%%%%%%%%%%%%%%%%%%%%%%%%%%%%%%%%%%%%%%%%%%
%%% Blocos:
\input{utils/blocos}

%%%%%%%%%%%%%%%%%%%%%%%%%%%%%%%%%%%%%%%%%%%%%%%%%%%%%%%%%%%%%%%%%%%%%%%%%%%%%%%%
%%% Boxes coloridos:
\input{utils/colorbox}

%%%%%%%%%%%%%%%%%%%%%%%%%%%%%%%%%%%%%%%%%%%%%%%%%%%%%%%%%%%%%%%%%%%%%%%%%%%%%%%%
%%% Tabelas:
\input{utils/tabelas}

%%%%%%%%%%%%%%%%%%%%%%%%%%%%%%%%%%%%%%%%%%%%%%%%%%%%%%%%%%%%%%%%%%%%%%%%%%%%%%%%
%%% Ambientes de listas:
\input{utils/listas}

%%%%%%%%%%%%%%%%%%%%%%%%%%%%%%%%%%%%%%%%%%%%%%%%%%%%%%%%%%%%%%%%%%%%%%%%%%%%%%%%
%%% Ambientes para quadros:
%\input{utils/quadros}

%%%%%%%%%%%%%%%%%%%%%%%%%%%%%%%%%%%%%%%%%%%%%%%%%%%%%%%%%%%%%%%%%%%%%%%%%%%%%%%%
%%% Ambientes floats e similares:
\input{utils/floats}

%%%%%%%%%%%%%%%%%%%%%%%%%%%%%%%%%%%%%%%%%%%%%%%%%%%%%%%%%%%%%%%%%%%%%%%%%%%%%%%%
%%% Notas de rodapé:
\input{utils/footnotes}

%%%%%%%%%%%%%%%%%%%%%%%%%%%%%%%%%%%%%%%%%%%%%%%%%%%%%%%%%%%%%%%%%%%%%%%%%%%%%%%%
%%% Referências cruzadas, links, citações:
\input{utils/crossrefs}

%%%%%%%%%%%%%%%%%%%%%%%%%%%%%%%%%%%%%%%%%%%%%%%%%%%%%%%%%%%%%%%%%%%%%%%%%%%%%%%%
%%% Configurações para os hiperlinks em PDF:
\input{utils/pdfconfig}

%%%%%%%%%%%%%%%%%%%%%%%%%%%%%%%%%%%%%%%%%%%%%%%%%%%%%%%%%%%%%%%%%%%%%%%%%%%%%%%%
%%% Glossário e índice remissivo:
\input{utils/glossindex}

%%%%%%%%%%%%%%%%%%%%%%%%%%%%%%%%%%%%%%%%%%%%%%%%%%%%%%%%%%%%%%%%%%%%%%%%%%%%%%%%
%%% Referências bibliográficas:
\input{utils/bibliografia}

%%%%%%%%%%%%%%%%%%%%%%%%%%%%%%%%%%%%%%%%%%%%%%%%%%%%%%%%%%%%%%%%%%%%%%%%%%%%%%%%
%%% Ajustes para a classe EXAM:
\input{utils/exam}

%%%%%%%%%%%%%%%%%%%%%%%%%%%%%%%%%%%%%%%%%%%%%%%%%%%%%%%%%%%%%%%%%%%%%%%%%%%%%%%%
%%% Ajustes para textos:
\input{utils/texto}

%%%%%%%%%%%%%%%%%%%%%%%%%%%%%%%%%%%%%%%%%%%%%%%%%%%%%%%%%%%%%%%%%%%%%%%%%%%%%%%%
%%% Ambientes, aliases, comandos e customizações em geral para serem
%%% utilizadas nos documentos. Defina aqui qualquer customização específica
%%% para seus documentos!
\input{utils/customizacoes}
\input{utils/juridico}
\input{utils/engenharia}
, acrescente:
%    %%%%%%%%%%%%%%%%%%%%%%%%%%%%%%%%%%%%%%%%%%%%%%%%%%%%%%%%%%%%%%%%%%%%%%%%%%%%%%%%
% engenharia.tex
%
% Configurações específicas do livro "Engenharia em 500 Gramas": pacotes de
% engenharia (circuitos, gráficos de dados, diagramas de tempo), caixas
% didáticas semânticas e comandos do livro.
%
% Deve ser carregado DEPOIS de utils/preambulo.tex (que já carrega tikz,
% tcolorbox[many], listings, siunitx, xcolor[svgnames] e cleveref). Para isso,
% em livro.tex, logo após \input{utils/preambulo.tex}, acrescente:
%    \input{utils/engenharia.tex}
%
% Nada aqui altera arquivos do modelo. Os capítulos devem usar SOMENTE os
% ambientes semânticos abaixo; a aparência é definida neste arquivo.
%%%%%%%%%%%%%%%%%%%%%%%%%%%%%%%%%%%%%%%%%%%%%%%%%%%%%%%%%%%%%%%%%%%%%%%%%%%%%%%%


%%%%%%%%%%%%%%%%%%%%%%%%%%%%%%%%%%%%%%%%%%%%%%%%%%%%%%%%%%%%%%%%%%%%%%%%%%%%%%%%
%%% Circuitos: esquemas elétricos. Estilo europeu (resistor retangular), o
%%% usado no Brasil. A opção siunitx faz os rótulos usarem \SI, herdando a
%%% vírgula decimal configurada em utils/matematica.tex.
\usepackage[siunitx, european]{circuitikz}
\ctikzset{bipoles/length=1.2cm}

%%%%%%%%%%%%%%%%%%%%%%%%%%%%%%%%%%%%%%%%%%%%%%%%%%%%%%%%%%%%%%%%%%%%%%%%%%%%%%%%
%%% Gráficos de dados: curvas dos ensaios (motor, sensores, bateria).
%%% Vírgula decimal e espaço fino como separador de milhar, como no siunitx.
\usepackage{pgfplots}
\usepackage{pgfplotstable}
\pgfplotsset{
  compat=1.18,
  /pgf/number format/.cd, use comma, 1000 sep={\,},
  every axis/.append style={grid=major, grid style={dotted, gray!60}},
}

%%%%%%%%%%%%%%%%%%%%%%%%%%%%%%%%%%%%%%%%%%%%%%%%%%%%%%%%%%%%%%%%%%%%%%%%%%%%%%%%
%%% Bibliotecas TikZ para diagramas de blocos, mecânica e o robô no dohyô.
%%% (utils/graficos.tex já carrega arrows, automata, backgrounds, matrix,
%%% patterns, positioning, shapes, shadows, calc e circuits.logic.US.)
\usetikzlibrary{arrows.meta, decorations.pathmorphing, decorations.markings,
                fit, chains, angles}
% Não carregue a biblioteca "quotes": a sintaxe ["rótulo"] em \draw e \pic
% falha neste modelo (o caractere " é ativo por causa do babel/polyglossia).
% Use "node[above]{$\theta$}" e "pic text=$\alpha$" no lugar. Alternativa:
% \shorthandoff{"} antes do tikzpicture (e a biblioteca quotes), mas isso
% desliga localmente os atalhos "" e "- do português.

%%%%%%%%%%%%%%%%%%%%%%%%%%%%%%%%%%%%%%%%%%%%%%%%%%%%%%%%%%%%%%%%%%%%%%%%%%%%%%%%
%%% Diagramas de tempo (PWM, pulso do HC-SR04, protocolo IR) e formatos de
%%% dados/registradores.
\usepackage{tikz-timing}
\usepackage{bytefield}

%%%%%%%%%%%%%%%%%%%%%%%%%%%%%%%%%%%%%%%%%%%%%%%%%%%%%%%%%%%%%%%%%%%%%%%%%%%%%%%%
%%% QR codes para os vídeos complementares.
\usepackage{qrcode}

%%% Árvores (classificações, hierarquias).
\usepackage[edges]{forest}


%%%%%%%%%%%%%%%%%%%%%%%%%%%%%%%%%%%%%%%%%%%%%%%%%%%%%%%%%%%%%%%%%%%%%%%%%%%%%%%%
%%% Tons de cinza usados pelas caixas (edição em preto e branco). Para uma
%%% edição colorida, basta redefinir estas cores.
\providecolor{cinzaClaro}{gray}{0.90}
\providecolor{cinzaMedio}{gray}{0.55}
\providecolor{cinzaEscuro}{gray}{0.25}


%%%%%%%%%%%%%%%%%%%%%%%%%%%%%%%%%%%%%%%%%%%%%%%%%%%%%%%%%%%%%%%%%%%%%%%%%%%%%%%%
%%% Comandos do livro

% Plataformas: usar sempre os comandos, nunca digitar à mão.
\newcommand{\RI}{R1}
\newcommand{\RII}{R2}

% Marcas de rascunho: visíveis enquanto \rascunhotrue; somem na edição final
% (troque para \rascunhofalse).
\newif\ifrascunho
\rascunhotrue
\newcommand{\TODO}[1]{\ifrascunho{\sffamily\footnotesize\color{cinzaEscuro}%
  [\textbf{A FAZER:} #1]}\fi}

% \codigo[opções listings]{arquivo}{marcador}
% Extrai do arquivo-fonte REAL o trecho delimitado por
%    // @livro: <marcador>
%    ...
%    // @livro: <marcador>_fim
% Assim o livro nunca diverge do firmware: se o código muda, o trecho muda
% junto. Usa o estilo CStyle definido em utils/computacao.tex.
%
% Limitações do listings (testadas no XeLaTeX):
%  1. Caracteres acima de U+00FF no arquivo-fonte (travessão, símbolo de
%     aproximação, emojis) geram linhas espúrias. Acentos do português são
%     seguros.
%  2. Comentário de bloco /* ... */ de várias linhas dentro de um trecho
%     marcado faz o marcador final ser ignorado. Dentro de trechos marcados,
%     use //.
\newcommand{\codigo}[3][]{%
  \lstinputlisting[style=CStyle,
    linerange={#3-#3_fim},
    rangeprefix=//\ @livro:\ , rangesuffix=,
    includerangemarker=false, #1]{#2}}


%%%%%%%%%%%%%%%%%%%%%%%%%%%%%%%%%%%%%%%%%%%%%%%%%%%%%%%%%%%%%%%%%%%%%%%%%%%%%%%%
%%% Caixas didáticas semânticas
% Recuo de parágrafo dentro das caixas: o mesmo do texto corrido. O tcolorbox
% zera \parindent no conteúdo; o primeiro parágrafo fica sem recuo.
\newlength{\recuocaixa}
\AtBeginDocument{\setlength{\recuocaixa}{\parindent}}
\tcbset{
  e500base/.style={
    enhanced, breakable, sharp corners,
    colback=white, colframe=black, coltitle=black,
    boxrule=0.4pt, left=2mm, right=2mm, top=1.2mm, bottom=1.2mm,
    fonttitle=\sffamily\bfseries\small,
    before skip=0.8\baselineskip, after skip=0.8\baselineskip,
    before upper={\setlength{\parindent}{\recuocaixa}\noindent},
  },
  % Título na primeira linha do conteúdo. Substitui "attach title to upper",
  % que redefine before upper e perderia o recuo acima.
  e500tituloInterno/.style={detach title,
    before upper={\setlength{\parindent}{\recuocaixa}%
      \noindent\tcbtitle\par\smallskip\noindent}},
}

% Fundamento: teoria transferível, independente do robô.
\newtcolorbox{fundamento}[1][]{e500base,
  frame hidden, borderline west={2.5pt}{0pt}{black},
  title={Fundamento\ifx&#1&\else\ --- #1\fi},
  e500tituloInterno}

% Modelo: desenvolvimento matemático.
\newtcolorbox{modelo}[1][]{e500base,
  colbacktitle=cinzaClaro, titlerule=0.4pt,
  title={Modelo\ifx&#1&\else\ --- #1\fi}}

% Projeto R1 / Projeto R2: decisão aplicada a cada plataforma.
\newtcolorbox{projetoRI}[1][]{e500base,
  colbacktitle=black, coltitle=white,
  title={Projeto R1\ifx&#1&\else\ --- #1\fi}}
\newtcolorbox{projetoRII}[1][]{e500base,
  colbacktitle=cinzaEscuro, coltitle=white, boxrule=0.8pt,
  title={Projeto R2\ifx&#1&\else\ --- #1\fi}}

% Laboratório: experimento obrigatório do capítulo.
\newtcolorbox{laboratorio}[1][]{e500base,
  frame hidden, borderline north={1.2pt}{0pt}{black},
  borderline south={1.2pt}{0pt}{black},
  title={Laboratório\ifx&#1&\else\ --- #1\fi},
  e500tituloInterno}

% Quando a engenharia dá errado: falha real, documentada.
\newtcolorbox{falhareal}[1][]{e500base,
  frame hidden, borderline west={5pt}{0pt}{black},
  title={Quando a engenharia dá errado\ifx&#1&\else\ --- #1\fi},
  e500tituloInterno}

% Regra -> requisito: regulamento traduzido em requisito verificável.
\newtcolorbox{regra}[1][]{e500base,
  frame hidden,
  borderline={0.6pt}{0pt}{black,dashed},
  title={Regra $\rightarrow$ requisito\ifx&#1&\else\ --- #1\fi},
  e500tituloInterno}

% Competição: impacto da decisão no desempenho.
\newtcolorbox{competicao}[1][]{e500base,
  colframe=cinzaMedio, boxrule=0.8pt,
  title={Competição\ifx&#1&\else\ --- #1\fi},
  e500tituloInterno}

% Avançado: conteúdo opcional; pode ser pulado sem perda de continuidade.
\newtcolorbox{avancado}[1][]{e500base,
  colframe=cinzaMedio, boxrule=0.4pt,
  borderline={0.4pt}{1.2pt}{cinzaMedio},
  title={Avançado (opcional)\ifx&#1&\else\ --- #1\fi},
  e500tituloInterno}

% Atenção: cuidado prático, risco de dano ou erro comum de uso.
\newtcolorbox{atencao}[1][]{e500base,
  boxrule=1.2pt,
  title={Atenção\ifx&#1&\else\ --- #1\fi},
  e500tituloInterno}

% Vídeo complementar: nunca indispensável para entender o texto.
%    \video{V05.1}{Caracterização do motor em bancada}{https://...}
\newcommand{\video}[3]{%
  \begin{tcolorbox}[e500base, colframe=cinzaMedio, boxrule=0.4pt,
                    sidebyside, lefthand width=18mm, sidebyside gap=3mm]
    \qrcode[height=16mm]{#3}
    \tcblower
    {\sffamily\bfseries\small Vídeo complementar #1}\par
    #2\par
    {\footnotesize\ttfamily\url{#3}}
  \end{tcolorbox}}

%
% Nada aqui altera arquivos do modelo. Os capítulos devem usar SOMENTE os
% ambientes semânticos abaixo; a aparência é definida neste arquivo.
%%%%%%%%%%%%%%%%%%%%%%%%%%%%%%%%%%%%%%%%%%%%%%%%%%%%%%%%%%%%%%%%%%%%%%%%%%%%%%%%


%%%%%%%%%%%%%%%%%%%%%%%%%%%%%%%%%%%%%%%%%%%%%%%%%%%%%%%%%%%%%%%%%%%%%%%%%%%%%%%%
%%% Circuitos: esquemas elétricos. Estilo europeu (resistor retangular), o
%%% usado no Brasil. A opção siunitx faz os rótulos usarem \SI, herdando a
%%% vírgula decimal configurada em utils/matematica.tex.
\usepackage[siunitx, european]{circuitikz}
\ctikzset{bipoles/length=1.2cm}

%%%%%%%%%%%%%%%%%%%%%%%%%%%%%%%%%%%%%%%%%%%%%%%%%%%%%%%%%%%%%%%%%%%%%%%%%%%%%%%%
%%% Gráficos de dados: curvas dos ensaios (motor, sensores, bateria).
%%% Vírgula decimal e espaço fino como separador de milhar, como no siunitx.
\usepackage{pgfplots}
\usepackage{pgfplotstable}
\pgfplotsset{
  compat=1.18,
  /pgf/number format/.cd, use comma, 1000 sep={\,},
  every axis/.append style={grid=major, grid style={dotted, gray!60}},
}

%%%%%%%%%%%%%%%%%%%%%%%%%%%%%%%%%%%%%%%%%%%%%%%%%%%%%%%%%%%%%%%%%%%%%%%%%%%%%%%%
%%% Bibliotecas TikZ para diagramas de blocos, mecânica e o robô no dohyô.
%%% (utils/graficos.tex já carrega arrows, automata, backgrounds, matrix,
%%% patterns, positioning, shapes, shadows, calc e circuits.logic.US.)
\usetikzlibrary{arrows.meta, decorations.pathmorphing, decorations.markings,
                fit, chains, angles}
% Não carregue a biblioteca "quotes": a sintaxe ["rótulo"] em \draw e \pic
% falha neste modelo (o caractere " é ativo por causa do babel/polyglossia).
% Use "node[above]{$\theta$}" e "pic text=$\alpha$" no lugar. Alternativa:
% \shorthandoff{"} antes do tikzpicture (e a biblioteca quotes), mas isso
% desliga localmente os atalhos "" e "- do português.

%%%%%%%%%%%%%%%%%%%%%%%%%%%%%%%%%%%%%%%%%%%%%%%%%%%%%%%%%%%%%%%%%%%%%%%%%%%%%%%%
%%% Diagramas de tempo (PWM, pulso do HC-SR04, protocolo IR) e formatos de
%%% dados/registradores.
\usepackage{tikz-timing}
\usepackage{bytefield}

%%%%%%%%%%%%%%%%%%%%%%%%%%%%%%%%%%%%%%%%%%%%%%%%%%%%%%%%%%%%%%%%%%%%%%%%%%%%%%%%
%%% QR codes para os vídeos complementares.
\usepackage{qrcode}

%%% Árvores (classificações, hierarquias).
\usepackage[edges]{forest}


%%%%%%%%%%%%%%%%%%%%%%%%%%%%%%%%%%%%%%%%%%%%%%%%%%%%%%%%%%%%%%%%%%%%%%%%%%%%%%%%
%%% Tons de cinza usados pelas caixas (edição em preto e branco). Para uma
%%% edição colorida, basta redefinir estas cores.
\providecolor{cinzaClaro}{gray}{0.90}
\providecolor{cinzaMedio}{gray}{0.55}
\providecolor{cinzaEscuro}{gray}{0.25}


%%%%%%%%%%%%%%%%%%%%%%%%%%%%%%%%%%%%%%%%%%%%%%%%%%%%%%%%%%%%%%%%%%%%%%%%%%%%%%%%
%%% Comandos do livro

% Plataformas: usar sempre os comandos, nunca digitar à mão.
\newcommand{\RI}{R1}
\newcommand{\RII}{R2}

% Marcas de rascunho: visíveis enquanto \rascunhotrue; somem na edição final
% (troque para \rascunhofalse).
\newif\ifrascunho
\rascunhotrue
\newcommand{\TODO}[1]{\ifrascunho{\sffamily\footnotesize\color{cinzaEscuro}%
  [\textbf{A FAZER:} #1]}\fi}

% \codigo[opções listings]{arquivo}{marcador}
% Extrai do arquivo-fonte REAL o trecho delimitado por
%    // @livro: <marcador>
%    ...
%    // @livro: <marcador>_fim
% Assim o livro nunca diverge do firmware: se o código muda, o trecho muda
% junto. Usa o estilo CStyle definido em utils/computacao.tex.
%
% Limitações do listings (testadas no XeLaTeX):
%  1. Caracteres acima de U+00FF no arquivo-fonte (travessão, símbolo de
%     aproximação, emojis) geram linhas espúrias. Acentos do português são
%     seguros.
%  2. Comentário de bloco /* ... */ de várias linhas dentro de um trecho
%     marcado faz o marcador final ser ignorado. Dentro de trechos marcados,
%     use //.
\newcommand{\codigo}[3][]{%
  \lstinputlisting[style=CStyle,
    linerange={#3-#3_fim},
    rangeprefix=//\ @livro:\ , rangesuffix=,
    includerangemarker=false, #1]{#2}}


%%%%%%%%%%%%%%%%%%%%%%%%%%%%%%%%%%%%%%%%%%%%%%%%%%%%%%%%%%%%%%%%%%%%%%%%%%%%%%%%
%%% Caixas didáticas semânticas
% Recuo de parágrafo dentro das caixas: o mesmo do texto corrido. O tcolorbox
% zera \parindent no conteúdo; o primeiro parágrafo fica sem recuo.
\newlength{\recuocaixa}
\AtBeginDocument{\setlength{\recuocaixa}{\parindent}}
\tcbset{
  e500base/.style={
    enhanced, breakable, sharp corners,
    colback=white, colframe=black, coltitle=black,
    boxrule=0.4pt, left=2mm, right=2mm, top=1.2mm, bottom=1.2mm,
    fonttitle=\sffamily\bfseries\small,
    before skip=0.8\baselineskip, after skip=0.8\baselineskip,
    before upper={\setlength{\parindent}{\recuocaixa}\noindent},
  },
  % Título na primeira linha do conteúdo. Substitui "attach title to upper",
  % que redefine before upper e perderia o recuo acima.
  e500tituloInterno/.style={detach title,
    before upper={\setlength{\parindent}{\recuocaixa}%
      \noindent\tcbtitle\par\smallskip\noindent}},
}

% Fundamento: teoria transferível, independente do robô.
\newtcolorbox{fundamento}[1][]{e500base,
  frame hidden, borderline west={2.5pt}{0pt}{black},
  title={Fundamento\ifx&#1&\else\ --- #1\fi},
  e500tituloInterno}

% Modelo: desenvolvimento matemático.
\newtcolorbox{modelo}[1][]{e500base,
  colbacktitle=cinzaClaro, titlerule=0.4pt,
  title={Modelo\ifx&#1&\else\ --- #1\fi}}

% Projeto R1 / Projeto R2: decisão aplicada a cada plataforma.
\newtcolorbox{projetoRI}[1][]{e500base,
  colbacktitle=black, coltitle=white,
  title={Projeto R1\ifx&#1&\else\ --- #1\fi}}
\newtcolorbox{projetoRII}[1][]{e500base,
  colbacktitle=cinzaEscuro, coltitle=white, boxrule=0.8pt,
  title={Projeto R2\ifx&#1&\else\ --- #1\fi}}

% Laboratório: experimento obrigatório do capítulo.
\newtcolorbox{laboratorio}[1][]{e500base,
  frame hidden, borderline north={1.2pt}{0pt}{black},
  borderline south={1.2pt}{0pt}{black},
  title={Laboratório\ifx&#1&\else\ --- #1\fi},
  e500tituloInterno}

% Quando a engenharia dá errado: falha real, documentada.
\newtcolorbox{falhareal}[1][]{e500base,
  frame hidden, borderline west={5pt}{0pt}{black},
  title={Quando a engenharia dá errado\ifx&#1&\else\ --- #1\fi},
  e500tituloInterno}

% Regra -> requisito: regulamento traduzido em requisito verificável.
\newtcolorbox{regra}[1][]{e500base,
  frame hidden,
  borderline={0.6pt}{0pt}{black,dashed},
  title={Regra $\rightarrow$ requisito\ifx&#1&\else\ --- #1\fi},
  e500tituloInterno}

% Competição: impacto da decisão no desempenho.
\newtcolorbox{competicao}[1][]{e500base,
  colframe=cinzaMedio, boxrule=0.8pt,
  title={Competição\ifx&#1&\else\ --- #1\fi},
  e500tituloInterno}

% Avançado: conteúdo opcional; pode ser pulado sem perda de continuidade.
\newtcolorbox{avancado}[1][]{e500base,
  colframe=cinzaMedio, boxrule=0.4pt,
  borderline={0.4pt}{1.2pt}{cinzaMedio},
  title={Avançado (opcional)\ifx&#1&\else\ --- #1\fi},
  e500tituloInterno}

% Atenção: cuidado prático, risco de dano ou erro comum de uso.
\newtcolorbox{atencao}[1][]{e500base,
  boxrule=1.2pt,
  title={Atenção\ifx&#1&\else\ --- #1\fi},
  e500tituloInterno}

% Vídeo complementar: nunca indispensável para entender o texto.
%    \video{V05.1}{Caracterização do motor em bancada}{https://...}
\newcommand{\video}[3]{%
  \begin{tcolorbox}[e500base, colframe=cinzaMedio, boxrule=0.4pt,
                    sidebyside, lefthand width=18mm, sidebyside gap=3mm]
    \qrcode[height=16mm]{#3}
    \tcblower
    {\sffamily\bfseries\small Vídeo complementar #1}\par
    #2\par
    {\footnotesize\ttfamily\url{#3}}
  \end{tcolorbox}}

, acrescente:
%    %%%%%%%%%%%%%%%%%%%%%%%%%%%%%%%%%%%%%%%%%%%%%%%%%%%%%%%%%%%%%%%%%%%%%%%%%%%%%%%%
% engenharia.tex
%
% Configurações específicas do livro "Engenharia em 500 Gramas": pacotes de
% engenharia (circuitos, gráficos de dados, diagramas de tempo), caixas
% didáticas semânticas e comandos do livro.
%
% Deve ser carregado DEPOIS de utils/preambulo.tex (que já carrega tikz,
% tcolorbox[many], listings, siunitx, xcolor[svgnames] e cleveref). Para isso,
% em livro.tex, logo após %%%%%%%%%%%%%%%%%%%%%%%%%%%%%%%%%%%%%%%%%%%%%%%%%%%%%%%%%%%%%%%%%%%%%%%%%%%%%%%%
% preambulo.tex
%
% Arquivo de configuração de packages para o LaTeX, considerando diversas
% classes que constumo utilizar.
%
% Para maiores informações, visite:
%    https://github.com/abrantesasf/latex
%    https://github.com/uvv-computacao/uvvtex2
%
% Para maiores informações, visite:
%    https://github.com/abrantesasf/latex
%
% NÃO ALTERE SE NÃO SOUBER O QUE ESTÁ FAZENDO!
%%%%%%%%%%%%%%%%%%%%%%%%%%%%%%%%%%%%%%%%%%%%%%%%%%%%%%%%%%%%%%%%%%%%%%%%%%%%%%%%


%%%%%%%%%%%%%%%%%%%%%%%%%%%%%%%%%%%%%%%%%%%%%%%%%%%%%%%%%%%%%%%%%%%%%%%%%%%%%%%%
%%% Estruturas de controle:
%%%%%%%%%%%%%%%%%%%%%%%%%%%%%%%%%%%%%%%%%%%%%%%%%%%%%%%%%%%%%%%%%%%%%%%%%%%%%%%%
% Por: Abrantes Araújo Silva Filho
%      abrantesasf@gmail.com
% URL: https://github.com/abrantesasf/latex


%%%%%%%%%%%%%%%%%%%%%%%%%%%%%%%%%%%%%%%%%%%%%%%%%%%%%%%%%%%%%%%%%%%%%%%%%%%%%%%%
%%% Carrega pacotes iniciais necessários para estrutura de controle e para a
%%% criação e o parse de novos comandos
\usepackage{ifthen}
\usepackage{xparse}
\usepackage{ifxetex}


%%%%%%%%%%%%%%%%%%%%%%%%%%%%%%%%%%%%%%%%%%%%%%%%%%%%%%%%%%%%%%%%%%%%%%%%%%%%%%%%
%%% Compilação condicional em PDF:
%%%%%%%%%%%%%%%%%%%%%%%%%%%%%%%%%%%%%%%%%%%%%%%%%%%%%%%%%%%%%%%%%%%%%%%%%%%%%%%%
% Por: Abrantes Araújo Silva Filho
%      abrantesasf@gmail.com
% URL: https://github.com/abrantesasf/latex


%%%%%%%%%%%%%%%%%%%%%%%%%%%%%%%%%%%%%%%%%%%%%%%%%%%%%%%%%%%%%%%%%%%%%%%%%%%%%%%%
%%% Pacotes para compilaçãocondicional em PDF:
\ifxetex
\else
   \usepackage{ifpdf}
\fi


%%%%%%%%%%%%%%%%%%%%%%%%%%%%%%%%%%%%%%%%%%%%%%%%%%%%%%%%%%%%%%%%%%%%%%%%%%%%%%%%
%%% Configurações de layout da página:
%%%%%%%%%%%%%%%%%%%%%%%%%%%%%%%%%%%%%%%%%%%%%%%%%%%%%%%%%%%%%%%%%%%%%%%%%%%%%%%%
% layout.tex
%
% Arquivo de configuração de packages para o LaTeX, considerando diversas
% classes que constumo utilizar.
%
% Para maiores informações, visite:
%    https://github.com/abrantesasf/latex
%    https://github.com/uvv-computacao/uvvtex2
%
% NÃO ALTERE SE NÃO SOUBER O QUE ESTÁ FAZENDO!
%%%%%%%%%%%%%%%%%%%%%%%%%%%%%%%%%%%%%%%%%%%%%%%%%%%%%%%%%%%%%%%%%%%%%%%%%%%%%%%%


%%%%%%%%%%%%%%%%%%%%%%%%%%%%%%%%%%%%%%%%%%%%%%%%%%%%%%%%%%%%%%%%%%%%%%%%%%%%%%%%
%%% Configuração do tamanho da página, margens, espaçamento entrelinhas e, se
%%% necessário, ativa a indentação dos primeiros parágrafos.

% Para article:
\makeatletter
\@ifclassloaded{article}{
  \usepackage{setspace}
  \ifxetex
    \usepackage[xetex]{geometry}
  \else
    \ifpdf
      \usepackage[pdftex]{geometry}
    \else
      \usepackage[dvips]{geometry}
    \fi
  \fi
  \usepackage{microtype}
  \usepackage{lastpage}
}{}
\makeatother

% Para book:
\makeatletter
\@ifclassloaded{book}{
  \usepackage{setspace}
  \ifxetex
    \usepackage[xetex]{geometry}
  \else
    \ifpdf
      \usepackage[pdftex]{geometry}
    \else
      \usepackage[dvips]{geometry}
    \fi
  \fi
  \usepackage{microtype}
  \usepackage{lastpage}
}{}
\makeatother

% Para exam:
\makeatletter
\@ifclassloaded{exam}{
  \usepackage{setspace}
  \ifxetex
    \usepackage[xetex]{geometry}
  \else
    \ifpdf
      \usepackage[pdftex]{geometry}
    \else
      \usepackage[dvips]{geometry}
    \fi
  \fi
  \usepackage{microtype}
  \usepackage{lastpage}
}{}
\makeatother

% Posiciona figuras e tabelas no topo da página quando adicionadas sozinhas
% em um página em branco. Ver https://github.com/abntex/abntex2/issues/170
\makeatletter
\@ifclassloaded{uvvtex2}{
  % Distância entre o topo da página e o primeiro float:
  \setlength{\@fptop}{5pt}
  % O tamanho do parágrafo é dado por:
  \setlength{\parindent}{1.3cm}
  % Controle do espaçamento entre um parágrafo e outro:
  \setlength{\parskip}{0.2cm}  % tente também \onelineskip
  % Para melhorias de justificação:
  \usepackage{microtype}
  % Para contar págians:
  \usepackage{lastpage}
}{}
\makeatother

% Posiciona figuras e tabelas no topo da página quando adicionadas sozinhas
% em um página em branco. Ver https://github.com/abntex/abntex2/issues/170
\makeatletter
\@ifclassloaded{faesatex2}{
  % Distância entre o topo da página e o primeiro float:
  \setlength{\@fptop}{5pt}
  % O tamanho do parágrafo é dado por:
  \setlength{\parindent}{1.3cm}
  % Controle do espaçamento entre um parágrafo e outro:
  \setlength{\parskip}{0.2cm}  % tente também \onelineskip
  % Para melhorias de justificação:
  \usepackage{microtype}
  % Para contar págians:
  \usepackage{lastpage}
}{}
\makeatother


%%%%%%%%%%%%%%%%%%%%%%%%%%%%%%%%%%%%%%%%%%%%%%%%%%%%%%%%%%%%%%%%%%%%%%%%%%%%%%%%
%%% Configurações para autores:
%%%%%%%%%%%%%%%%%%%%%%%%%%%%%%%%%%%%%%%%%%%%%%%%%%%%%%%%%%%%%%%%%%%%%%%%%%%%%%%%
% Por: Abrantes Araújo Silva Filho
%      abrantesasf@gmail.com
% URL: https://github.com/abrantesasf/latex


%%%%%%%%%%%%%%%%%%%%%%%%%%%%%%%%%%%%%%%%%%%%%%%%%%%%%%%%%%%%%%%%%%%%%%%%%%%%%%%%
%%% Ajustes para os autores e afiliações, em artigos
\makeatletter
\@ifclassloaded{article}{
\usepackage{authblk}
}{}
\makeatother


%%%%%%%%%%%%%%%%%%%%%%%%%%%%%%%%%%%%%%%%%%%%%%%%%%%%%%%%%%%%%%%%%%%%%%%%%%%%%%%%
%%% Cabeçalho e rodapé:
%%% Atenção! É MUITO DIFÍCIL ajustar apropriadamente os running headers
%%% e running footers da primeira vez. Você pode confiar no LaTeX e deixar que
%%% ele faça o ajuste automaticamente ou pode quebrar a cabeça e
%%% tentar melhorar algo que o LaTeX já faz muito bem, mesmo que não fique
%%% exatamente do jeito que você quer. O que você prefere? Se quiser ajustar
%%% manualmente, descomente a linha abaixo e faça as configurações no arquivo
%%% "utils/headings.tex".
%%%%%%%%%%%%%%%%%%%%%%%%%%%%%%%%%%%%%%%%%%%%%%%%%%%%%%%%%%%%%%%%%%%%%%%%%%%%%%%%%
% Por: Abrantes Araújo Silva Filho
%      abrantesasf@gmail.com
% URL: https://github.com/abrantesasf/latex


%%%%%%%%%%%%%%%%%%%%%%%%%%%%%%%%%%%%%%%%%%%%%%%%%%%%%%%%%%%%%%%%%%%%%%%%%%%%%%%%
%%% Configurações de cabeçalho e rodapé:
\usepackage{fancyhdr}
\setlength{\headheight}{1cm}
\setlength{\footskip}{1.5cm}
\renewcommand{\headrulewidth}{0.3pt}
\renewcommand{\footrulewidth}{0.0pt}
\pagestyle{fancy}
\renewcommand{\sectionmark}[1]{%
  \markboth{\uppercase{#1}}{}}
\renewcommand{\subsectionmark}[1]{%
  \markright{\uppercase{\thesubsection \hspace{0.1cm} #1}}{}}
\fancyhead{}
\fancyfoot{}
\newcommand{\diminuifonte}{%
    \fontsize{9pt}{9}\selectfont
}
\newcommand{\aumentafonte}{%
    \fontsize{12}{12}\selectfont
}
% Configura cabeçalho e rodapé para documentos TWOSIDE
\fancyhead[EL]{\textbf{\thepage}}
\fancyhead[EC]{}
\fancyhead[ER]{\diminuifonte \textbf{\leftmark}}
\fancyhead[OR]{\textbf{\thepage}}
\fancyhead[OC]{}
\fancyhead[OL]{\diminuifonte \textbf{\rightmark}}
\fancyfoot[EL,EC,ER,OR,OC,OL]{}
% Configura cabeçalho e rodapé para documentos ONESIDE
%%\lhead{ \fancyplain{}{sup esquerdo} }
%%\chead{ \fancyplain{}{sup centro} }
%%\rhead{ \fancyplain{}{\thesection} }
%%\lfoot{ \fancyplain{}{inf esquerdo} }
%%\cfoot{ \fancyplain{}{inf centro} }
%%\rfoot{ \fancyplain{}{\thepage} }


%%%%%%%%%%%%%%%%%%%%%%%%%%%%%%%%%%%%%%%%%%%%%%%%%%%%%%%%%%%%%%%%%%%%%%%%%%%%%%%%
%%% Configuração para as fontes do documento:
%%% Se você quiser usar fontes próprias ou do sistema, precisará ajustar as
%%% configurações no arquivo "utils/fontes.tex".
%%% ATENÇÃO: por padrão eu utilizo XeLaTeX com fontes proprietárias projetadas
%%% por Matthew Butterick, adquiridas em https://mbtype.com, que não podem ser
%%% distribuídas devido à licença de utilização. Se você usar XeLaTeX, você
%%% DEVERÁ OBRIGATORIAMENTE ajustar as fontes no arquivo "utils/fontes.tex".
%%%%%%%%%%%%%%%%%%%%%%%%%%%%%%%%%%%%%%%%%%%%%%%%%%%%%%%%%%%%%%%%%%%%%%%%%%%%%%%%
% fontes.tex
%
% Arquivo de configuração de packages para uso o LaTeX, conforme minhas
% preferências e modelos pessoais.
%
% Para maiores informações, visite:
%    https://github.com/abrantesasf/latex
%
% NÃO ALTERE SE NÃO SOUBER O QUE ESTÁ FAZENDO!
%%%%%%%%%%%%%%%%%%%%%%%%%%%%%%%%%%%%%%%%%%%%%%%%%%%%%%%%%%%%%%%%%%%%%%%%%%%%%%%%


%%%%%%%%%%%%%%%%%%%%%%%%%%%%%%%%%%%%%%%%%%%%%%%%%%%%%%%%%%%%%%%%%%%%%%%%%%%%%%%%
%%% Configurações de encodings e fontes:
\ifxetex
  % Se usar XeLaTeX, usa fontes específicas:
  %\usepackage[tuenc,no-math]{fontspec}
  \usepackage[tuenc]{fontspec}
  \setmainfont{equity-text-a-regular.otf}[
    Path           = /home/abrantesasf/.local/share/fonts/ ,
    BoldFont       = equity-text-a-bold.otf                ,
    ItalicFont     = equity-text-a-italic.otf              ,
    BoldItalicFont = equity-text-a-bold-italic.otf         ,
    SmallCapsFont  = equity-caps-a-regular.otf]
  \setsansfont{equity-text-b-regular.otf}[
    Path           = /home/abrantesasf/.local/share/fonts/ ,
    BoldFont       = equity-text-b-bold.otf                ,
    ItalicFont     = equity-text-b-italic.otf              ,
    BoldItalicFont = equity-text-b-bold-italic.otf]      
  \setmonofont{courier_prime.ttf}[
    Path           = /home/abrantesasf/.local/share/fonts/ ,
    BoldFont       = courier_prime_bold.ttf                ,
    ItalicFont     = courier_prime_italic.ttf              ,
    BoldItalicFont = courier_prime_bold_italic.ttf]
%  \setmonofont{triplicate-t4-regular.otf}[
%    Path           = /home/abrantesasf/.local/share/fonts/ ,
%    BoldFont       = triplicate-t4-bold.otf                ,
%    ItalicFont     = triplicate-t4-italic.otf              ,
%    BoldItalicFont = triplicate-t4-bold-italic.otf]
  \usepackage{fontawesome5}
\else
  % Se não for XeLaTeX, vai com as normais mesmo:
  \usepackage[T1]{fontenc}
  \usepackage[utf8]{inputenc}
  \usepackage{lmodern}
  % Altera a fonte padrão do documento (nem todas funcionam em modo math):
  %   phv = Helvetica
  %   ptm = Times
  %   ppl = Palatino
  %   pbk = bookman
  %   pag = AdobeAvantGarde
  %   pnc = Adobe NewCenturySchoolBook
  %\renewcommand{\familydefault}{ppl}
\fi


%%%%%%%%%%%%%%%%%%%%%%%%%%%%%%%%%%%%%%%%%%%%%%%%%%%%%%%%%%%%%%%%%%%%%%%%%%%%%%%%
%%% Linguagem e hifenização:
%%%%%%%%%%%%%%%%%%%%%%%%%%%%%%%%%%%%%%%%%%%%%%%%%%%%%%%%%%%%%%%%%%%%%%%%%%%%%%%%
% Por: Abrantes Araújo Silva Filho
%      abrantesasf@gmail.com
% URL: https://github.com/abrantesasf/latex


%%%%%%%%%%%%%%%%%%%%%%%%%%%%%%%%%%%%%%%%%%%%%%%%%%%%%%%%%%%%%%%%%%%%%%%%%%%%%%%%
%%% Configurações de linguagem e hifenização
\usepackage{polyglossia}
\setdefaultlanguage[variant=brazilian]{portuguese}
\setotherlanguage{english}

%%%%%%%%%%%%%%%%%%%%%%%%%%%%%%%%%%%%%%%%%%%%%%%%%%%%%%%%%%%%%%%%%%%%%%%%%%%%%%%%
%%% Hifenização específica de alguams palavras quando a hifenização padrão não
%%% estiver hifenizando corretamente:
\hyphenation{%
Git-Hub
LaTeX
po-ly-glos-sia
trans-gre-din-do
}


%%%%%%%%%%%%%%%%%%%%%%%%%%%%%%%%%%%%%%%%%%%%%%%%%%%%%%%%%%%%%%%%%%%%%%%%%%%%%%%%
%%% Sumário:
%%%%%%%%%%%%%%%%%%%%%%%%%%%%%%%%%%%%%%%%%%%%%%%%%%%%%%%%%%%%%%%%%%%%%%%%%%%%%%%%
% Por: Abrantes Araújo Silva Filho
%      abrantesasf@gmail.com
% URL: https://github.com/abrantesasf/latex


%%%%%%%%%%%%%%%%%%%%%%%%%%%%%%%%%%%%%%%%%%%%%%%%%%%%%%%%%%%%%%%%%%%%%%%%%%%%%%%%
%%% Ajustes de sumário:
\makeatletter
\@ifclassloaded{article}{
   \usepackage[nottoc]{tocbibind}
}{}
\@ifclassloaded{book}{
   \usepackage{tocbibind}
}{}
\makeatother


%%%%%%%%%%%%%%%%%%%%%%%%%%%%%%%%%%%%%%%%%%%%%%%%%%%%%%%%%%%%%%%%%%%%%%%%%%%%%%%%
%%% Matemática:
%%%%%%%%%%%%%%%%%%%%%%%%%%%%%%%%%%%%%%%%%%%%%%%%%%%%%%%%%%%%%%%%%%%%%%%%%%%%%%%%
% matematica.tex
%
% Arquivo de configuração de packages para uso o LaTeX, conforme minhas
% preferências e modelos pessoais.
%
% Para maiores informações, visite:
%    https://github.com/abrantesasf/latex
%
% NÃO ALTERE SE NÃO SOUBER O QUE ESTÁ FAZENDO!
%%%%%%%%%%%%%%%%%%%%%%%%%%%%%%%%%%%%%%%%%%%%%%%%%%%%%%%%%%%%%%%%%%%%%%%%%%%%%%%%


%%%%%%%%%%%%%%%%%%%%%%%%%%%%%%%%%%%%%%%%%%%%%%%%%%%%%%%%%%%%%%%%%%%%%%%%%%%%%%%%
%%% Carrega bibliotecas e símbolos matemáticos, fontes adicionais e configura
%%% algumas outras opções
\usepackage{amsmath}
\usepackage{amssymb}
\usepackage{amsthm}
\usepackage{amsfonts}
\usepackage{mathrsfs}
\usepackage{proof}
\usepackage{siunitx}
  \sisetup{group-separator = {\,}}
  \sisetup{group-digits = {integer}}
  \sisetup{output-decimal-marker = {,}}
  \sisetup{separate-uncertainty}
  \sisetup{multi-part-units = single}
  \sisetup{binary-units = true}
\usepackage{syllogism}
  \setsylpuncpa{}
  \setsylpuncpb{}
  \setsylpuncc{}
  \setsylergosign{}
\usepackage{bm}
\usepackage{cancel}
\usepackage{esvect}
\usepackage{mathtools}
\usepackage{icomma}
\usepackage{nicefrac}
%\usepackage{units}

% Altera separador decimal via comando, se necessário (prefira o siunitx):
%\mathchardef\period=\mathcode`.
%\DeclareMathSymbol{.}{\mathord}{letters}{"3B}


%%%%%%%%%%%%%%%%%%%%%%%%%%%%%%%%%%%%%%%%%%%%%%%%%%%%%%%%%%%%%%%%%%%%%%%%%%%%%%%%
%%% Definições para teoremas, etc.

% Para article:
\makeatletter
\@ifclassloaded{article}{
  \theoremstyle{definition}
  \newtheorem{definicao}{Definição}[section]
  \newtheorem{conjecture}{Conjectura}[section]
  \newtheorem{teorema}{Teorema}[section]
  \newtheorem{lemma}{Lema}[section]
  \newtheorem{corolario}{Corolário}[section]
  \theoremstyle{remark}
  \newtheorem*{nota}{Nota}
  \newtheorem*{observacao}{Observação:}
}{}

% Para book:
\@ifclassloaded{book}{
  \theoremstyle{definition}
  \newtheorem{definicao}{Definição}[chapter]
  \newtheorem{conjecture}{Conjectura}[chapter]
  \newtheorem{teorema}{Teorema}[chapter]
  \newtheorem{lemma}{Lema}[chapter]
  \newtheorem{corolario}{Corolário}[chapter]
  \theoremstyle{remark}
  \newtheorem*{nota}{Nota}
  \newtheorem*{observacao}{Observação:}
}{}
\makeatother



%%%%%%%%%%%%%%%%%%%%%%%%%%%%%%%%%%%%%%%%%%%%%%%%%%%%%%%%%%%%%%%%%%%%%%%%%%%%%%%%
%%% Cores:
%%%%%%%%%%%%%%%%%%%%%%%%%%%%%%%%%%%%%%%%%%%%%%%%%%%%%%%%%%%%%%%%%%%%%%%%%%%%%%%%
% cores.tex
%
% Arquivo de configuração de packages para uso o LaTeX, conforme minhas
% preferências e modelos pessoais.
%
% Para maiores informações, visite:
%    https://github.com/abrantesasf/latex
%
% NÃO ALTERE SE NÃO SOUBER O QUE ESTÁ FAZENDO!
%%%%%%%%%%%%%%%%%%%%%%%%%%%%%%%%%%%%%%%%%%%%%%%%%%%%%%%%%%%%%%%%%%%%%%%%%%%%%%%%


%%%%%%%%%%%%%%%%%%%%%%%%%%%%%%%%%%%%%%%%%%%%%%%%%%%%%%%%%%%%%%%%%%%%%%%%%%%%%%%%
%%% Ativa suporte extendido a cores
%%% Opções comuns:
%%%    usenames (16)
%%%    dvipsnames (64)
%%%    svgnames (150)
%%%    x11names (300)
\makeatletter
\@ifclassloaded{beamer}{
   % Beamer sem opções
  \usepackage{xcolor} 
}{
   % Resto opção svgnames
  \usepackage[svgnames, table]{xcolor}
}
\makeatother



%%%%%%%%%%%%%%%%%%%%%%%%%%%%%%%%%%%%%%%%%%%%%%%%%%%%%%%%%%%%%%%%%%%%%%%%%%%%%%%%
%%% Computação:
%%%%%%%%%%%%%%%%%%%%%%%%%%%%%%%%%%%%%%%%%%%%%%%%%%%%%%%%%%%%%%%%%%%%%%%%%%%%%%%%
% Por: Abrantes Araújo Silva Filho
%      abrantesasf@gmail.com
% URL: https://github.com/abrantesasf/latex


%%%%%%%%%%%%%%%%%%%%%%%%%%%%%%%%%%%%%%%%%%%%%%%%%%%%%%%%%%%%%%%%%%%%%%%%%%%%%%%%
%%% Carrega packages relacionados à computação
\usepackage{algorithm2e}
\usepackage{algorithmicx}
\usepackage{algpseudocode}
\usepackage{listings}
  \lstset{literate=
    {á}{{\'a}}1 {é}{{\'e}}1 {í}{{\'i}}1 {ó}{{\'o}}1 {ú}{{\'u}}1
    {Á}{{\'A}}1 {É}{{\'E}}1 {Í}{{\'I}}1 {Ó}{{\'O}}1 {Ú}{{\'U}}1
    {à}{{\`a}}1 {è}{{\`e}}1 {ì}{{\`i}}1 {ò}{{\`o}}1 {ù}{{\`u}}1
    {À}{{\`A}}1 {È}{{\'E}}1 {Ì}{{\`I}}1 {Ò}{{\`O}}1 {Ù}{{\`U}}1
    {ä}{{\"a}}1 {ë}{{\"e}}1 {ï}{{\"i}}1 {ö}{{\"o}}1 {ü}{{\"u}}1
    {Ä}{{\"A}}1 {Ë}{{\"E}}1 {Ï}{{\"I}}1 {Ö}{{\"O}}1 {Ü}{{\"U}}1
    {â}{{\^a}}1 {ê}{{\^e}}1 {î}{{\^i}}1 {ô}{{\^o}}1 {û}{{\^u}}1
    {Â}{{\^A}}1 {Ê}{{\^E}}1 {Î}{{\^I}}1 {Ô}{{\^O}}1 {Û}{{\^U}}1
    {œ}{{\oe}}1 {Œ}{{\OE}}1 {æ}{{\ae}}1 {Æ}{{\AE}}1 {ß}{{\ss}}1
    {ű}{{\H{u}}}1 {Ű}{{\H{U}}}1 {ő}{{\H{o}}}1 {Ő}{{\H{O}}}1
    {ç}{{\c c}}1 {Ç}{{\c C}}1 {ø}{{\o}}1 {å}{{\r a}}1 {Å}{{\r A}}1
    {€}{{\euro}}1 {£}{{\pounds}}1 {«}{{\guillemotleft}}1
    {»}{{\guillemotright}}1 {ñ}{{\~n}}1 {Ñ}{{\~N}}1 {¿}{{?`}}1
  }
\definecolor{mGreen}{rgb}{0,0.6,0}
\definecolor{mGray}{rgb}{0.5,0.5,0.5}
\definecolor{mPurple}{rgb}{0.58,0,0.82}
\definecolor{backgroundColour}{rgb}{0.95,0.95,0.92}
\lstdefinestyle{CStyle}{
    backgroundcolor=\color{backgroundColour},   
    commentstyle=\color{green}\ttfamily,
    keywordstyle=\color{blue}\ttfamily,
    numberstyle=\tiny\color{mGray},
    stringstyle=\color{red}\ttfamily,
    basicstyle=\ttfamily,
    morecomment=[l][\color{magenta}]{\#},
    breakatwhitespace=false,         
    breaklines=true,                 
    captionpos=b,                    
    keepspaces=true,                 
    numbers=left,                    
    numbersep=5pt,                  
    showspaces=false,                
    showstringspaces=false,
    showtabs=false,                  
    tabsize=4,
    language=C
}


%%%%%%%%%%%%%%%%%%%%%%%%%%%%%%%%%%%%%%%%%%%%%%%%%%%%%%%%%%%%%%%%%%%%%%%%%%%%%%%%
%%% Gráficos:
%%%%%%%%%%%%%%%%%%%%%%%%%%%%%%%%%%%%%%%%%%%%%%%%%%%%%%%%%%%%%%%%%%%%%%%%%%%%%%%%
% graficos.tex
%
% Arquivo de configuração de packages para uso o LaTeX, conforme minhas
% preferências e modelos pessoais.
%
% Para maiores informações, visite:
%    https://github.com/abrantesasf/latex
%
% NÃO ALTERE SE NÃO SOUBER O QUE ESTÁ FAZENDO!
%%%%%%%%%%%%%%%%%%%%%%%%%%%%%%%%%%%%%%%%%%%%%%%%%%%%%%%%%%%%%%%%%%%%%%%%%%%%%%%%


%%%%%%%%%%%%%%%%%%%%%%%%%%%%%%%%%%%%%%%%%%%%%%%%%%%%%%%%%%%%%%%%%%%%%%%%%%%%%%%%
%%% Suporte à importação de gráficos externos
\makeatletter
\@ifclassloaded{beamer}{
  \usepackage{graphicx}
}{
  \ifxetex
    \usepackage{graphicx}
  \else
    \ifpdf
      \usepackage[pdftex]{graphicx}
    \else
      \usepackage[dvips]{graphicx}
    \fi
  \fi
}
\makeatother


%%%%%%%%%%%%%%%%%%%%%%%%%%%%%%%%%%%%%%%%%%%%%%%%%%%%%%%%%%%%%%%%%%%%%%%%%%%%%%%%
%%% Suporte à criação de gráficos proceduralmente no LaTeX:
\usepackage{tikz}
\usetikzlibrary{arrows, automata, backgrounds, matrix, patterns, positioning,
                shapes, shadows, calc, circuits.logic.US}



%%%%%%%%%%%%%%%%%%%%%%%%%%%%%%%%%%%%%%%%%%%%%%%%%%%%%%%%%%%%%%%%%%%%%%%%%%%%%%%%
%%% Ícones:
%%%%%%%%%%%%%%%%%%%%%%%%%%%%%%%%%%%%%%%%%%%%%%%%%%%%%%%%%%%%%%%%%%%%%%%%%%%%%%%%
% Por: Abrantes Araújo Silva Filho
%      abrantesasf@gmail.com
% URL: https://github.com/abrantesasf/latex


%%%%%%%%%%%%%%%%%%%%%%%%%%%%%%%%%%%%%%%%%%%%%%%%%%%%%%%%%%%%%%%%%%%%%%%%%%%%%%%%
%%% Pacotes com ícones, figuras e gráficos extras

% Ícones da Creative Commons:
\usepackage{ccicons}


%%%%%%%%%%%%%%%%%%%%%%%%%%%%%%%%%%%%%%%%%%%%%%%%%%%%%%%%%%%%%%%%%%%%%%%%%%%%%%%%
%%% Blocos:
%%%%%%%%%%%%%%%%%%%%%%%%%%%%%%%%%%%%%%%%%%%%%%%%%%%%%%%%%%%%%%%%%%%%%%%%%%%%%%%%
% Por: Abrantes Araújo Silva Filho
%      abrantesasf@gmail.com
% URL: https://github.com/abrantesasf/latex


%%%%%%%%%%%%%%%%%%%%%%%%%%%%%%%%%%%%%%%%%%%%%%%%%%%%%%%%%%%%%%%%%%%%%%%%%%%%%%%%
%%% Ambiente para boxes em posições arbitrárias no Beamer
\makeatletter
\@ifclassloaded{beamer}{
   \usepackage{textpos}
}{}
\makeatother


%%%%%%%%%%%%%%%%%%%%%%%%%%%%%%%%%%%%%%%%%%%%%%%%%%%%%%%%%%%%%%%%%%%%%%%%%%%%%%%%
%%% Boxes coloridos:
%%%%%%%%%%%%%%%%%%%%%%%%%%%%%%%%%%%%%%%%%%%%%%%%%%%%%%%%%%%%%%%%%%%%%%%%%%%%%%%%
% colorbox.tex
%
% Arquivo de configuração de packages para uso o LaTeX, conforme minhas
% preferências e modelos pessoais.
%
% Para maiores informações, visite:
%    https://github.com/abrantesasf/latex
%
% NÃO ALTERE SE NÃO SOUBER O QUE ESTÁ FAZENDO!
%%%%%%%%%%%%%%%%%%%%%%%%%%%%%%%%%%%%%%%%%%%%%%%%%%%%%%%%%%%%%%%%%%%%%%%%%%%%%%%%


%%%%%%%%%%%%%%%%%%%%%%%%%%%%%%%%%%%%%%%%%%%%%%%%%%%%%%%%%%%%%%%%%%%%%%%%%%%%%%%%
%%% Ativa suporte para boxes coloridos
\usepackage[many]{tcolorbox}


%%%%%%%%%%%%%%%%%%%%%%%%%%%%%%%%%%%%%%%%%%%%%%%%%%%%%%%%%%%%%%%%%%%%%%%%%%%%%%%%
%%% Tabelas:
%%%%%%%%%%%%%%%%%%%%%%%%%%%%%%%%%%%%%%%%%%%%%%%%%%%%%%%%%%%%%%%%%%%%%%%%%%%%%%%%
% tabelas.tex
%
% Arquivo de configuração de packages para o LaTeX, considerando diversas
% classes que constumo utilizar.
%
% Para maiores informações, visite:
%    https://github.com/abrantesasf/latex
%    https://github.com/uvv-computacao/uvvtex2
%
% NÃO ALTERE SE NÃO SOUBER O QUE ESTÁ FAZENDO!
%%%%%%%%%%%%%%%%%%%%%%%%%%%%%%%%%%%%%%%%%%%%%%%%%%%%%%%%%%%%%%%%%%%%%%%%%%%%%%%%


%%%%%%%%%%%%%%%%%%%%%%%%%%%%%%%%%%%%%%%%%%%%%%%%%%%%%%%%%%%%%%%%%%%%%%%%%%%%%%%%
%%% Packages para tabelas
\usepackage{array}
\usepackage{tabularray}
\DefTblrTemplate{contfoot-text}{default}{\footnotesize (continua na próxima página)}
\DefTblrTemplate{conthead-text}{default}{(continuação)}
\ExplSyntaxOn
\RenewDocumentCommand \TblrNote { m }
  {
    \cs_if_exist:NT \hypersetup
      { \ExpTblrTemplate { note-border }{ default } }
    \TblrOverlap
      {
        \__tblr_hyper_link:nn {#1}
          { \kern0.03em \textsuperscript { \rmfamily \UseTblrFont { note-tag } #1 } }
      }
  }
\ExplSyntaxOff
\DeclareTblrTemplate{note-tag}{normal}{
  \textsuperscript{\rmfamily \UseTblrFont{note-tag}\InsertTblrNoteTag}
}
\SetTblrTemplate{note-tag}{normal}
\usepackage{longtable}
\usepackage{tabularx}
%\usepackage{tabu}
\usepackage{lscape}
\usepackage{colortbl}  
\usepackage{booktabs}
\usepackage{multirow}
\usepackage{makecell}



%%%%%%%%%%%%%%%%%%%%%%%%%%%%%%%%%%%%%%%%%%%%%%%%%%%%%%%%%%%%%%%%%%%%%%%%%%%%%%%%
%%% Ambientes de listas:
%%%%%%%%%%%%%%%%%%%%%%%%%%%%%%%%%%%%%%%%%%%%%%%%%%%%%%%%%%%%%%%%%%%%%%%%%%%%%%%%
% Por: Abrantes Araújo Silva Filho
%      abrantesasf@gmail.com
% URL: https://github.com/abrantesasf/latex


%%%%%%%%%%%%%%%%%%%%%%%%%%%%%%%%%%%%%%%%%%%%%%%%%%%%%%%%%%%%%%%%%%%%%%%%%%%%%%%%
%%% Packages para ambientes de listas
\makeatletter
\@ifclassloaded{beamer}{
  \usepackage{listings}
}{
  \usepackage{enumitem}
  \usepackage[ampersand]{easylist}
  \usepackage{fancyvrb, bera}
}
\makeatother



%%%%%%%%%%%%%%%%%%%%%%%%%%%%%%%%%%%%%%%%%%%%%%%%%%%%%%%%%%%%%%%%%%%%%%%%%%%%%%%%
%%% Ambientes para quadros:
%%%%%%%%%%%%%%%%%%%%%%%%%%%%%%%%%%%%%%%%%%%%%%%%%%%%%%%%%%%%%%%%%%%%%%%%%%%%%%%%%
% quadros.tex
%
% Arquivo de configuração de packages para uso com a classe uvvTeX2, para a
% formatação de trabalhos acadêmicos na Universidade Vila Velha (UVV)
% (https://www.uvv.br).
%
% Para maiores informações, visite:
%    https://github.com/uvv-computacao/uvvtex2
%
% NÃO ALTERE NADA AQUI!
%%%%%%%%%%%%%%%%%%%%%%%%%%%%%%%%%%%%%%%%%%%%%%%%%%%%%%%%%%%%%%%%%%%%%%%%%%%%%%%%

% Possibilita criação de quadros
% Ver https://github.com/abntex/abntex2/issues/176

\makeatletter
\@ifclassloaded{uvvtex2}{
  \newcommand{\quadroname}{Quadro}
  \newcommand{\listofquadrosname}{Lista de quadros}

  \newfloat[chapter]{quadro}{loq}{\quadroname}
  \newlistof{listofquadros}{loq}{\listofquadrosname}
  \newlistentry{quadro}{loq}{0}

  \setfloatadjustment{quadro}{\centering}
  \counterwithout{quadro}{chapter}
  \renewcommand{\cftquadroname}{\quadroname\space} 
  \renewcommand*{\cftquadroaftersnum}{\hfill--\hfill}

  \setfloatlocations{quadro}{hbtp}
}{}
\makeatother


%%%%%%%%%%%%%%%%%%%%%%%%%%%%%%%%%%%%%%%%%%%%%%%%%%%%%%%%%%%%%%%%%%%%%%%%%%%%%%%%
%%% Ambientes floats e similares:
%%%%%%%%%%%%%%%%%%%%%%%%%%%%%%%%%%%%%%%%%%%%%%%%%%%%%%%%%%%%%%%%%%%%%%%%%%%%%%%%
% Por: Abrantes Araújo Silva Filho
%      abrantesasf@gmail.com
% URL: https://github.com/abrantesasf/latex


%%%%%%%%%%%%%%%%%%%%%%%%%%%%%%%%%%%%%%%%%%%%%%%%%%%%%%%%%%%%%%%%%%%%%%%%%%%%%%%%
%%% Packages para suporte a ambientes floats, captions, etc.:
\usepackage{float}
\usepackage{wrapfig}
\usepackage{placeins}
\usepackage[justification=centering]{caption}
\usepackage{sidecap}
\usepackage{subcaption}


%%%%%%%%%%%%%%%%%%%%%%%%%%%%%%%%%%%%%%%%%%%%%%%%%%%%%%%%%%%%%%%%%%%%%%%%%%%%%%%%
%%% Notas de rodapé:
%%%%%%%%%%%%%%%%%%%%%%%%%%%%%%%%%%%%%%%%%%%%%%%%%%%%%%%%%%%%%%%%%%%%%%%%%%%%%%%%
% Por: Abrantes Araújo Silva Filho
%      abrantesasf@gmail.com
% URL: https://github.com/abrantesasf/latex


%%%%%%%%%%%%%%%%%%%%%%%%%%%%%%%%%%%%%%%%%%%%%%%%%%%%%%%%%%%%%%%%%%%%%%%%%%%%%%%%%
%%% Carrega pacotes para maior controle de notas de rodapé, inclusive dando
%%% um label para uma nota e fazendo referência cruzada.
%%% NÃO ALTERE a ordem de carregamento dos packages!
\usepackage{footmisc}


%%%%%%%%%%%%%%%%%%%%%%%%%%%%%%%%%%%%%%%%%%%%%%%%%%%%%%%%%%%%%%%%%%%%%%%%%%%%%%%%
%%% Referências cruzadas, links, citações:
%%%%%%%%%%%%%%%%%%%%%%%%%%%%%%%%%%%%%%%%%%%%%%%%%%%%%%%%%%%%%%%%%%%%%%%%%%%%%%%%
% crossrefs.tex
%
% Arquivo de configuração de packages para o LaTeX, considerando diversas
% classes que constumo utilizar.
%
% Para maiores informações, visite:
%    https://github.com/abrantesasf/latex
%    https://github.com/uvv-computacao/uvvtex2
%
% NÃO ALTERE SE NÃO SOUBER O QUE ESTÁ FAZENDO!
%%%%%%%%%%%%%%%%%%%%%%%%%%%%%%%%%%%%%%%%%%%%%%%%%%%%%%%%%%%%%%%%%%%%%%%%%%%%%%%%


%%%%%%%%%%%%%%%%%%%%%%%%%%%%%%%%%%%%%%%%%%%%%%%%%%%%%%%%%%%%%%%%%%%%%%%%%%%%%%%%
%%% Carrega pacotes para referências cruzadas, citações dentro do documento,
%%% links para internet e outros.Configura algumas opções.
%%% NÃO ALTERE a ordem de carregamento dos packages!

% Para article:
\makeatletter
\@ifclassloaded{article}{
  \usepackage{varioref}
  \ifxetex
    \usepackage{hyperref}
    \usepackage[brazilian,nameinlink,noabbrev]{cleveref}
  \else
    \ifpdf
      \usepackage[pdftex]{hyperref}
    \else
      \usepackage[dvips]{hyperref}
    \fi
      \usepackage{cleveref}
  \fi
  \usepackage{xurl}
}{}
\makeatother

% Para book:
\makeatletter
\@ifclassloaded{book}{
  \usepackage{varioref}
  \ifxetex
    \usepackage{hyperref}
    \usepackage[brazilian,nameinlink,noabbrev]{cleveref}
  \else
    \ifpdf
      \usepackage[pdftex]{hyperref}
    \else
      \usepackage[dvips]{hyperref}
    \fi
      \usepackage{cleveref}
  \fi
  \usepackage{xurl}
}{}
\makeatother

% Ajustes para citações com o cleveref:
\newcommand{\SeCref}[1]{%
  ``\hyperref[#1]{\Cref*{#1}: \nameref*{#1}}''%
}


%%%%%%%%%%%%%%%%%%%%%%%%%%%%%%%%%%%%%%%%%%%%%%%%%%%%%%%%%%%%%%%%%%%%%%%%%%%%%%%%
%%% Configurações para os hiperlinks em PDF:
%%%%%%%%%%%%%%%%%%%%%%%%%%%%%%%%%%%%%%%%%%%%%%%%%%%%%%%%%%%%%%%%%%%%%%%%%%%%%%%%
% pdfconfig.tex
%
% Arquivo de configuração de packages para uso o LaTeX, conforme minhas
% preferências e modelos pessoais.
%
% Para maiores informações, visite:
%    https://github.com/abrantesasf/latex
%
% NÃO ALTERE SE NÃO SOUBER O QUE ESTÁ FAZENDO!
%%%%%%%%%%%%%%%%%%%%%%%%%%%%%%%%%%%%%%%%%%%%%%%%%%%%%%%%%%%%%%%%%%%%%%%%%%%%%%%%

% Configurações básicas do PDF
\hypersetup{
   unicode=true,
   pdflang={pt-BR},
   bookmarksopen=false,
   bookmarksnumbered=true,
   bookmarksopenlevel=5,
   pdfdisplaydoctitle=true,
   pdfpagemode=UseOutlines,
   pdfstartview=FitH,
   pdfnewwindow=true,
   %pagebackref=true,
   pdfcreator={TeX Live 2021: XeTeX (0.999993) + LaTeX2e (2021-06-01) + eTeX (2.6) + TeX (3.141592653)},
   pdfproducer={xdvipdfmx (20210318)},
   %pdftitle={\@title}, 
   %pdfauthor={\@author},
   %pdfsubject={\imprimirpreambulo},
   %pdfkeywords={\imprimirpalavraschave, \imprimirkeywords}, 
   bookmarksdepth=5,
   citecolor=cyan,      % cor dos links para as referências bibliográficas
   linkcolor=red,       % cor dos links internos (sumário, figuras, rodapé)
   filecolor=magenta,   % dor dos links para arquivos
   urlcolor=blue,       % cor para links URL
   anchorcolor=black    % cor para o texto âncora (alvo)
}



%%%%%%%%%%%%%%%%%%%%%%%%%%%%%%%%%%%%%%%%%%%%%%%%%%%%%%%%%%%%%%%%%%%%%%%%%%%%%%%%
%%% Glossário e índice remissivo:
%%%%%%%%%%%%%%%%%%%%%%%%%%%%%%%%%%%%%%%%%%%%%%%%%%%%%%%%%%%%%%%%%%%%%%%%%%%%%%%%
% Por: Abrantes Araújo Silva Filho
%      abrantesasf@gmail.com
% URL: https://github.com/abrantesasf/latex


%%%%%%%%%%%%%%%%%%%%%%%%%%%%%%%%%%%%%%%%%%%%%%%%%%%%%%%%%%%%%%%%%%%%%%%%%%%%%%%%
%%% Pacote para índice remissivo
\makeatletter
\@ifclassloaded{beamer}{}
{
   \usepackage{makeidx}
   %\makeindex
}
\makeatother


%%%%%%%%%%%%%%%%%%%%%%%%%%%%%%%%%%%%%%%%%%%%%%%%%%%%%%%%%%%%%%%%%%%%%%%%%%%%%%%%
%%% Pacote para glossário
\makeatletter
\@ifclassloaded{beamer}{}
{
   \usepackage[toc]{glossaries}
   %\newglossary[nlg]{notation}{not}{ntn}{Notação}
   %\newglossaryentry{not:set}{
   %   type = notation,
   %   name = {$\mathbb{N}$},
   %   description = conjunto dos números naturais,
   %   sort = {N}}
   %\makeglossaries
}
\makeatother


%%%%%%%%%%%%%%%%%%%%%%%%%%%%%%%%%%%%%%%%%%%%%%%%%%%%%%%%%%%%%%%%%%%%%%%%%%%%%%%%
%%% Referências bibliográficas:
%%%%%%%%%%%%%%%%%%%%%%%%%%%%%%%%%%%%%%%%%%%%%%%%%%%%%%%%%%%%%%%%%%%%%%%%%%%%%%%%
% Por: Abrantes Araújo Silva Filho
%      abrantesasf@gmail.com
% URL: https://github.com/abrantesasf/latex


%%%%%%%%%%%%%%%%%%%%%%%%%%%%%%%%%%%%%%%%%%%%%%%%%%%%%%%%%%%%%%%%%%%%%%%%%%%%%%%%
%%% Referências bibliográficas
%\usepackage[square,numbers]{natbib}
%\bibliographystyle{plainnat}

\usepackage[alf]{abntex2cite}
%\bibliographystyle{abnt-alf}


%%%%%%%%%%%%%%%%%%%%%%%%%%%%%%%%%%%%%%%%%%%%%%%%%%%%%%%%%%%%%%%%%%%%%%%%%%%%%%%%
%%% Ajustes para a classe EXAM:
%%%%%%%%%%%%%%%%%%%%%%%%%%%%%%%%%%%%%%%%%%%%%%%%%%%%%%%%%%%%%%%%%%%%%%%%%%%%%%%%
% exam.tex
%
% Arquivo de configuração para a classe EXAM, conforme minhas
% preferências e modelos pessoais.
%
% Para maiores informações, visite:
%    https://github.com/abrantesasf/latex
%
% NÃO ALTERE SE NÃO SOUBER O QUE ESTÁ FAZENDO!
%%%%%%%%%%%%%%%%%%%%%%%%%%%%%%%%%%%%%%%%%%%%%%%%%%%%%%%%%%%%%%%%%%%%%%%%%%%%%%%%


% Utilitários para a classe Exam:
\makeatletter
\@ifclassloaded{exam}{%
   % Linhas para respostas
   \setlength\linefillthickness{0.2pt}
   \newcommand{\umalinha}{\fillwithlines{0.25in}}
   \newcommand{\duaslinhas}{\fillwithlines{0.50in}}
   \newcommand{\treslinhas}{\fillwithlines{0.75in}}
   \newcommand{\quatrolinhas}{\fillwithlines{1.00in}}
   \newcommand{\cincolinhas}{\fillwithlines{1.25in}}
   \newcommand{\seislinhas}{\fillwithlines{1.50in}}
   \newcommand{\setelinhas}{\fillwithlines{1.75in}}
   \newcommand{\oitolinhas}{\fillwithlines{2.00in}}
   \newcommand{\novelinhas}{\fillwithlines{2.25in}}
   \newcommand{\dezlinhas}{\fillwithlines{2.50in}}
   
   % Verdadeiro ou Falso
   \newcommand{\vf}[1][{}]{%
      \fillin[#1][0.25in]%
   }
   
   % Título das respostas:
   %\renewcommand{\solutiontitle}{\noindent\textbf{Resposta:}\par\noindent}
   \renewcommand{\solutiontitle}{\noindent}
   %\shadedsolutions
   %\SolutionEmphasis{\itshape}
   
   % Altera a marca de questão correta e a ênfase:
   \checkedchar{$\Rightarrow$}
   %\CorrectChoiceEmphasis{\color{red}}
}{}
\makeatother



%%%%%%%%%%%%%%%%%%%%%%%%%%%%%%%%%%%%%%%%%%%%%%%%%%%%%%%%%%%%%%%%%%%%%%%%%%%%%%%%
%%% Ajustes para textos:
%%%%%%%%%%%%%%%%%%%%%%%%%%%%%%%%%%%%%%%%%%%%%%%%%%%%%%%%%%%%%%%%%%%%%%%%%%%%%%%%
% texto.tex
%
% Arquivo de configuração de packages para o LaTeX, considerando diversas
% classes que constumo utilizar.
%
% Para maiores informações, visite:
%    https://github.com/abrantesasf/latex
%    https://github.com/uvv-computacao/uvvtex2
%
% NÃO ALTERE SE NÃO SOUBER O QUE ESTÁ FAZENDO!
%%%%%%%%%%%%%%%%%%%%%%%%%%%%%%%%%%%%%%%%%%%%%%%%%%%%%%%%%%%%%%%%%%%%%%%%%%%%%%%%


%%%%%%%%%%%%%%%%%%%%%%%%%%%%%%%%%%%%%%%%%%%%%%%%%%%%%%%%%%%%%%%%%%%%%%%%%%%%%%%%
%%% Pacote para texto dummy:
\usepackage{blindtext}


%%%%%%%%%%%%%%%%%%%%%%%%%%%%%%%%%%%%%%%%%%%%%%%%%%%%%%%%%%%%%%%%%%%%%%%%%%%%%%%%
%%% Ambientes, aliases, comandos e customizações em geral para serem
%%% utilizadas nos documentos. Defina aqui qualquer customização específica
%%% para seus documentos!
%%%%%%%%%%%%%%%%%%%%%%%%%%%%%%%%%%%%%%%%%%%%%%%%%%%%%%%%%%%%%%%%%%%%%%%%%%%%%%%%
% customizacoes.tex
%
% Arquivo de configuração de packages para uso o LaTeX, conforme minhas
% preferências e modelos pessoais.
%
% Para maiores informações, visite:
%    https://github.com/abrantesasf/latex
%
% NÃO ALTERE SE NÃO SOUBER O QUE ESTÁ FAZENDO!
%%%%%%%%%%%%%%%%%%%%%%%%%%%%%%%%%%%%%%%%%%%%%%%%%%%%%%%%%%%%%%%%%%%%%%%%%%%%%%%%


%%%%%%%%%%%%%%%%%%%%%%%%%%%%%%%%%%%%%%%%%%%%%%%%%%%%%%%%%%%%%%%%%%%%%%%%%%%%%%%%
%%% Vários comandinhos úteis e outras gambiarras

% Commando para ``italizar´´ palavras em inglês (e outras línguas!):
\newcommand{\ingles}[1]{\textit{#1}}

% Comando para escrever uma função e símbolos em fonte monoespaçada:
\newcommand{\funcao}[1]{\textbf{\texttt{#1}}}
\newcommand{\simbolo}[1]{\texttt{#1}}

% Commando para colocar o espaço correto entre um número e sua unidade:
\newcommand{\unidade}[2]{\ensuremath{#1\,\mathrm{#2}}}
\newcommand{\unidado}[2]{{#1}\,{#2}}

% Produz ordinal masculino ou feminino dependendo do segundo argumento:
\newcommand{\ordinal}[2]{%
#1%
\ifthenelse{\equal{a}{#2}}%
{\textordfeminine}%
{\textordmasculine}}

% Comando para colocar o autor da citação corretamente justificado à direita:
% Modificado de: https://latex.org/forum/viewtopic.php?p=15605&sid=216895679326b531e85caec779e1d710#p15605
%\newcommand*{\fontequote}[2]{%
%   \unskip\hspace*{1em plus 1fill}%
%   \linebreak\hspace*{\fill}\mbox{#1}\vspace{-0.15cm}
%   \linebreak\hspace*{\fill}\mbox{#2}
%}

% Comandos para os quadros padronizados dos livros de matemática:
\newcommand{\abrevideo}[1]{%
\begin{tcolorbox}[colback=yellow!5!white,colframe=yellow!75!black,
       title=\textbf{Vídeo(s) de referência: #1}]}
\newcommand{\fechavideo}{\end{tcolorbox}}

% Comandos para os ícones comuns:
\newcommand{\youtube}{\faYoutube}

% Comandos para álgebra relacional (testados em XeLaTeX):
\def\ojoin{\setbox0=\hbox{$\Join$}%
\rule[0.02ex]{.4em}{.5pt}\llap{\rule[0.95ex]{.4em}{.5pt}}}
\def\leftouterjoin{\mathbin{\ojoin\mkern-7.8mu\Join}}
\def\rightouterjoin{\mathbin{\Join\mkern-7.8mu\ojoin}}
\def\fullouterjoin{\mathbin{\ojoin\mkern-7.8mu\Join\mkern-7.8mu\ojoin}}

\usepackage{needspace}

% ============================================================
% Estrutura jurídica: artigos, cláusulas, parágrafos,
% incisos e alíneas
% ============================================================

\newcounter{artigo}
\newcounter{clausula}
\newcounter{paragrafo}
\newcounter{inciso}
\newcounter{alinea}

% ============================================================
% Medidas
% ============================================================

% Parágrafos numerados
\newlength{\recuoparagrafo}
\newlength{\larguraparagrafo}

\setlength{\recuoparagrafo}{\parindent}
\setlength{\larguraparagrafo}{2.5em}

% Parágrafo único
\newlength{\larguraparagrafounico}

\settowidth{\larguraparagrafounico}{%
  \textbf{Parágrafo único.}\enspace
}

% Incisos
\newlength{\recuoinciso}
\newlength{\largurainciso}

\setlength{\recuoinciso}{%
  \dimexpr\recuoparagrafo+\larguraparagrafo\relax
}
\setlength{\largurainciso}{3em}

% Alíneas
\newlength{\recuoalinea}
\newlength{\larguraalinea}

\setlength{\recuoalinea}{%
  \dimexpr\recuoinciso+\largurainciso\relax
}
\setlength{\larguraalinea}{2em}


% ------------------------------------------------------------
% Formatação dos números
% ------------------------------------------------------------

\renewcommand{\theartigo}{\arabic{artigo}}
\renewcommand{\theclausula}{\arabic{clausula}}
\renewcommand{\theparagrafo}{\arabic{paragrafo}}
\renewcommand{\theinciso}{\Roman{inciso}}
\renewcommand{\thealinea}{\alph{alinea}}

% ------------------------------------------------------------
% Artigo
%
% Uso:
%   \artigo
%   \artigo[Do objeto]
% ------------------------------------------------------------

\newcommand{\artigo}[1][]{%
  \par
  \Needspace{3\baselineskip}%
  \vspace{1.5em}%
  \refstepcounter{artigo}%
  \setcounter{paragrafo}{0}%
  \setcounter{inciso}{0}%
  \setcounter{alinea}{0}%
  \noindent
  \textbf{%
     Artigo~\theartigo
     \ifnum\value{artigo}<10 \textsuperscript{o}\fi%
     \if\relax\detokenize{#1}\relax
     \else\space-- #1%
     \fi
  }%
  \par\nobreak
}

% ------------------------------------------------------------
% Cláusula
%
% Uso:
%   \clausula
%   \clausula[Do objeto]
% ------------------------------------------------------------

\newcommand{\clausula}[1][]{%
  \par
  \Needspace{3\baselineskip}%
  \vspace{1.5em}%
  \refstepcounter{clausula}%
  \setcounter{paragrafo}{0}%
  \setcounter{inciso}{0}%
  \setcounter{alinea}{0}%
  \Needspace{3\baselineskip}
  \par\vspace{1.5em}
  \noindent
  \textbf{%
    Cláusula~\theclausula
    \ifnum\value{clausula}<10 \textsuperscript{a}\fi%
    \if\relax\detokenize{#1}\relax
    \else\space-- #1%
    \fi
  }%
  \par\nobreak
}


% ------------------------------------------------------------
% Parágrafo numerado
%
% Uso:
%   \paragrafo Texto...
%
% Produz:
%   § 1º Texto...
% ------------------------------------------------------------
\setlength{\recuoparagrafo}{\parindent}
\setlength{\larguraparagrafo}{2.5em}

\newcommand{\paragrafo}[1]{%
  \refstepcounter{paragrafo}%
  \setcounter{inciso}{0}%
  \setcounter{alinea}{0}%
  \par
  \begingroup
    \setlength{\parindent}{0pt}%
    \advance\leftskip by
      \dimexpr\recuoparagrafo+\larguraparagrafo\relax
    \noindent
    \llap{%
      \makebox[\larguraparagrafo][l]{%
        \textbf{%
          §~\theparagrafo
          \ifnum\value{paragrafo}<10 º\fi
        }%
      }%
    }%
    \ignorespaces#1%
    \par
  \endgroup
}

% ------------------------------------------------------------
% Parágrafo único
%
% Uso:
%   \paragrafounico Texto...
% ------------------------------------------------------------

\newcommand{\paragrafounico}[1]{%
  \setcounter{inciso}{0}%
  \setcounter{alinea}{0}%
  \par
  \begingroup
    \setlength{\parindent}{0pt}%
    \advance\leftskip by
      \dimexpr\recuoparagrafo+\larguraparagrafounico\relax
    \noindent
    \llap{%
      \makebox[\larguraparagrafounico][l]{%
        \textbf{§ único.}%
      }%
    }%
    \ignorespaces#1%
    \par
  \endgroup
}

% ------------------------------------------------------------
% Inciso
%
% Uso:
%   \inciso Texto...
%
% Produz:
%   I – Texto...
% ------------------------------------------------------------

\newcommand{\inciso}[1]{%
  \refstepcounter{inciso}%
  \setcounter{alinea}{0}%
  \par
  \begingroup
    \setlength{\parindent}{0pt}%
    \advance\leftskip by
      \dimexpr\recuoinciso+\largurainciso\relax
    \noindent
    \llap{%
      \makebox[\largurainciso][l]{%
        \textbf{\theinciso\space--}%
      }%
    }%
    \ignorespaces#1%
    \par
  \endgroup
}

% ------------------------------------------------------------
% Alínea
%
% Uso:
%   \alinea Texto...
%
% Produz:
%   a) Texto...
% ------------------------------------------------------------

\newcommand{\alinea}[1]{%
  \refstepcounter{alinea}%
  \par
  \begingroup
    \setlength{\parindent}{0pt}%
    \advance\leftskip by
      \dimexpr\recuoalinea+\larguraalinea\relax
    \noindent
    \llap{%
      \makebox[\larguraalinea][l]{%
        \textbf{\thealinea)}%
      }%
    }%
    \ignorespaces#1%
    \par
  \endgroup
}

%%%%%%%%%%%%%%%%%%%%%%%%%%%%%%%%%%%%%%%%%%%%%%%%%%%%%%%%%%%%%%%%%%%%%%%%%%%%%%%%
% engenharia.tex
%
% Configurações específicas do livro "Engenharia em 500 Gramas": pacotes de
% engenharia (circuitos, gráficos de dados, diagramas de tempo), caixas
% didáticas semânticas e comandos do livro.
%
% Deve ser carregado DEPOIS de utils/preambulo.tex (que já carrega tikz,
% tcolorbox[many], listings, siunitx, xcolor[svgnames] e cleveref). Para isso,
% em livro.tex, logo após \input{utils/preambulo.tex}, acrescente:
%    \input{utils/engenharia.tex}
%
% Nada aqui altera arquivos do modelo. Os capítulos devem usar SOMENTE os
% ambientes semânticos abaixo; a aparência é definida neste arquivo.
%%%%%%%%%%%%%%%%%%%%%%%%%%%%%%%%%%%%%%%%%%%%%%%%%%%%%%%%%%%%%%%%%%%%%%%%%%%%%%%%


%%%%%%%%%%%%%%%%%%%%%%%%%%%%%%%%%%%%%%%%%%%%%%%%%%%%%%%%%%%%%%%%%%%%%%%%%%%%%%%%
%%% Circuitos: esquemas elétricos. Estilo europeu (resistor retangular), o
%%% usado no Brasil. A opção siunitx faz os rótulos usarem \SI, herdando a
%%% vírgula decimal configurada em utils/matematica.tex.
\usepackage[siunitx, european]{circuitikz}
\ctikzset{bipoles/length=1.2cm}

%%%%%%%%%%%%%%%%%%%%%%%%%%%%%%%%%%%%%%%%%%%%%%%%%%%%%%%%%%%%%%%%%%%%%%%%%%%%%%%%
%%% Gráficos de dados: curvas dos ensaios (motor, sensores, bateria).
%%% Vírgula decimal e espaço fino como separador de milhar, como no siunitx.
\usepackage{pgfplots}
\usepackage{pgfplotstable}
\pgfplotsset{
  compat=1.18,
  /pgf/number format/.cd, use comma, 1000 sep={\,},
  every axis/.append style={grid=major, grid style={dotted, gray!60}},
}

%%%%%%%%%%%%%%%%%%%%%%%%%%%%%%%%%%%%%%%%%%%%%%%%%%%%%%%%%%%%%%%%%%%%%%%%%%%%%%%%
%%% Bibliotecas TikZ para diagramas de blocos, mecânica e o robô no dohyô.
%%% (utils/graficos.tex já carrega arrows, automata, backgrounds, matrix,
%%% patterns, positioning, shapes, shadows, calc e circuits.logic.US.)
\usetikzlibrary{arrows.meta, decorations.pathmorphing, decorations.markings,
                fit, chains, angles}
% Não carregue a biblioteca "quotes": a sintaxe ["rótulo"] em \draw e \pic
% falha neste modelo (o caractere " é ativo por causa do babel/polyglossia).
% Use "node[above]{$\theta$}" e "pic text=$\alpha$" no lugar. Alternativa:
% \shorthandoff{"} antes do tikzpicture (e a biblioteca quotes), mas isso
% desliga localmente os atalhos "" e "- do português.

%%%%%%%%%%%%%%%%%%%%%%%%%%%%%%%%%%%%%%%%%%%%%%%%%%%%%%%%%%%%%%%%%%%%%%%%%%%%%%%%
%%% Diagramas de tempo (PWM, pulso do HC-SR04, protocolo IR) e formatos de
%%% dados/registradores.
\usepackage{tikz-timing}
\usepackage{bytefield}

%%%%%%%%%%%%%%%%%%%%%%%%%%%%%%%%%%%%%%%%%%%%%%%%%%%%%%%%%%%%%%%%%%%%%%%%%%%%%%%%
%%% QR codes para os vídeos complementares.
\usepackage{qrcode}

%%% Árvores (classificações, hierarquias).
\usepackage[edges]{forest}


%%%%%%%%%%%%%%%%%%%%%%%%%%%%%%%%%%%%%%%%%%%%%%%%%%%%%%%%%%%%%%%%%%%%%%%%%%%%%%%%
%%% Tons de cinza usados pelas caixas (edição em preto e branco). Para uma
%%% edição colorida, basta redefinir estas cores.
\providecolor{cinzaClaro}{gray}{0.90}
\providecolor{cinzaMedio}{gray}{0.55}
\providecolor{cinzaEscuro}{gray}{0.25}


%%%%%%%%%%%%%%%%%%%%%%%%%%%%%%%%%%%%%%%%%%%%%%%%%%%%%%%%%%%%%%%%%%%%%%%%%%%%%%%%
%%% Comandos do livro

% Plataformas: usar sempre os comandos, nunca digitar à mão.
\newcommand{\RI}{R1}
\newcommand{\RII}{R2}

% Marcas de rascunho: visíveis enquanto \rascunhotrue; somem na edição final
% (troque para \rascunhofalse).
\newif\ifrascunho
\rascunhotrue
\newcommand{\TODO}[1]{\ifrascunho{\sffamily\footnotesize\color{cinzaEscuro}%
  [\textbf{A FAZER:} #1]}\fi}

% \codigo[opções listings]{arquivo}{marcador}
% Extrai do arquivo-fonte REAL o trecho delimitado por
%    // @livro: <marcador>
%    ...
%    // @livro: <marcador>_fim
% Assim o livro nunca diverge do firmware: se o código muda, o trecho muda
% junto. Usa o estilo CStyle definido em utils/computacao.tex.
%
% Limitações do listings (testadas no XeLaTeX):
%  1. Caracteres acima de U+00FF no arquivo-fonte (travessão, símbolo de
%     aproximação, emojis) geram linhas espúrias. Acentos do português são
%     seguros.
%  2. Comentário de bloco /* ... */ de várias linhas dentro de um trecho
%     marcado faz o marcador final ser ignorado. Dentro de trechos marcados,
%     use //.
\newcommand{\codigo}[3][]{%
  \lstinputlisting[style=CStyle,
    linerange={#3-#3_fim},
    rangeprefix=//\ @livro:\ , rangesuffix=,
    includerangemarker=false, #1]{#2}}


%%%%%%%%%%%%%%%%%%%%%%%%%%%%%%%%%%%%%%%%%%%%%%%%%%%%%%%%%%%%%%%%%%%%%%%%%%%%%%%%
%%% Caixas didáticas semânticas
% Recuo de parágrafo dentro das caixas: o mesmo do texto corrido. O tcolorbox
% zera \parindent no conteúdo; o primeiro parágrafo fica sem recuo.
\newlength{\recuocaixa}
\AtBeginDocument{\setlength{\recuocaixa}{\parindent}}
\tcbset{
  e500base/.style={
    enhanced, breakable, sharp corners,
    colback=white, colframe=black, coltitle=black,
    boxrule=0.4pt, left=2mm, right=2mm, top=1.2mm, bottom=1.2mm,
    fonttitle=\sffamily\bfseries\small,
    before skip=0.8\baselineskip, after skip=0.8\baselineskip,
    before upper={\setlength{\parindent}{\recuocaixa}\noindent},
  },
  % Título na primeira linha do conteúdo. Substitui "attach title to upper",
  % que redefine before upper e perderia o recuo acima.
  e500tituloInterno/.style={detach title,
    before upper={\setlength{\parindent}{\recuocaixa}%
      \noindent\tcbtitle\par\smallskip\noindent}},
}

% Fundamento: teoria transferível, independente do robô.
\newtcolorbox{fundamento}[1][]{e500base,
  frame hidden, borderline west={2.5pt}{0pt}{black},
  title={Fundamento\ifx&#1&\else\ --- #1\fi},
  e500tituloInterno}

% Modelo: desenvolvimento matemático.
\newtcolorbox{modelo}[1][]{e500base,
  colbacktitle=cinzaClaro, titlerule=0.4pt,
  title={Modelo\ifx&#1&\else\ --- #1\fi}}

% Projeto R1 / Projeto R2: decisão aplicada a cada plataforma.
\newtcolorbox{projetoRI}[1][]{e500base,
  colbacktitle=black, coltitle=white,
  title={Projeto R1\ifx&#1&\else\ --- #1\fi}}
\newtcolorbox{projetoRII}[1][]{e500base,
  colbacktitle=cinzaEscuro, coltitle=white, boxrule=0.8pt,
  title={Projeto R2\ifx&#1&\else\ --- #1\fi}}

% Laboratório: experimento obrigatório do capítulo.
\newtcolorbox{laboratorio}[1][]{e500base,
  frame hidden, borderline north={1.2pt}{0pt}{black},
  borderline south={1.2pt}{0pt}{black},
  title={Laboratório\ifx&#1&\else\ --- #1\fi},
  e500tituloInterno}

% Quando a engenharia dá errado: falha real, documentada.
\newtcolorbox{falhareal}[1][]{e500base,
  frame hidden, borderline west={5pt}{0pt}{black},
  title={Quando a engenharia dá errado\ifx&#1&\else\ --- #1\fi},
  e500tituloInterno}

% Regra -> requisito: regulamento traduzido em requisito verificável.
\newtcolorbox{regra}[1][]{e500base,
  frame hidden,
  borderline={0.6pt}{0pt}{black,dashed},
  title={Regra $\rightarrow$ requisito\ifx&#1&\else\ --- #1\fi},
  e500tituloInterno}

% Competição: impacto da decisão no desempenho.
\newtcolorbox{competicao}[1][]{e500base,
  colframe=cinzaMedio, boxrule=0.8pt,
  title={Competição\ifx&#1&\else\ --- #1\fi},
  e500tituloInterno}

% Avançado: conteúdo opcional; pode ser pulado sem perda de continuidade.
\newtcolorbox{avancado}[1][]{e500base,
  colframe=cinzaMedio, boxrule=0.4pt,
  borderline={0.4pt}{1.2pt}{cinzaMedio},
  title={Avançado (opcional)\ifx&#1&\else\ --- #1\fi},
  e500tituloInterno}

% Atenção: cuidado prático, risco de dano ou erro comum de uso.
\newtcolorbox{atencao}[1][]{e500base,
  boxrule=1.2pt,
  title={Atenção\ifx&#1&\else\ --- #1\fi},
  e500tituloInterno}

% Vídeo complementar: nunca indispensável para entender o texto.
%    \video{V05.1}{Caracterização do motor em bancada}{https://...}
\newcommand{\video}[3]{%
  \begin{tcolorbox}[e500base, colframe=cinzaMedio, boxrule=0.4pt,
                    sidebyside, lefthand width=18mm, sidebyside gap=3mm]
    \qrcode[height=16mm]{#3}
    \tcblower
    {\sffamily\bfseries\small Vídeo complementar #1}\par
    #2\par
    {\footnotesize\ttfamily\url{#3}}
  \end{tcolorbox}}

, acrescente:
%    %%%%%%%%%%%%%%%%%%%%%%%%%%%%%%%%%%%%%%%%%%%%%%%%%%%%%%%%%%%%%%%%%%%%%%%%%%%%%%%%
% engenharia.tex
%
% Configurações específicas do livro "Engenharia em 500 Gramas": pacotes de
% engenharia (circuitos, gráficos de dados, diagramas de tempo), caixas
% didáticas semânticas e comandos do livro.
%
% Deve ser carregado DEPOIS de utils/preambulo.tex (que já carrega tikz,
% tcolorbox[many], listings, siunitx, xcolor[svgnames] e cleveref). Para isso,
% em livro.tex, logo após %%%%%%%%%%%%%%%%%%%%%%%%%%%%%%%%%%%%%%%%%%%%%%%%%%%%%%%%%%%%%%%%%%%%%%%%%%%%%%%%
% preambulo.tex
%
% Arquivo de configuração de packages para o LaTeX, considerando diversas
% classes que constumo utilizar.
%
% Para maiores informações, visite:
%    https://github.com/abrantesasf/latex
%    https://github.com/uvv-computacao/uvvtex2
%
% Para maiores informações, visite:
%    https://github.com/abrantesasf/latex
%
% NÃO ALTERE SE NÃO SOUBER O QUE ESTÁ FAZENDO!
%%%%%%%%%%%%%%%%%%%%%%%%%%%%%%%%%%%%%%%%%%%%%%%%%%%%%%%%%%%%%%%%%%%%%%%%%%%%%%%%


%%%%%%%%%%%%%%%%%%%%%%%%%%%%%%%%%%%%%%%%%%%%%%%%%%%%%%%%%%%%%%%%%%%%%%%%%%%%%%%%
%%% Estruturas de controle:
\input{utils/controles}

%%%%%%%%%%%%%%%%%%%%%%%%%%%%%%%%%%%%%%%%%%%%%%%%%%%%%%%%%%%%%%%%%%%%%%%%%%%%%%%%
%%% Compilação condicional em PDF:
\input{utils/pdf}

%%%%%%%%%%%%%%%%%%%%%%%%%%%%%%%%%%%%%%%%%%%%%%%%%%%%%%%%%%%%%%%%%%%%%%%%%%%%%%%%
%%% Configurações de layout da página:
\input{utils/layout}

%%%%%%%%%%%%%%%%%%%%%%%%%%%%%%%%%%%%%%%%%%%%%%%%%%%%%%%%%%%%%%%%%%%%%%%%%%%%%%%%
%%% Configurações para autores:
\input{utils/autores}

%%%%%%%%%%%%%%%%%%%%%%%%%%%%%%%%%%%%%%%%%%%%%%%%%%%%%%%%%%%%%%%%%%%%%%%%%%%%%%%%
%%% Cabeçalho e rodapé:
%%% Atenção! É MUITO DIFÍCIL ajustar apropriadamente os running headers
%%% e running footers da primeira vez. Você pode confiar no LaTeX e deixar que
%%% ele faça o ajuste automaticamente ou pode quebrar a cabeça e
%%% tentar melhorar algo que o LaTeX já faz muito bem, mesmo que não fique
%%% exatamente do jeito que você quer. O que você prefere? Se quiser ajustar
%%% manualmente, descomente a linha abaixo e faça as configurações no arquivo
%%% "utils/headings.tex".
%\input{utils/headings}

%%%%%%%%%%%%%%%%%%%%%%%%%%%%%%%%%%%%%%%%%%%%%%%%%%%%%%%%%%%%%%%%%%%%%%%%%%%%%%%%
%%% Configuração para as fontes do documento:
%%% Se você quiser usar fontes próprias ou do sistema, precisará ajustar as
%%% configurações no arquivo "utils/fontes.tex".
%%% ATENÇÃO: por padrão eu utilizo XeLaTeX com fontes proprietárias projetadas
%%% por Matthew Butterick, adquiridas em https://mbtype.com, que não podem ser
%%% distribuídas devido à licença de utilização. Se você usar XeLaTeX, você
%%% DEVERÁ OBRIGATORIAMENTE ajustar as fontes no arquivo "utils/fontes.tex".
\input{utils/fontes}

%%%%%%%%%%%%%%%%%%%%%%%%%%%%%%%%%%%%%%%%%%%%%%%%%%%%%%%%%%%%%%%%%%%%%%%%%%%%%%%%
%%% Linguagem e hifenização:
\input{utils/linguagem}

%%%%%%%%%%%%%%%%%%%%%%%%%%%%%%%%%%%%%%%%%%%%%%%%%%%%%%%%%%%%%%%%%%%%%%%%%%%%%%%%
%%% Sumário:
\input{utils/toc}

%%%%%%%%%%%%%%%%%%%%%%%%%%%%%%%%%%%%%%%%%%%%%%%%%%%%%%%%%%%%%%%%%%%%%%%%%%%%%%%%
%%% Matemática:
\input{utils/matematica}

%%%%%%%%%%%%%%%%%%%%%%%%%%%%%%%%%%%%%%%%%%%%%%%%%%%%%%%%%%%%%%%%%%%%%%%%%%%%%%%%
%%% Cores:
\input{utils/cores}

%%%%%%%%%%%%%%%%%%%%%%%%%%%%%%%%%%%%%%%%%%%%%%%%%%%%%%%%%%%%%%%%%%%%%%%%%%%%%%%%
%%% Computação:
\input{utils/computacao}

%%%%%%%%%%%%%%%%%%%%%%%%%%%%%%%%%%%%%%%%%%%%%%%%%%%%%%%%%%%%%%%%%%%%%%%%%%%%%%%%
%%% Gráficos:
\input{utils/graficos}

%%%%%%%%%%%%%%%%%%%%%%%%%%%%%%%%%%%%%%%%%%%%%%%%%%%%%%%%%%%%%%%%%%%%%%%%%%%%%%%%
%%% Ícones:
\input{utils/icones}

%%%%%%%%%%%%%%%%%%%%%%%%%%%%%%%%%%%%%%%%%%%%%%%%%%%%%%%%%%%%%%%%%%%%%%%%%%%%%%%%
%%% Blocos:
\input{utils/blocos}

%%%%%%%%%%%%%%%%%%%%%%%%%%%%%%%%%%%%%%%%%%%%%%%%%%%%%%%%%%%%%%%%%%%%%%%%%%%%%%%%
%%% Boxes coloridos:
\input{utils/colorbox}

%%%%%%%%%%%%%%%%%%%%%%%%%%%%%%%%%%%%%%%%%%%%%%%%%%%%%%%%%%%%%%%%%%%%%%%%%%%%%%%%
%%% Tabelas:
\input{utils/tabelas}

%%%%%%%%%%%%%%%%%%%%%%%%%%%%%%%%%%%%%%%%%%%%%%%%%%%%%%%%%%%%%%%%%%%%%%%%%%%%%%%%
%%% Ambientes de listas:
\input{utils/listas}

%%%%%%%%%%%%%%%%%%%%%%%%%%%%%%%%%%%%%%%%%%%%%%%%%%%%%%%%%%%%%%%%%%%%%%%%%%%%%%%%
%%% Ambientes para quadros:
%\input{utils/quadros}

%%%%%%%%%%%%%%%%%%%%%%%%%%%%%%%%%%%%%%%%%%%%%%%%%%%%%%%%%%%%%%%%%%%%%%%%%%%%%%%%
%%% Ambientes floats e similares:
\input{utils/floats}

%%%%%%%%%%%%%%%%%%%%%%%%%%%%%%%%%%%%%%%%%%%%%%%%%%%%%%%%%%%%%%%%%%%%%%%%%%%%%%%%
%%% Notas de rodapé:
\input{utils/footnotes}

%%%%%%%%%%%%%%%%%%%%%%%%%%%%%%%%%%%%%%%%%%%%%%%%%%%%%%%%%%%%%%%%%%%%%%%%%%%%%%%%
%%% Referências cruzadas, links, citações:
\input{utils/crossrefs}

%%%%%%%%%%%%%%%%%%%%%%%%%%%%%%%%%%%%%%%%%%%%%%%%%%%%%%%%%%%%%%%%%%%%%%%%%%%%%%%%
%%% Configurações para os hiperlinks em PDF:
\input{utils/pdfconfig}

%%%%%%%%%%%%%%%%%%%%%%%%%%%%%%%%%%%%%%%%%%%%%%%%%%%%%%%%%%%%%%%%%%%%%%%%%%%%%%%%
%%% Glossário e índice remissivo:
\input{utils/glossindex}

%%%%%%%%%%%%%%%%%%%%%%%%%%%%%%%%%%%%%%%%%%%%%%%%%%%%%%%%%%%%%%%%%%%%%%%%%%%%%%%%
%%% Referências bibliográficas:
\input{utils/bibliografia}

%%%%%%%%%%%%%%%%%%%%%%%%%%%%%%%%%%%%%%%%%%%%%%%%%%%%%%%%%%%%%%%%%%%%%%%%%%%%%%%%
%%% Ajustes para a classe EXAM:
\input{utils/exam}

%%%%%%%%%%%%%%%%%%%%%%%%%%%%%%%%%%%%%%%%%%%%%%%%%%%%%%%%%%%%%%%%%%%%%%%%%%%%%%%%
%%% Ajustes para textos:
\input{utils/texto}

%%%%%%%%%%%%%%%%%%%%%%%%%%%%%%%%%%%%%%%%%%%%%%%%%%%%%%%%%%%%%%%%%%%%%%%%%%%%%%%%
%%% Ambientes, aliases, comandos e customizações em geral para serem
%%% utilizadas nos documentos. Defina aqui qualquer customização específica
%%% para seus documentos!
\input{utils/customizacoes}
\input{utils/juridico}
\input{utils/engenharia}
, acrescente:
%    %%%%%%%%%%%%%%%%%%%%%%%%%%%%%%%%%%%%%%%%%%%%%%%%%%%%%%%%%%%%%%%%%%%%%%%%%%%%%%%%
% engenharia.tex
%
% Configurações específicas do livro "Engenharia em 500 Gramas": pacotes de
% engenharia (circuitos, gráficos de dados, diagramas de tempo), caixas
% didáticas semânticas e comandos do livro.
%
% Deve ser carregado DEPOIS de utils/preambulo.tex (que já carrega tikz,
% tcolorbox[many], listings, siunitx, xcolor[svgnames] e cleveref). Para isso,
% em livro.tex, logo após \input{utils/preambulo.tex}, acrescente:
%    \input{utils/engenharia.tex}
%
% Nada aqui altera arquivos do modelo. Os capítulos devem usar SOMENTE os
% ambientes semânticos abaixo; a aparência é definida neste arquivo.
%%%%%%%%%%%%%%%%%%%%%%%%%%%%%%%%%%%%%%%%%%%%%%%%%%%%%%%%%%%%%%%%%%%%%%%%%%%%%%%%


%%%%%%%%%%%%%%%%%%%%%%%%%%%%%%%%%%%%%%%%%%%%%%%%%%%%%%%%%%%%%%%%%%%%%%%%%%%%%%%%
%%% Circuitos: esquemas elétricos. Estilo europeu (resistor retangular), o
%%% usado no Brasil. A opção siunitx faz os rótulos usarem \SI, herdando a
%%% vírgula decimal configurada em utils/matematica.tex.
\usepackage[siunitx, european]{circuitikz}
\ctikzset{bipoles/length=1.2cm}

%%%%%%%%%%%%%%%%%%%%%%%%%%%%%%%%%%%%%%%%%%%%%%%%%%%%%%%%%%%%%%%%%%%%%%%%%%%%%%%%
%%% Gráficos de dados: curvas dos ensaios (motor, sensores, bateria).
%%% Vírgula decimal e espaço fino como separador de milhar, como no siunitx.
\usepackage{pgfplots}
\usepackage{pgfplotstable}
\pgfplotsset{
  compat=1.18,
  /pgf/number format/.cd, use comma, 1000 sep={\,},
  every axis/.append style={grid=major, grid style={dotted, gray!60}},
}

%%%%%%%%%%%%%%%%%%%%%%%%%%%%%%%%%%%%%%%%%%%%%%%%%%%%%%%%%%%%%%%%%%%%%%%%%%%%%%%%
%%% Bibliotecas TikZ para diagramas de blocos, mecânica e o robô no dohyô.
%%% (utils/graficos.tex já carrega arrows, automata, backgrounds, matrix,
%%% patterns, positioning, shapes, shadows, calc e circuits.logic.US.)
\usetikzlibrary{arrows.meta, decorations.pathmorphing, decorations.markings,
                fit, chains, angles}
% Não carregue a biblioteca "quotes": a sintaxe ["rótulo"] em \draw e \pic
% falha neste modelo (o caractere " é ativo por causa do babel/polyglossia).
% Use "node[above]{$\theta$}" e "pic text=$\alpha$" no lugar. Alternativa:
% \shorthandoff{"} antes do tikzpicture (e a biblioteca quotes), mas isso
% desliga localmente os atalhos "" e "- do português.

%%%%%%%%%%%%%%%%%%%%%%%%%%%%%%%%%%%%%%%%%%%%%%%%%%%%%%%%%%%%%%%%%%%%%%%%%%%%%%%%
%%% Diagramas de tempo (PWM, pulso do HC-SR04, protocolo IR) e formatos de
%%% dados/registradores.
\usepackage{tikz-timing}
\usepackage{bytefield}

%%%%%%%%%%%%%%%%%%%%%%%%%%%%%%%%%%%%%%%%%%%%%%%%%%%%%%%%%%%%%%%%%%%%%%%%%%%%%%%%
%%% QR codes para os vídeos complementares.
\usepackage{qrcode}

%%% Árvores (classificações, hierarquias).
\usepackage[edges]{forest}


%%%%%%%%%%%%%%%%%%%%%%%%%%%%%%%%%%%%%%%%%%%%%%%%%%%%%%%%%%%%%%%%%%%%%%%%%%%%%%%%
%%% Tons de cinza usados pelas caixas (edição em preto e branco). Para uma
%%% edição colorida, basta redefinir estas cores.
\providecolor{cinzaClaro}{gray}{0.90}
\providecolor{cinzaMedio}{gray}{0.55}
\providecolor{cinzaEscuro}{gray}{0.25}


%%%%%%%%%%%%%%%%%%%%%%%%%%%%%%%%%%%%%%%%%%%%%%%%%%%%%%%%%%%%%%%%%%%%%%%%%%%%%%%%
%%% Comandos do livro

% Plataformas: usar sempre os comandos, nunca digitar à mão.
\newcommand{\RI}{R1}
\newcommand{\RII}{R2}

% Marcas de rascunho: visíveis enquanto \rascunhotrue; somem na edição final
% (troque para \rascunhofalse).
\newif\ifrascunho
\rascunhotrue
\newcommand{\TODO}[1]{\ifrascunho{\sffamily\footnotesize\color{cinzaEscuro}%
  [\textbf{A FAZER:} #1]}\fi}

% \codigo[opções listings]{arquivo}{marcador}
% Extrai do arquivo-fonte REAL o trecho delimitado por
%    // @livro: <marcador>
%    ...
%    // @livro: <marcador>_fim
% Assim o livro nunca diverge do firmware: se o código muda, o trecho muda
% junto. Usa o estilo CStyle definido em utils/computacao.tex.
%
% Limitações do listings (testadas no XeLaTeX):
%  1. Caracteres acima de U+00FF no arquivo-fonte (travessão, símbolo de
%     aproximação, emojis) geram linhas espúrias. Acentos do português são
%     seguros.
%  2. Comentário de bloco /* ... */ de várias linhas dentro de um trecho
%     marcado faz o marcador final ser ignorado. Dentro de trechos marcados,
%     use //.
\newcommand{\codigo}[3][]{%
  \lstinputlisting[style=CStyle,
    linerange={#3-#3_fim},
    rangeprefix=//\ @livro:\ , rangesuffix=,
    includerangemarker=false, #1]{#2}}


%%%%%%%%%%%%%%%%%%%%%%%%%%%%%%%%%%%%%%%%%%%%%%%%%%%%%%%%%%%%%%%%%%%%%%%%%%%%%%%%
%%% Caixas didáticas semânticas
% Recuo de parágrafo dentro das caixas: o mesmo do texto corrido. O tcolorbox
% zera \parindent no conteúdo; o primeiro parágrafo fica sem recuo.
\newlength{\recuocaixa}
\AtBeginDocument{\setlength{\recuocaixa}{\parindent}}
\tcbset{
  e500base/.style={
    enhanced, breakable, sharp corners,
    colback=white, colframe=black, coltitle=black,
    boxrule=0.4pt, left=2mm, right=2mm, top=1.2mm, bottom=1.2mm,
    fonttitle=\sffamily\bfseries\small,
    before skip=0.8\baselineskip, after skip=0.8\baselineskip,
    before upper={\setlength{\parindent}{\recuocaixa}\noindent},
  },
  % Título na primeira linha do conteúdo. Substitui "attach title to upper",
  % que redefine before upper e perderia o recuo acima.
  e500tituloInterno/.style={detach title,
    before upper={\setlength{\parindent}{\recuocaixa}%
      \noindent\tcbtitle\par\smallskip\noindent}},
}

% Fundamento: teoria transferível, independente do robô.
\newtcolorbox{fundamento}[1][]{e500base,
  frame hidden, borderline west={2.5pt}{0pt}{black},
  title={Fundamento\ifx&#1&\else\ --- #1\fi},
  e500tituloInterno}

% Modelo: desenvolvimento matemático.
\newtcolorbox{modelo}[1][]{e500base,
  colbacktitle=cinzaClaro, titlerule=0.4pt,
  title={Modelo\ifx&#1&\else\ --- #1\fi}}

% Projeto R1 / Projeto R2: decisão aplicada a cada plataforma.
\newtcolorbox{projetoRI}[1][]{e500base,
  colbacktitle=black, coltitle=white,
  title={Projeto R1\ifx&#1&\else\ --- #1\fi}}
\newtcolorbox{projetoRII}[1][]{e500base,
  colbacktitle=cinzaEscuro, coltitle=white, boxrule=0.8pt,
  title={Projeto R2\ifx&#1&\else\ --- #1\fi}}

% Laboratório: experimento obrigatório do capítulo.
\newtcolorbox{laboratorio}[1][]{e500base,
  frame hidden, borderline north={1.2pt}{0pt}{black},
  borderline south={1.2pt}{0pt}{black},
  title={Laboratório\ifx&#1&\else\ --- #1\fi},
  e500tituloInterno}

% Quando a engenharia dá errado: falha real, documentada.
\newtcolorbox{falhareal}[1][]{e500base,
  frame hidden, borderline west={5pt}{0pt}{black},
  title={Quando a engenharia dá errado\ifx&#1&\else\ --- #1\fi},
  e500tituloInterno}

% Regra -> requisito: regulamento traduzido em requisito verificável.
\newtcolorbox{regra}[1][]{e500base,
  frame hidden,
  borderline={0.6pt}{0pt}{black,dashed},
  title={Regra $\rightarrow$ requisito\ifx&#1&\else\ --- #1\fi},
  e500tituloInterno}

% Competição: impacto da decisão no desempenho.
\newtcolorbox{competicao}[1][]{e500base,
  colframe=cinzaMedio, boxrule=0.8pt,
  title={Competição\ifx&#1&\else\ --- #1\fi},
  e500tituloInterno}

% Avançado: conteúdo opcional; pode ser pulado sem perda de continuidade.
\newtcolorbox{avancado}[1][]{e500base,
  colframe=cinzaMedio, boxrule=0.4pt,
  borderline={0.4pt}{1.2pt}{cinzaMedio},
  title={Avançado (opcional)\ifx&#1&\else\ --- #1\fi},
  e500tituloInterno}

% Atenção: cuidado prático, risco de dano ou erro comum de uso.
\newtcolorbox{atencao}[1][]{e500base,
  boxrule=1.2pt,
  title={Atenção\ifx&#1&\else\ --- #1\fi},
  e500tituloInterno}

% Vídeo complementar: nunca indispensável para entender o texto.
%    \video{V05.1}{Caracterização do motor em bancada}{https://...}
\newcommand{\video}[3]{%
  \begin{tcolorbox}[e500base, colframe=cinzaMedio, boxrule=0.4pt,
                    sidebyside, lefthand width=18mm, sidebyside gap=3mm]
    \qrcode[height=16mm]{#3}
    \tcblower
    {\sffamily\bfseries\small Vídeo complementar #1}\par
    #2\par
    {\footnotesize\ttfamily\url{#3}}
  \end{tcolorbox}}

%
% Nada aqui altera arquivos do modelo. Os capítulos devem usar SOMENTE os
% ambientes semânticos abaixo; a aparência é definida neste arquivo.
%%%%%%%%%%%%%%%%%%%%%%%%%%%%%%%%%%%%%%%%%%%%%%%%%%%%%%%%%%%%%%%%%%%%%%%%%%%%%%%%


%%%%%%%%%%%%%%%%%%%%%%%%%%%%%%%%%%%%%%%%%%%%%%%%%%%%%%%%%%%%%%%%%%%%%%%%%%%%%%%%
%%% Circuitos: esquemas elétricos. Estilo europeu (resistor retangular), o
%%% usado no Brasil. A opção siunitx faz os rótulos usarem \SI, herdando a
%%% vírgula decimal configurada em utils/matematica.tex.
\usepackage[siunitx, european]{circuitikz}
\ctikzset{bipoles/length=1.2cm}

%%%%%%%%%%%%%%%%%%%%%%%%%%%%%%%%%%%%%%%%%%%%%%%%%%%%%%%%%%%%%%%%%%%%%%%%%%%%%%%%
%%% Gráficos de dados: curvas dos ensaios (motor, sensores, bateria).
%%% Vírgula decimal e espaço fino como separador de milhar, como no siunitx.
\usepackage{pgfplots}
\usepackage{pgfplotstable}
\pgfplotsset{
  compat=1.18,
  /pgf/number format/.cd, use comma, 1000 sep={\,},
  every axis/.append style={grid=major, grid style={dotted, gray!60}},
}

%%%%%%%%%%%%%%%%%%%%%%%%%%%%%%%%%%%%%%%%%%%%%%%%%%%%%%%%%%%%%%%%%%%%%%%%%%%%%%%%
%%% Bibliotecas TikZ para diagramas de blocos, mecânica e o robô no dohyô.
%%% (utils/graficos.tex já carrega arrows, automata, backgrounds, matrix,
%%% patterns, positioning, shapes, shadows, calc e circuits.logic.US.)
\usetikzlibrary{arrows.meta, decorations.pathmorphing, decorations.markings,
                fit, chains, angles}
% Não carregue a biblioteca "quotes": a sintaxe ["rótulo"] em \draw e \pic
% falha neste modelo (o caractere " é ativo por causa do babel/polyglossia).
% Use "node[above]{$\theta$}" e "pic text=$\alpha$" no lugar. Alternativa:
% \shorthandoff{"} antes do tikzpicture (e a biblioteca quotes), mas isso
% desliga localmente os atalhos "" e "- do português.

%%%%%%%%%%%%%%%%%%%%%%%%%%%%%%%%%%%%%%%%%%%%%%%%%%%%%%%%%%%%%%%%%%%%%%%%%%%%%%%%
%%% Diagramas de tempo (PWM, pulso do HC-SR04, protocolo IR) e formatos de
%%% dados/registradores.
\usepackage{tikz-timing}
\usepackage{bytefield}

%%%%%%%%%%%%%%%%%%%%%%%%%%%%%%%%%%%%%%%%%%%%%%%%%%%%%%%%%%%%%%%%%%%%%%%%%%%%%%%%
%%% QR codes para os vídeos complementares.
\usepackage{qrcode}

%%% Árvores (classificações, hierarquias).
\usepackage[edges]{forest}


%%%%%%%%%%%%%%%%%%%%%%%%%%%%%%%%%%%%%%%%%%%%%%%%%%%%%%%%%%%%%%%%%%%%%%%%%%%%%%%%
%%% Tons de cinza usados pelas caixas (edição em preto e branco). Para uma
%%% edição colorida, basta redefinir estas cores.
\providecolor{cinzaClaro}{gray}{0.90}
\providecolor{cinzaMedio}{gray}{0.55}
\providecolor{cinzaEscuro}{gray}{0.25}


%%%%%%%%%%%%%%%%%%%%%%%%%%%%%%%%%%%%%%%%%%%%%%%%%%%%%%%%%%%%%%%%%%%%%%%%%%%%%%%%
%%% Comandos do livro

% Plataformas: usar sempre os comandos, nunca digitar à mão.
\newcommand{\RI}{R1}
\newcommand{\RII}{R2}

% Marcas de rascunho: visíveis enquanto \rascunhotrue; somem na edição final
% (troque para \rascunhofalse).
\newif\ifrascunho
\rascunhotrue
\newcommand{\TODO}[1]{\ifrascunho{\sffamily\footnotesize\color{cinzaEscuro}%
  [\textbf{A FAZER:} #1]}\fi}

% \codigo[opções listings]{arquivo}{marcador}
% Extrai do arquivo-fonte REAL o trecho delimitado por
%    // @livro: <marcador>
%    ...
%    // @livro: <marcador>_fim
% Assim o livro nunca diverge do firmware: se o código muda, o trecho muda
% junto. Usa o estilo CStyle definido em utils/computacao.tex.
%
% Limitações do listings (testadas no XeLaTeX):
%  1. Caracteres acima de U+00FF no arquivo-fonte (travessão, símbolo de
%     aproximação, emojis) geram linhas espúrias. Acentos do português são
%     seguros.
%  2. Comentário de bloco /* ... */ de várias linhas dentro de um trecho
%     marcado faz o marcador final ser ignorado. Dentro de trechos marcados,
%     use //.
\newcommand{\codigo}[3][]{%
  \lstinputlisting[style=CStyle,
    linerange={#3-#3_fim},
    rangeprefix=//\ @livro:\ , rangesuffix=,
    includerangemarker=false, #1]{#2}}


%%%%%%%%%%%%%%%%%%%%%%%%%%%%%%%%%%%%%%%%%%%%%%%%%%%%%%%%%%%%%%%%%%%%%%%%%%%%%%%%
%%% Caixas didáticas semânticas
% Recuo de parágrafo dentro das caixas: o mesmo do texto corrido. O tcolorbox
% zera \parindent no conteúdo; o primeiro parágrafo fica sem recuo.
\newlength{\recuocaixa}
\AtBeginDocument{\setlength{\recuocaixa}{\parindent}}
\tcbset{
  e500base/.style={
    enhanced, breakable, sharp corners,
    colback=white, colframe=black, coltitle=black,
    boxrule=0.4pt, left=2mm, right=2mm, top=1.2mm, bottom=1.2mm,
    fonttitle=\sffamily\bfseries\small,
    before skip=0.8\baselineskip, after skip=0.8\baselineskip,
    before upper={\setlength{\parindent}{\recuocaixa}\noindent},
  },
  % Título na primeira linha do conteúdo. Substitui "attach title to upper",
  % que redefine before upper e perderia o recuo acima.
  e500tituloInterno/.style={detach title,
    before upper={\setlength{\parindent}{\recuocaixa}%
      \noindent\tcbtitle\par\smallskip\noindent}},
}

% Fundamento: teoria transferível, independente do robô.
\newtcolorbox{fundamento}[1][]{e500base,
  frame hidden, borderline west={2.5pt}{0pt}{black},
  title={Fundamento\ifx&#1&\else\ --- #1\fi},
  e500tituloInterno}

% Modelo: desenvolvimento matemático.
\newtcolorbox{modelo}[1][]{e500base,
  colbacktitle=cinzaClaro, titlerule=0.4pt,
  title={Modelo\ifx&#1&\else\ --- #1\fi}}

% Projeto R1 / Projeto R2: decisão aplicada a cada plataforma.
\newtcolorbox{projetoRI}[1][]{e500base,
  colbacktitle=black, coltitle=white,
  title={Projeto R1\ifx&#1&\else\ --- #1\fi}}
\newtcolorbox{projetoRII}[1][]{e500base,
  colbacktitle=cinzaEscuro, coltitle=white, boxrule=0.8pt,
  title={Projeto R2\ifx&#1&\else\ --- #1\fi}}

% Laboratório: experimento obrigatório do capítulo.
\newtcolorbox{laboratorio}[1][]{e500base,
  frame hidden, borderline north={1.2pt}{0pt}{black},
  borderline south={1.2pt}{0pt}{black},
  title={Laboratório\ifx&#1&\else\ --- #1\fi},
  e500tituloInterno}

% Quando a engenharia dá errado: falha real, documentada.
\newtcolorbox{falhareal}[1][]{e500base,
  frame hidden, borderline west={5pt}{0pt}{black},
  title={Quando a engenharia dá errado\ifx&#1&\else\ --- #1\fi},
  e500tituloInterno}

% Regra -> requisito: regulamento traduzido em requisito verificável.
\newtcolorbox{regra}[1][]{e500base,
  frame hidden,
  borderline={0.6pt}{0pt}{black,dashed},
  title={Regra $\rightarrow$ requisito\ifx&#1&\else\ --- #1\fi},
  e500tituloInterno}

% Competição: impacto da decisão no desempenho.
\newtcolorbox{competicao}[1][]{e500base,
  colframe=cinzaMedio, boxrule=0.8pt,
  title={Competição\ifx&#1&\else\ --- #1\fi},
  e500tituloInterno}

% Avançado: conteúdo opcional; pode ser pulado sem perda de continuidade.
\newtcolorbox{avancado}[1][]{e500base,
  colframe=cinzaMedio, boxrule=0.4pt,
  borderline={0.4pt}{1.2pt}{cinzaMedio},
  title={Avançado (opcional)\ifx&#1&\else\ --- #1\fi},
  e500tituloInterno}

% Atenção: cuidado prático, risco de dano ou erro comum de uso.
\newtcolorbox{atencao}[1][]{e500base,
  boxrule=1.2pt,
  title={Atenção\ifx&#1&\else\ --- #1\fi},
  e500tituloInterno}

% Vídeo complementar: nunca indispensável para entender o texto.
%    \video{V05.1}{Caracterização do motor em bancada}{https://...}
\newcommand{\video}[3]{%
  \begin{tcolorbox}[e500base, colframe=cinzaMedio, boxrule=0.4pt,
                    sidebyside, lefthand width=18mm, sidebyside gap=3mm]
    \qrcode[height=16mm]{#3}
    \tcblower
    {\sffamily\bfseries\small Vídeo complementar #1}\par
    #2\par
    {\footnotesize\ttfamily\url{#3}}
  \end{tcolorbox}}

%
% Nada aqui altera arquivos do modelo. Os capítulos devem usar SOMENTE os
% ambientes semânticos abaixo; a aparência é definida neste arquivo.
%%%%%%%%%%%%%%%%%%%%%%%%%%%%%%%%%%%%%%%%%%%%%%%%%%%%%%%%%%%%%%%%%%%%%%%%%%%%%%%%


%%%%%%%%%%%%%%%%%%%%%%%%%%%%%%%%%%%%%%%%%%%%%%%%%%%%%%%%%%%%%%%%%%%%%%%%%%%%%%%%
%%% Circuitos: esquemas elétricos. Estilo europeu (resistor retangular), o
%%% usado no Brasil. A opção siunitx faz os rótulos usarem \SI, herdando a
%%% vírgula decimal configurada em utils/matematica.tex.
\usepackage[siunitx, european]{circuitikz}
\ctikzset{bipoles/length=1.2cm}

%%%%%%%%%%%%%%%%%%%%%%%%%%%%%%%%%%%%%%%%%%%%%%%%%%%%%%%%%%%%%%%%%%%%%%%%%%%%%%%%
%%% Gráficos de dados: curvas dos ensaios (motor, sensores, bateria).
%%% Vírgula decimal e espaço fino como separador de milhar, como no siunitx.
\usepackage{pgfplots}
\usepackage{pgfplotstable}
\pgfplotsset{
  compat=1.18,
  /pgf/number format/.cd, use comma, 1000 sep={\,},
  every axis/.append style={grid=major, grid style={dotted, gray!60}},
}

%%%%%%%%%%%%%%%%%%%%%%%%%%%%%%%%%%%%%%%%%%%%%%%%%%%%%%%%%%%%%%%%%%%%%%%%%%%%%%%%
%%% Bibliotecas TikZ para diagramas de blocos, mecânica e o robô no dohyô.
%%% (utils/graficos.tex já carrega arrows, automata, backgrounds, matrix,
%%% patterns, positioning, shapes, shadows, calc e circuits.logic.US.)
\usetikzlibrary{arrows.meta, decorations.pathmorphing, decorations.markings,
                fit, chains, angles}
% Não carregue a biblioteca "quotes": a sintaxe ["rótulo"] em \draw e \pic
% falha neste modelo (o caractere " é ativo por causa do babel/polyglossia).
% Use "node[above]{$\theta$}" e "pic text=$\alpha$" no lugar. Alternativa:
% \shorthandoff{"} antes do tikzpicture (e a biblioteca quotes), mas isso
% desliga localmente os atalhos "" e "- do português.

%%%%%%%%%%%%%%%%%%%%%%%%%%%%%%%%%%%%%%%%%%%%%%%%%%%%%%%%%%%%%%%%%%%%%%%%%%%%%%%%
%%% Diagramas de tempo (PWM, pulso do HC-SR04, protocolo IR) e formatos de
%%% dados/registradores.
\usepackage{tikz-timing}
\usepackage{bytefield}

%%%%%%%%%%%%%%%%%%%%%%%%%%%%%%%%%%%%%%%%%%%%%%%%%%%%%%%%%%%%%%%%%%%%%%%%%%%%%%%%
%%% QR codes para os vídeos complementares.
\usepackage{qrcode}

%%% Árvores (classificações, hierarquias).
\usepackage[edges]{forest}


%%%%%%%%%%%%%%%%%%%%%%%%%%%%%%%%%%%%%%%%%%%%%%%%%%%%%%%%%%%%%%%%%%%%%%%%%%%%%%%%
%%% Tons de cinza usados pelas caixas (edição em preto e branco). Para uma
%%% edição colorida, basta redefinir estas cores.
\providecolor{cinzaClaro}{gray}{0.90}
\providecolor{cinzaMedio}{gray}{0.55}
\providecolor{cinzaEscuro}{gray}{0.25}


%%%%%%%%%%%%%%%%%%%%%%%%%%%%%%%%%%%%%%%%%%%%%%%%%%%%%%%%%%%%%%%%%%%%%%%%%%%%%%%%
%%% Comandos do livro

% Plataformas: usar sempre os comandos, nunca digitar à mão.
\newcommand{\RI}{R1}
\newcommand{\RII}{R2}

% Marcas de rascunho: visíveis enquanto \rascunhotrue; somem na edição final
% (troque para \rascunhofalse).
\newif\ifrascunho
\rascunhotrue
\newcommand{\TODO}[1]{\ifrascunho{\sffamily\footnotesize\color{cinzaEscuro}%
  [\textbf{A FAZER:} #1]}\fi}

% \codigo[opções listings]{arquivo}{marcador}
% Extrai do arquivo-fonte REAL o trecho delimitado por
%    // @livro: <marcador>
%    ...
%    // @livro: <marcador>_fim
% Assim o livro nunca diverge do firmware: se o código muda, o trecho muda
% junto. Usa o estilo CStyle definido em utils/computacao.tex.
%
% Limitações do listings (testadas no XeLaTeX):
%  1. Caracteres acima de U+00FF no arquivo-fonte (travessão, símbolo de
%     aproximação, emojis) geram linhas espúrias. Acentos do português são
%     seguros.
%  2. Comentário de bloco /* ... */ de várias linhas dentro de um trecho
%     marcado faz o marcador final ser ignorado. Dentro de trechos marcados,
%     use //.
\newcommand{\codigo}[3][]{%
  \lstinputlisting[style=CStyle,
    linerange={#3-#3_fim},
    rangeprefix=//\ @livro:\ , rangesuffix=,
    includerangemarker=false, #1]{#2}}


%%%%%%%%%%%%%%%%%%%%%%%%%%%%%%%%%%%%%%%%%%%%%%%%%%%%%%%%%%%%%%%%%%%%%%%%%%%%%%%%
%%% Caixas didáticas semânticas
% Recuo de parágrafo dentro das caixas: o mesmo do texto corrido. O tcolorbox
% zera \parindent no conteúdo; o primeiro parágrafo fica sem recuo.
\newlength{\recuocaixa}
\AtBeginDocument{\setlength{\recuocaixa}{\parindent}}
\tcbset{
  e500base/.style={
    enhanced, breakable, sharp corners,
    colback=white, colframe=black, coltitle=black,
    boxrule=0.4pt, left=2mm, right=2mm, top=1.2mm, bottom=1.2mm,
    fonttitle=\sffamily\bfseries\small,
    before skip=0.8\baselineskip, after skip=0.8\baselineskip,
    before upper={\setlength{\parindent}{\recuocaixa}\noindent},
  },
  % Título na primeira linha do conteúdo. Substitui "attach title to upper",
  % que redefine before upper e perderia o recuo acima.
  e500tituloInterno/.style={detach title,
    before upper={\setlength{\parindent}{\recuocaixa}%
      \noindent\tcbtitle\par\smallskip\noindent}},
}

% Fundamento: teoria transferível, independente do robô.
\newtcolorbox{fundamento}[1][]{e500base,
  frame hidden, borderline west={2.5pt}{0pt}{black},
  title={Fundamento\ifx&#1&\else\ --- #1\fi},
  e500tituloInterno}

% Modelo: desenvolvimento matemático.
\newtcolorbox{modelo}[1][]{e500base,
  colbacktitle=cinzaClaro, titlerule=0.4pt,
  title={Modelo\ifx&#1&\else\ --- #1\fi}}

% Projeto R1 / Projeto R2: decisão aplicada a cada plataforma.
\newtcolorbox{projetoRI}[1][]{e500base,
  colbacktitle=black, coltitle=white,
  title={Projeto R1\ifx&#1&\else\ --- #1\fi}}
\newtcolorbox{projetoRII}[1][]{e500base,
  colbacktitle=cinzaEscuro, coltitle=white, boxrule=0.8pt,
  title={Projeto R2\ifx&#1&\else\ --- #1\fi}}

% Laboratório: experimento obrigatório do capítulo.
\newtcolorbox{laboratorio}[1][]{e500base,
  frame hidden, borderline north={1.2pt}{0pt}{black},
  borderline south={1.2pt}{0pt}{black},
  title={Laboratório\ifx&#1&\else\ --- #1\fi},
  e500tituloInterno}

% Quando a engenharia dá errado: falha real, documentada.
\newtcolorbox{falhareal}[1][]{e500base,
  frame hidden, borderline west={5pt}{0pt}{black},
  title={Quando a engenharia dá errado\ifx&#1&\else\ --- #1\fi},
  e500tituloInterno}

% Regra -> requisito: regulamento traduzido em requisito verificável.
\newtcolorbox{regra}[1][]{e500base,
  frame hidden,
  borderline={0.6pt}{0pt}{black,dashed},
  title={Regra $\rightarrow$ requisito\ifx&#1&\else\ --- #1\fi},
  e500tituloInterno}

% Competição: impacto da decisão no desempenho.
\newtcolorbox{competicao}[1][]{e500base,
  colframe=cinzaMedio, boxrule=0.8pt,
  title={Competição\ifx&#1&\else\ --- #1\fi},
  e500tituloInterno}

% Avançado: conteúdo opcional; pode ser pulado sem perda de continuidade.
\newtcolorbox{avancado}[1][]{e500base,
  colframe=cinzaMedio, boxrule=0.4pt,
  borderline={0.4pt}{1.2pt}{cinzaMedio},
  title={Avançado (opcional)\ifx&#1&\else\ --- #1\fi},
  e500tituloInterno}

% Atenção: cuidado prático, risco de dano ou erro comum de uso.
\newtcolorbox{atencao}[1][]{e500base,
  boxrule=1.2pt,
  title={Atenção\ifx&#1&\else\ --- #1\fi},
  e500tituloInterno}

% Vídeo complementar: nunca indispensável para entender o texto.
%    \video{V05.1}{Caracterização do motor em bancada}{https://...}
\newcommand{\video}[3]{%
  \begin{tcolorbox}[e500base, colframe=cinzaMedio, boxrule=0.4pt,
                    sidebyside, lefthand width=18mm, sidebyside gap=3mm]
    \qrcode[height=16mm]{#3}
    \tcblower
    {\sffamily\bfseries\small Vídeo complementar #1}\par
    #2\par
    {\footnotesize\ttfamily\url{#3}}
  \end{tcolorbox}}

%
% Nada aqui altera arquivos do modelo. Os capítulos devem usar SOMENTE os
% ambientes semânticos abaixo; a aparência é definida neste arquivo.
%%%%%%%%%%%%%%%%%%%%%%%%%%%%%%%%%%%%%%%%%%%%%%%%%%%%%%%%%%%%%%%%%%%%%%%%%%%%%%%%


%%%%%%%%%%%%%%%%%%%%%%%%%%%%%%%%%%%%%%%%%%%%%%%%%%%%%%%%%%%%%%%%%%%%%%%%%%%%%%%%
%%% Circuitos: esquemas elétricos. Estilo europeu (resistor retangular), o
%%% usado no Brasil. A opção siunitx faz os rótulos usarem \SI, herdando a
%%% vírgula decimal configurada em utils/matematica.tex.
\usepackage[siunitx, european]{circuitikz}
\ctikzset{bipoles/length=1.2cm}

%%%%%%%%%%%%%%%%%%%%%%%%%%%%%%%%%%%%%%%%%%%%%%%%%%%%%%%%%%%%%%%%%%%%%%%%%%%%%%%%
%%% Gráficos de dados: curvas dos ensaios (motor, sensores, bateria).
%%% Vírgula decimal e espaço fino como separador de milhar, como no siunitx.
\usepackage{pgfplots}
\usepackage{pgfplotstable}
\pgfplotsset{
  compat=1.18,
  /pgf/number format/.cd, use comma, 1000 sep={\,},
  every axis/.append style={grid=major, grid style={dotted, gray!60}},
}

%%%%%%%%%%%%%%%%%%%%%%%%%%%%%%%%%%%%%%%%%%%%%%%%%%%%%%%%%%%%%%%%%%%%%%%%%%%%%%%%
%%% Bibliotecas TikZ para diagramas de blocos, mecânica e o robô no dohyô.
%%% (utils/graficos.tex já carrega arrows, automata, backgrounds, matrix,
%%% patterns, positioning, shapes, shadows, calc e circuits.logic.US.)
\usetikzlibrary{arrows.meta, decorations.pathmorphing, decorations.markings,
                fit, chains, angles}
% Não carregue a biblioteca "quotes": a sintaxe ["rótulo"] em \draw e \pic
% falha neste modelo (o caractere " é ativo por causa do babel/polyglossia).
% Use "node[above]{$\theta$}" e "pic text=$\alpha$" no lugar. Alternativa:
% \shorthandoff{"} antes do tikzpicture (e a biblioteca quotes), mas isso
% desliga localmente os atalhos "" e "- do português.

%%%%%%%%%%%%%%%%%%%%%%%%%%%%%%%%%%%%%%%%%%%%%%%%%%%%%%%%%%%%%%%%%%%%%%%%%%%%%%%%
%%% Diagramas de tempo (PWM, pulso do HC-SR04, protocolo IR) e formatos de
%%% dados/registradores.
\usepackage{tikz-timing}
\usepackage{bytefield}

%%%%%%%%%%%%%%%%%%%%%%%%%%%%%%%%%%%%%%%%%%%%%%%%%%%%%%%%%%%%%%%%%%%%%%%%%%%%%%%%
%%% QR codes para os vídeos complementares.
\usepackage{qrcode}

%%% Árvores (classificações, hierarquias).
\usepackage[edges]{forest}


%%%%%%%%%%%%%%%%%%%%%%%%%%%%%%%%%%%%%%%%%%%%%%%%%%%%%%%%%%%%%%%%%%%%%%%%%%%%%%%%
%%% Tons de cinza usados pelas caixas (edição em preto e branco). Para uma
%%% edição colorida, basta redefinir estas cores.
\providecolor{cinzaClaro}{gray}{0.90}
\providecolor{cinzaMedio}{gray}{0.55}
\providecolor{cinzaEscuro}{gray}{0.25}


%%%%%%%%%%%%%%%%%%%%%%%%%%%%%%%%%%%%%%%%%%%%%%%%%%%%%%%%%%%%%%%%%%%%%%%%%%%%%%%%
%%% Comandos do livro

% Plataformas: usar sempre os comandos, nunca digitar à mão.
\newcommand{\RI}{R1}
\newcommand{\RII}{R2}

% Marcas de rascunho: visíveis enquanto \rascunhotrue; somem na edição final
% (troque para \rascunhofalse).
\newif\ifrascunho
\rascunhotrue
\newcommand{\TODO}[1]{\ifrascunho{\sffamily\footnotesize\color{cinzaEscuro}%
  [\textbf{A FAZER:} #1]}\fi}

% \codigo[opções listings]{arquivo}{marcador}
% Extrai do arquivo-fonte REAL o trecho delimitado por
%    // @livro: <marcador>
%    ...
%    // @livro: <marcador>_fim
% Assim o livro nunca diverge do firmware: se o código muda, o trecho muda
% junto. Usa o estilo CStyle definido em utils/computacao.tex.
%
% Limitações do listings (testadas no XeLaTeX):
%  1. Caracteres acima de U+00FF no arquivo-fonte (travessão, símbolo de
%     aproximação, emojis) geram linhas espúrias. Acentos do português são
%     seguros.
%  2. Comentário de bloco /* ... */ de várias linhas dentro de um trecho
%     marcado faz o marcador final ser ignorado. Dentro de trechos marcados,
%     use //.
\newcommand{\codigo}[3][]{%
  \lstinputlisting[style=CStyle,
    linerange={#3-#3_fim},
    rangeprefix=//\ @livro:\ , rangesuffix=,
    includerangemarker=false, #1]{#2}}


%%%%%%%%%%%%%%%%%%%%%%%%%%%%%%%%%%%%%%%%%%%%%%%%%%%%%%%%%%%%%%%%%%%%%%%%%%%%%%%%
%%% Caixas didáticas semânticas
% Recuo de parágrafo dentro das caixas: o mesmo do texto corrido. O tcolorbox
% zera \parindent no conteúdo; o primeiro parágrafo fica sem recuo.
\newlength{\recuocaixa}
\AtBeginDocument{\setlength{\recuocaixa}{\parindent}}
\tcbset{
  e500base/.style={
    enhanced, breakable, sharp corners,
    colback=white, colframe=black, coltitle=black,
    boxrule=0.4pt, left=2mm, right=2mm, top=1.2mm, bottom=1.2mm,
    fonttitle=\sffamily\bfseries\small,
    before skip=0.8\baselineskip, after skip=0.8\baselineskip,
    before upper={\setlength{\parindent}{\recuocaixa}\noindent},
  },
  % Título na primeira linha do conteúdo. Substitui "attach title to upper",
  % que redefine before upper e perderia o recuo acima.
  e500tituloInterno/.style={detach title,
    before upper={\setlength{\parindent}{\recuocaixa}%
      \noindent\tcbtitle\par\smallskip\noindent}},
}

% Fundamento: teoria transferível, independente do robô.
\newtcolorbox{fundamento}[1][]{e500base,
  frame hidden, borderline west={2.5pt}{0pt}{black},
  title={Fundamento\ifx&#1&\else\ --- #1\fi},
  e500tituloInterno}

% Modelo: desenvolvimento matemático.
\newtcolorbox{modelo}[1][]{e500base,
  colbacktitle=cinzaClaro, titlerule=0.4pt,
  title={Modelo\ifx&#1&\else\ --- #1\fi}}

% Projeto R1 / Projeto R2: decisão aplicada a cada plataforma.
\newtcolorbox{projetoRI}[1][]{e500base,
  colbacktitle=black, coltitle=white,
  title={Projeto R1\ifx&#1&\else\ --- #1\fi}}
\newtcolorbox{projetoRII}[1][]{e500base,
  colbacktitle=cinzaEscuro, coltitle=white, boxrule=0.8pt,
  title={Projeto R2\ifx&#1&\else\ --- #1\fi}}

% Laboratório: experimento obrigatório do capítulo.
\newtcolorbox{laboratorio}[1][]{e500base,
  frame hidden, borderline north={1.2pt}{0pt}{black},
  borderline south={1.2pt}{0pt}{black},
  title={Laboratório\ifx&#1&\else\ --- #1\fi},
  e500tituloInterno}

% Quando a engenharia dá errado: falha real, documentada.
\newtcolorbox{falhareal}[1][]{e500base,
  frame hidden, borderline west={5pt}{0pt}{black},
  title={Quando a engenharia dá errado\ifx&#1&\else\ --- #1\fi},
  e500tituloInterno}

% Regra -> requisito: regulamento traduzido em requisito verificável.
\newtcolorbox{regra}[1][]{e500base,
  frame hidden,
  borderline={0.6pt}{0pt}{black,dashed},
  title={Regra $\rightarrow$ requisito\ifx&#1&\else\ --- #1\fi},
  e500tituloInterno}

% Competição: impacto da decisão no desempenho.
\newtcolorbox{competicao}[1][]{e500base,
  colframe=cinzaMedio, boxrule=0.8pt,
  title={Competição\ifx&#1&\else\ --- #1\fi},
  e500tituloInterno}

% Avançado: conteúdo opcional; pode ser pulado sem perda de continuidade.
\newtcolorbox{avancado}[1][]{e500base,
  colframe=cinzaMedio, boxrule=0.4pt,
  borderline={0.4pt}{1.2pt}{cinzaMedio},
  title={Avançado (opcional)\ifx&#1&\else\ --- #1\fi},
  e500tituloInterno}

% Atenção: cuidado prático, risco de dano ou erro comum de uso.
\newtcolorbox{atencao}[1][]{e500base,
  boxrule=1.2pt,
  title={Atenção\ifx&#1&\else\ --- #1\fi},
  e500tituloInterno}

% Vídeo complementar: nunca indispensável para entender o texto.
%    \video{V05.1}{Caracterização do motor em bancada}{https://...}
\newcommand{\video}[3]{%
  \begin{tcolorbox}[e500base, colframe=cinzaMedio, boxrule=0.4pt,
                    sidebyside, lefthand width=18mm, sidebyside gap=3mm]
    \qrcode[height=16mm]{#3}
    \tcblower
    {\sffamily\bfseries\small Vídeo complementar #1}\par
    #2\par
    {\footnotesize\ttfamily\url{#3}}
  \end{tcolorbox}}
